% concatenated from Continuous_linear_series_II.tex
\documentclass[a4paper,11pt]{amsart}
\usepackage[utf8]{inputenc}
\usepackage{amsmath,amsfonts,amssymb,anysize,amscd}
\usepackage[all,cmtip]{xy}
\usepackage[onehalfspacing]{setspace}
\usepackage{amsthm}
\usepackage{graphicx}
\usepackage{pstricks-add}
\usepackage{pst-all}
\usepackage{caption}
%\usepackage{draftwatermark}
%\usepackage{showkeys}

%\newpsobject{grilla}{psgrid}{subgriddiv=10,griddots=10,gridlabels=6pt}

\theoremstyle{plain}
\newtheorem{thm}{Theorem}[section]
\newtheorem{cor}[thm]{Corollary}
\newtheorem{lem}[thm]{Lemma}
\newtheorem{prop}[thm]{Proposition}

\theoremstyle{definition}
\newtheorem{defn}[thm]{Definition}

\theoremstyle{remark}
\newtheorem{rem}[thm]{Remark}
\newtheorem{Ex}{Example}[section]

%watermark
%\SetWatermarkScale{6}
%\SetWatermarkText{Borrador}
%\SetWatermarkLightness{0.5}


% Short Command.
\def\r{\mathbb{R}}
\def\p{\mathbb{P}}
\def\z{\mathbf{Z}}
\def\cc{{\mathbf G}_{\mathbf m}}
\def\q{\mathbb{Q}}
\def\n{\mathbf{N}}
\def\a{\mathbb{A}}
\def\x{\mathbb{X}}
\def\y{\mathbb{Y}}
\def\j{JC}
\def\J{\overline{JC}}
\def\L{\mathcal{L}}
\def\X{\mathcal{X}}
\def\O{\mathcal{O}}
\def\I{\mathcal{I}}
\def\N{\mathcal{N}}
\def\V{\mathcal{V}}
\def\ol{\overline}
\def\Pic{{\rm Pic}}
\def\P{\mathcal{P}}
\def\W{\mathcal{W}}
\def\rk{{\rm rk}}
\def\Ker{{\rm Ker}}
\def\Im{{\rm Im}}
\def\g{{\mathfrak g}}
\def\IP{{\mathbf P}}
\def\div{\mathrm{div}}
\def\lra{\longrightarrow}
\def\F{\mathcal{F}}
\def\E{\mathcal{E}}
%\def\S{\mathcal{S}}
\def\:{\colon}
\def\ox{\otimes}
\def\wh{\widehat}
\def\red{\text{\rm red}}

\def\uDelta{{\underline\delta}}

%opening
\title{Continuous linear series}
%\author{E.~Esteves, A.~Nigro and P.~Rizzo}

\author[Esteves]{Eduardo Esteves}
\address{Instituto de Matem\'atica Pura e Aplicada, IMPA, Rio de Janeiro, RJ\\Brazil}
\email{esteves@impa.br}

\author[Nigro]{Antonio Nigro}
\address{Universidade Federal Fluminense, Instituto de Matem\'atica e Estat\'istica, Niter\'oi, RJ, Brazil}
\email{antonio.nigro@gmail.com}

\author[Rizzo]{Pedro  Rizzo}
\address{Instituto de Matemáticas\\Universidad de Antioquia\\ Medell\'in, Colombia}
\email{pedro.hernandez@udea.edu.co }

\keywords{(Limit) Linear series on Curves, Hilbert Schemes, Moduli Problems, Compactification}

\date{}
\begin{document}

\begin{abstract}
We parameterize by a fine moduli space all degenerations of linear series to a singular curve which is the union of two smooth components meeting transversally at a single point. For this we introduce a novel object in the study of degenerations of linear series, which is the \emph{continuous linear series}. Our moduli space can be regarded as a Hilbert quotient, in the terminology introduced by Kapranov, and is a new compactification of Osserman moduli space of exact limit linear series, and consequently, of Eisenbud and Harris moduli space of refined limit linear series on the curve.
%E< This paper introduces the concept of \textit{continuous linear series} on a singular curve $X$, where $X$ is the union of two smooth components $Y$ and $Z$ intersecting transversally in a unique node $P$. We construct a projective moduli space that provides a compactification of Osserman exact limit linear series, and consequently, of Eisenbud-Harris refined limit linear series, employing exclusively exact limit linear series—those limits with desirable properties. Furthermore, our construction provides novel insights into the theory of \textit{moduli spaces of complete collineations}, as presented by Thaddeus.
\end{abstract}
\subjclass[2020]{14H10, 14C05, 14D22}
\maketitle

\section{Introduction.}
Fix a singular curve $X$ consisting of two smooth components $Y$ and $Z$ meeting transversally at a point $P$. A \emph{generalized linear series} on $X$ is a space of sections of a torsion-free, rank-1 sheaf on $X$. For short, a \emph{continuous linear series} on $X$ is a \emph{maximal primitive} (definitions given later) family of generalized linear series on $X$ parameterized by an Artin stack which is the quotient of a two-punctured semistable chain of rational curves modulo an action of the one-dimensional torus $\cc$. We parameterize them by a projective scheme. 

The history of the study and applications of degenerations of linear series goes a long way back, even predating work by the Italian School of Algebraic Geometry of Severi and others from the turn to the 20th Century. The modern study arose from applications to Brill--Noether Theory of the theory of limit linear series developed by Eisenbud and Harris in the 1980's. 

Degenerating linear series along families of smooth curves degenerating to a singular curve have many limits. For limit curves of compact type, Eisenbud and Harris \cite{EH1} took those limits whose restrictions to the components of the limit curve have degree zero but for a component, the focus, and verified the collection of their restrictions to their respective foci is a \emph{limit linear series}. A limit linear series, according to Eisenbud and Harris \cite{EH1}, pp.~345--6, is a collection of linear series of same rank and degree, one on each component of the limit curve, satisfying certain inequalities of order sequences. It is called \emph{refined} if all the inequalities are equalities. Eisenbud and Harris proceeded to produce a moduli space parameterizing refined limit linear series over curves of compact type, \cite{EH1}, Thm.~3.3, p.~354. 
%{\color{red} Refinados são mais fáceis de trabalhar e tem propriedades importantes.} 
Limit linear series arising from degenerations may not be refined, so their space is only quasi-projective.

Multiple applications of their theory were found to this day; see \cite{EH2,EH3,EH4,Cot,Tar,Cast,Fark1,Fark2} for a sample. On the other hand, multiple efforts (including by the authors) have not yet produced a natural, convenient and convincing generalization of the theory to all nodal curves; see \cite{Ran,Ee,EM,ES,Oss1,Oss2,ENR,ESVI,ESVII} for a sample. Attempts were even made using Tropical Geometry (the connection explained in \cite{OssT}), which have nonetheless led to new applications; see \cite{BN,AB,CDP,JP1,JP2} for a sample.

In the 2000's Ossermann \cite{Oss1} argued for another approach for curves of compact type, which is more functorial and works in positive characteristic. Instead of taking limits of linear series whose multidegrees were focused on a component, he considered the collection of all the limits that had effective multidegrees, and verified it is an \emph{Osserman limit linear series}. Over our curve $X$, an Osserman limit linear series of rank $r$ and degree $d$ is a collection of $d+1$ linear series $\mathfrak g=\{(L^{(0)},V^{(0)}),\dots,(L^{(d)},V^{(d)})\}$ of rank $r$ and degree $d$, with $\deg(L^{(0)}_{|Y})=d$ and, for each $i$,
\begin{enumerate}
    \item $\big(L^{(i)}_{|Y}\,,\, L^{(i)}_{|Z}\big)=\big(L^{(0)}_{|Y}(-iP)\,,\, L^{(0)}_{|Z}(+iP)\big)$;
    \item $\varphi^i(V^{(i)})\subseteq V^{(i+1)}$ and $\varphi_i(V^{(i+1)})\subseteq V^{(i)}$, where $\varphi^i\colon L^{(i)}\to L^{(i+1)}$ and $\varphi_i\colon L^{(i+1)}\to L^{(i)}$ are the unique (modulo scalar) nontrivial maps.
\end{enumerate}
Osserman \cite{Oss1}, Thm.~5.3, p.~1178, went on to produce a projective moduli space $G^{r,\text{Oss}}_d(X)$ parameterizing those $\mathfrak g$. He worked only over our curve $X$, but later \cite{Oss2} extended the theory to all nodal curves.

In \cite{E-Oss} Osserman and the first author showed another advantage of Osserman's approach. They showed how to produce a collection $\mathbf P(\mathfrak g)$ of degree-$d$ divisors of $X$ out of an Osserman limit linear series $\mathfrak g=\{(L^{(0)},V^{(0)}),\dots,(L^{(d)},V^{(d)})\}$ over $X$, and showed it coincides with the collection of limits of divisors if $\mathfrak g$ arises from a degeneration to $X$, as long as $\mathfrak g$ is \emph{exact}, a term coined by Osserman in \cite{Oss1}, Def.~A10, p.~1197, meaning that the induced sequences 
$$
V^{(i)}\to V^{(i+1)}\to V^{(i)}\to V^{(i+1)}
$$
are exact for all $i$; see \cite{E-Oss}, Def.~4.1, p.~83 and Thm.~5.2, p.~90. (Exactness is also necessary by \cite{ENR}, Prop.~5.4, p.~846.)
%{\color{red} Reformular a frase pra incluir se e seomente se.}
This does not hold for Eisenbud and Harris limit linear series, unless they are refined, in which case they are essentially exact Osserman limit linear series; see \cite{Oss1}, Cor.~6.8, p.~1189 and Rmk.~6.16, p.~1192. (This correspondence is not true for curves with more components; see \cite{M}.) But exactness is an open condition. It turned out then that Osserman moduli space $G^{r,\text{Oss}}_d(X)$ is projective, but only an open subset of it --- though dense for general $X$ by Liu \cite{Liu} --- parameterizes enough data to determine the limits of divisors.

Degenerating linear series along a degenerating family of curves do yield at the limit exact Osserman limit linear series, as long as the total space of the family is smooth; see \cite{E-Oss}, \S~5, p.~90. For our curve $X$ as limit curve, smoothness is controlled by a single number $\delta$, called the \emph{singularity degree} at $P$, smoothness corresponding to $\delta=1$. We \cite{ENR} have realized that one could replace a nonexact limit linear series by an exact one, by passing to the semistable reduction of the family, albeit one that generalized Osserman limit linear series, by allowing for generalized linear series and for a number of them in the collection depending on $\delta$. Following Osserman, we \cite{ENR}, Prop.~3.2, p.~834, constructed a projective variety $G^r_{d,\delta}(X)$ parameterizing what we termed \emph{level-$\delta$ limit linear series}; see \cite{ENR}, Def.~3.1, p.~833, or Section~\ref{deltalls} for a generalization. Not all of them are exact, and as before only the exact ones yield the right collection of divisors; see \cite{ENR}, Thms.~5.5~and~5.6, p.~846. So, again, we were stuck with a quasi-projective variety parameterizing enough data.

As promised in \cite{ENR}, p.~829, we are able to put together all the exact loci of the varieties $G^r_{d,\delta}(X)$ for all $\delta$, and produce, under appropriate identifications, a projective variety $G^r_d(X)$; see Theorem~\ref{allevels}. We realized that the collection of linear series in a limit linear series $\mathfrak g$ of any level can be put together as a $\cc$-invariant family of generalized linear series over $X$ parameterized by a two-punctured semistable chain of rational curves $(S,N_0,N_{\infty})$ if and only if $\mathfrak g$ is exact, a fact essential in the proof of our Theorem~\ref{allevels}. The two punctures $N_0$ and $N_{\infty}$ are on the initial and final curve of the chain, and by $\cc$-invariant we mean invariant under an action of $\cc$ on $(S,N_0,N_{\infty})$. Of course, such a family corresponds to one parameterized by the quotient stack $[S-\{N_0,N_{\infty}\}/\cc]$. The latter family is \emph{maximal} if it contains linear series with sections on invertible sheaves of all possible nonnegative degrees on $Y$ (and on $Z$) and \emph{primitive} if the family is not the pullback of a family over a chain with less components. It is interesting to observe that $[S-\{N_0,N_{\infty}\}/\cc]$ is the limit of a point! 

We avoid stacks after this introduction. In this paper, the families we consider, which we do call \emph{continuous linear series}, are defined as $\cc$-invariant families of generalized linear series over two-punctured semistable chains of rational curves; see Definition~\ref{inv-const}. The notion of maximality is embedded in Definition~\ref{tgls} of twisted generalized linear series, and primitiveness is everywhere nonconstancy. Definition~\ref{thefunctor} is that of a functor $\mathfrak{G}_d^r(X)$ parameterizing families of continuous linear series on $X$, and Theorem~\ref{funct-equiv} states that the functor is represented by a projective scheme $G^r_d(X)$. Finally, Theorem~\ref{allevels} states that every point on $G^r_d(X)$ can be viewed as a level-$\delta$ limit linear series for some $\delta$.

In Proposition~\ref{betaclass} we show that the chain of rational curves $S$ parameterizing a continuous linear series can be embedded as a $\cc$-invariant subscheme of a certain fiber $H^r_d(X,L)$ of a $\cc$-invariant map $H^r_d(X)\to J$ constructed in Section~\ref{gls}, where $J$ is a component of the Picard scheme of $X$ and $H^r_d(X)$ is a scheme parameterizing generalized linear series. Furthermore, the fiber $H^r_d(X,L)$ is a $\cc$-invariant subscheme of a product $\text{Grass}\big(r+1,W_1\oplus W_2\big)\times T$, where $T$ is the rational chain of $d+1$ rational curves; see Section~\ref{tG}. Were it not for $T$, we would be in the context of work by Giansiracusa and Wu~\cite{GW22}, who generalized seminal work by Kapranov~\cite{Kap-chow, GKZ}. In fact, stripping a level-$\delta$ limit linear series of its sheaves, we are left with a linked sequence of vector spaces, using the terminology in \cite{ENR}, Def.~3.5, p.~835, which can be regarded as a complete colineation; see \cite{Tha} and \cite{VE}. Thaddeus \cite{Tha} showed how the moduli space of complete colineations can be realized as a Chow, Hilbert or Mumford quotient of the action of $\cc$ on a Grassmannian. Our approach is that of constructing a Hilbert quotient, but instead of simply for the action of $\cc$ on the Grassmannian $\text{Grass}\big(r+1,W_1\oplus W_2\big)$, the action on its product by $T$.

The crux of the paper is the proof of Theorem~\ref{chainmaps},
%{\color{red} escrever o que ele diz} 
the longest of the paper. Using it, not only do we prove representability for our functor $\mathfrak{G}_d^r(X)$, thus constructing the fine moduli space of continuous linear series, but we show the space is an open and closed subscheme of the Hilbert scheme of $\cc$-invariant subschemes of appropriate Hilbert polynomial of $H^r_d(X)$, our Theorem~\ref{funct-equiv}. 

It might be possible to generalize the approach in this paper to all nodal curves $X$. The space of embedded sheaves $\mathbf R^{\mathbf J^{\mathbf b}}$ in \cite{AEc}, \S~2, would replace $J\times T$, and a (to-be-defined) scheme $\mathbf H^r_d(X)$ fibered by Grassmannians over $\mathbf R^{\mathbf J^{\mathbf b}}$ parameterizing generalized linear series would replace $H^r_d(X)$. In \cite{AEc}, Thm.~4.7, degenerations of degree-$d$ invertible sheaves are parameterized by certain subschemes $Y^{a,b}_{\ell,\mathfrak m}$ of $\mathbf R^{\mathbf J^{\mathbf b}}$, which must be images of (to-be-defined) subschemes $Z^{a,b}_{\ell,\mathfrak m}$ of $\mathbf H^r_d(X)$ if one considers degenerations of $(r+1)$-dimensional spaces of sections of these sheaves. One should then look into the piece of the Hilbert scheme of $\mathbf H^r_d(X)$ parameterizing the $Z^{a,b}_{\ell,\mathfrak m}$. The difficulty is that instead of chains of rational curves, we have the  $Y^{a,b}_{\ell,\mathfrak m}$, which are more generally unions of toric varieties induced by tilings of an Euclidean space by polytopes, as described in \cite{AEb}, \S~4.2. The corresponding subschemes $Z^{a,b}_{\ell,\mathfrak m}$ of $\mathbf H^r_d(X)$ should be as complicated, hence a theme for another paper. 

A logarithmic approach may be possible as well. It might replace handling the toric varieties in the approach delineated above. Carocci has recently announced at the Isaac Newton Institute joint work with Battistella and Wise which looks like a log version of the approach above. On the other hand, a discrete approach has recently appeared in work by Amini and the first author \cite{AE1}, which led to applications to canonical series \cite{AE2}.  %{\color{red} melhore!}

This work is based on the doctor thesis by the third author \cite{Rizzo}, where an approach using stable maps rather than Hilbert schemes was developed. We would like to thank Omid Amini, Margarida Melo, Brian Osserman, Marco Pacini, Eduardo Vital and Filippo Viviani for helpful discussions on the subject.

%{\color{red} Old introduction below.}

%The main goal of this paper is to introduce the notion of \textit{continuous linear series} on singular curves and to present the construction of a projective moduli space parameterizing them. This space glues together the loci of exact limit linear series of all levels (this terminology was introduced in \cite{Oss1} and \cite{ENR}), which are precisely those limits with desirable properties. Each collection of limit linear series can be interpreted as members of an equivariant family of linear series parameterized by chains of rational curves. A remarkable feature of this new moduli space is its projectivity, providing a compactification of Osserman exact limit linear series, and consequently, of Eisenbud-Harris refined limit linear series, using exact limit linear series. Furthermore, we establish a relationship between this moduli space and the ``discrete'' moduli spaces of level-$\delta$ limit linear series on the curve, for each $\delta$ (see \cite{ENR}).

%Our construction has two prominent features. First, we provide an elegant and efficient way to build a single moduli space taking into account degenerations coming from families with any possible degree of singularity. Second, we obtain a novel functorial version of Hilbert's quotient, as a moduli space, providing new insights on the theory of \textit{moduli spaces of complete collineations}, as presented by Thaddeus in \cite{Tha}. Conversely, from this equivalence with Hilbert quotients, our construction of continuous linear series acquires an interpretation as a moduli space constructed via GIT theory.

%\subsection{Context}
%Eisenbud and Harris' theory of limit linear series \cite{EH1}, is a powerful framework for studying degenerations of linear series along families of smooth curves degenerating to singular curves of compact type. It has been used to answer many important questions about curves and their moduli spaces. Indeed, since its introduction the theory has found extensive applications, as demonstrated in, e.g., \cite{EH2}, \cite{EH3} and \cite{EH4}. Subsequent research by authors such as Cotterill \cite{Cot}, Tarasca \cite{Tar}, Castorena et al \cite{Cast}, Farkas et al \cite{Fark1, Fark2} has further expanded its scope. On the other hand, generalizations of the theory have been proposed by Ran \cite{Ran}, even predating Eisenbud and Harris's work. These efforts were independently rediscovered by the first author \cite{Ee} and applied in \cite{EM} and \cite{ES}. Moreover, Tropical Geometry has offered a new approach, pioneered by Baker and Norine \cite{BN} and Amini and Baker \cite{AB}, leading to substantial progress, as exemplified by \cite{CDP}, \cite{JP1, JP2}.

%In a significant breakthrough, Osserman in \cite{Oss1} achieved substantial progress towards generalizing limit linear series. He introduced a novel, functorial approach that proved more amenable to generalizations. Specifically, he constructed a projective moduli space $G_d^{r,\text{Oss}}(X)$ parameterizing all (bi)degrees $d$, rank $r$ limit linear series on a singular curve $X$, where $X$ is the union of two smooth components $Y$ and $Z$ intersecting transversally in a unique node $P$. Interesting applications of this perspective, as seen in \cite{Liu}, \cite{M}, \cite{JP2}, \cite{OssT} have opened new avenues for studying limit linear series and their geometric consequences. While Osserman's work in \cite{Oss1} has inspired generalizations of limit linear series to more general singular curves (see \cite{Oss2}, \cite{M}, \cite{ESVI, ESVII}), these generalizations often lack a corresponding moduli space construction. Recent developments, such as those presented in \cite{AEa,AEb,AEc,AE1}, are introducing new ideas to address the lack of a moduli space construction in many generalizations of limit linear series. These advancements aim to pave the way for constructing moduli spaces for limit linear series on general singular curves, count on the foundational work of Osserman and subsequent research. Given the connection of our work to the constructions and applications in \cite{Oss1} and \cite{ENR}, we focus our discussion in this paper on the ``toy case'' of the nodal curve $X$.

%When studying Osserman's work, it is inevitable to compare Osserman and Eisenbud-Harris limit linear series. Osserman limit linear series are recognized as carrying more information. More precisely, there exists a forgetful map $G^{r,\text{Oss}}_d(X)\to G^{r,E-H}_d(X)$, where $G^{r,E-H}_d(X)$ denotes the projective moduli space of limits of Eisenbud--Harris type. This map is an isomorphism over the \textit{refined locus} of $G^{r,E-H}_d(X)$ and thus $G^{r,\text{Oss}}_d(X)$ can be considered as a distinct natural compactification of this locus; see \cite{Oss1}, Section~6, p.~1183. A fundamental limitation of the Eisenbud-Harris approach is the lack of a direct connection between its limit linear series and ``schematic limits of divisors''. In contrast, Osserman and the first author explored this connection, yielding significant applications by defining a natural subscheme $\p(\mathfrak{g})$ of the symmetric product $X^{(d)}$ associated to a ``special'' Osserman limit linear series $\mathfrak{g}$ on $X$. We will elaborate on this idea below.

%Among Osserman limit linear series there are some special ones which he called \emph{exact}. They form an open subscheme $G_d^{r,\ast}(X)$ of $G^{r,\text{Oss}}_d(X)$. Exact limit linear series enjoy some important properties:
%\begin{itemize}
%\item If $X$ is general, they are dense among all limit linear series, by \cite{Liu}, Cor.~1.4, p.~4034.
%\item All limits of linear series along families of smooth curves degenerating to $X$ are exact, \emph{as long as} the total space of the family is regular; see \cite{E-Oss}, Section~5, p.~90.
%\item The subscheme $\p(\mathfrak{g})$ has the expected Hilbert polynomial, as if $\mathfrak g$ were a limit, and in this case $\p(\mathfrak{g})$ is actually the limit of the schemes corresponding to the linear series degenerating to $\mathfrak g$; see \cite{E-Oss}, Thms.~4.3 and 5.2.
%\end{itemize}

%In a more general setting, the authors introduced in \cite{ENR} the notion of {\em level-$\delta$ limit linear series} on $X$. As in Osserman's work, there are some special level-$\delta$ limit linear series, which we also called {\em exact}. These exact level-$\delta$ limit linear series arise as limits of linear series along \textit{non-regular smoothings of} $X$, that is, families of smooth curves degenerating to $X$ whose total space has singularity degree $\delta$ at $P$. The authors constructed projective moduli schemes $G^r_{d,\delta}(X)$ parameterizing level-$\delta$ limit linear series. Essentially, while $G^{r,\text{Oss}}_d(X)$ is the appropriate space to describe limits of linear series of smooth curves degenerating to $X$ along regular smoothings, the space of level-$\delta$ limit linear series, $G^r_{d,\delta}(X)$, provides the correct framework for describing limits along non-regular smoothings.

%Analogous to the comparison between Eisenbud-Harris and Osserman limits, level-$\delta$ limit linear series encode richer information than Osserman limit linear series. A forgetful map $G^r_{d,\delta}(X)\to G^{r,\text{Oss}}_d(X)$ exists. It turns out that, for each $\delta$, we may view $G^r_{d,\delta}(X)$ as a compactification of the moduli space of Osserman {\em exact} limit linear series. As in the level-$1$ case, we associate to each level-$\delta$ limit linear series $\mathfrak g$ a subscheme $\mathbb P(\mathfrak g)$ of $X^{(d)}$. Furthermore, we showed that if $\mathfrak g$ is exact, then $\p(\mathfrak{g})$ has the expected Hilbert polynomial, as if $\mathfrak g$ were a limit. Moreover, in this case, $\p(\mathfrak{g})$ accurately represents the limit of the schemes corresponding to the linear series degenerating to $\mathfrak g$.

%According to Theorem 4.1 and Example 4.1 in \cite{ENR} we need a method to establish an equivalence relation among the various $G_{d,\delta}^{r,\ast}(X)$ to reduce their size as $\delta$ grows. Additionally, we aim to concentrate within these same loci all the exact limit linear series of all levels, since are precisely those limits with good properties. The concept of \textit{continuous linear series} on $X$ elegantly and efficiently provides this equivalence relation. The core idea is remarkably simple: Osserman's exactness condition precisely characterizes the necessary and sufficient condition for the limits of two consecutive linear series in an Osserman limit linear series to coincide. This observation enables us to represent the discrete data defining exact limit linear series as a single continuous object: a family of linear series parameterized by a chain of rational smooth curves. Families of linear series over chains in the same or different levels are identified by requiring that their stabilizations coincide. This principle underlies the concept of continuous linear series.

%In conclusion, our first goal is to construct a moduli space that unifies the exact loci $G_{d,\delta}^{r,\ast}(X)$ for all values of $\delta$. This new moduli space, denoted by $G^r_d(X)$, parametrizes continuous linear series on $X$. It is a projective space, giving rise to a compactification of the locus of Osserman exact limit linear series, and thus of Eisenbud-Harris refined limit linear series, by exact limit linear series. From a broader perspective, this paper aims to introduce a novel notion of limit linear series on $X$, partially motivated by the desire to construct a meaningful resolution of a natural rational Abel-type map.
%\begin{eqnarray*}
%\mathbb{P}\colon G_d^{r,\text{Oss}}(X) & \dashrightarrow &
%\text{Hilb}_{A_d}^{\binom{r+s+t}{r}},\\
%\mathfrak{g} & \longmapsto & \mathbb{P}(\mathfrak{g})
%\end{eqnarray*}
%which is only defined on $G_d^{r,\ast}(X)$ (see \cite{E-Oss}). 

%Consequently, our second goal is to establish a geometric and comparative framework between continuous linear series and their ``discrete'' counterparts, namely, level-$\delta$ limit linear series of degree $d$ on $X$, for any value of $\delta$. This connection can be interpreted as a natural and meaningful partition of the moduli space of continuous linear series by numerical discrete data. In essence, it provides a mechanism for understanding each continuous limit linear series from a more simplified and manageable perspective.

%Our construction of the moduli space $G^r_d(X)$ of continuous linear series on $X$ offers a novel interpretation of Hilbert's quotient, which parameterizes complete collineations (see \cite[Prop. 7.3]{Tha}).Specifically, there exists a one-to-one correspondence between (classes of) complete collineations of depth $d$ and (classes of) exact linked nets of length $d$ (see \cite[Chap. 6]{VE}). Consequently, our Hilbert scheme, which parameterizes torus-invariant subschemes of the generalized linear series space on $X$ can be viewed as the representative projective scheme of a functor associated with the moduli problem of finding a meaningful closure of the space of complete collineations. We interpret continuous linear series as the continuous counterpart of a complete collineation. Furthermore, since the presentation of Hilbert's quotient in \cite{Tha} uses GIT theory, this connection with the moduli space of continuous linear series on $X$ provides an interpretation of this construction within the framework of GIT theory.

%\subsection{Roadmap and overview of the results}
%As in \cite{Oss1} and \cite{ENR}, we work with the ``toy case'' ofa nodal curve $X$ with only two components $Y$ and $Z$, which are smooth and meet transversally in a unique point $P$.

%The second section focuses on the study of \textit{continuous linear series} on the curve $X$. Informally, a continuous linear series on $X$ can be understood as a continuous representation of data that parameterizes \textit{twisted generalized linear series} on $X$.  These twisted generalized linear series are defined by a chain of rational curves associated with a \textit{family of twisters} on $X$. The concept of continuous linear series arises from the need to establish an equivalence relation among the increasingly numerous connected components of exact level-$\delta$ limit linear series as $\delta$ approaches infinity. We conclude this section by introducing the functor of continuous linear series of degree $d$ and rank $r$ on $X$.

%Section 3 focuses on the construction of a moduli space, denoted by $\mathcal{H}^{r}_d(X)$, for \textit{generalized linear series} on $X$. This moduli space parameterizes pairs $(\L,V)$, where $\mathcal{L}$ is a torsion-free, rank-$1$ sheaf on $X$ of degree $d$ and $V$ is a $(r+1)$-dimensional vector subspace of $\Gamma(X,\mathcal{L})$. In this section, we present two useful results. The first is the equivariant immersion of $\mathcal{H}^{r}_d(X)$ into a projective space. This allows us to define an induced natural action on the Hilbert scheme and identify the invariant subscheme under this action. The second result is the study of invariant curves on the Grassmannian under torus actions. Both of these results are crucial for proving the representability of the functor defining the continuous linear series of degree $d$ and rank $r$ on $X$. They enable us to associate each continuous linear series with a corresponding object within a specific invariant Hilbert scheme.

%In Sections 4 and 5, we present the first main result of this paper: the representability of the functor of continuous linear series of degree $d$ and rank $r$ on $X$ by a projective space, denoted by $G_d^r(X)$. Two crucial technical results, \textit{embeddedness} and \textit{rigidity}, are essential for establishing this representability. These concepts can be informally interpreted as follows: each continuous linear series on $X$ corresponds to an invariant torus action subscheme of $\mathcal{H}^{r}_d(X)$ and viceversa. 

%It is noteworthy that establishing embeddedness and rigidity, and consequently proving representability, presents a significant challenge when attempting to extend these results to more general curves. Omini and Esteves (see \cite{AEc, AE1}) have generalized the deformation theory of line bundles in an effort to extend, for example, the construction and existence of a moduli space for novel notions of limit linear series on any singular curve. Despite their efforts, such a construction remains elusive, likely due to the difficulties associated with extending these results to more intricate combinatorial objects like polytopes.


%Finally, in Section 6, we establish the second main theorem of this paper: a correspondence between the distinct loci of level-$\delta$ exact limit linear series and their continuous counterparts. This section introduces a more general setting than that presented in \cite{ENR}. Specifically, a more general notion of level-$\delta$ limit linear series is introduced with a combinatorial perspective, laying the groundwork for the next section, where we establish relations with the theory of divisors.

%This section introduces a more general framework for level-$\delta$ limit linear series extending beyond the framework presented in \cite{ENR}, by incorporating a combinatorial perspective. This generalized framework provides a foundation for future research, particularly in connection with emerging perspectives on the construction of such moduli spaces (see \cite{AE1}).

%Part of this work appeared in the third author's doctoral thesis at IMPA \cite{Rizzo}. We would like to thank Margarida Melo, Brian Osserman, Marco Pacini and Filippo Viviani for helpful discussions on the subject.

\section{Continuous linear series}\label{sec:stablells}

\textbf{Fix} an algebraically closed field $k$. All our schemes will be $k$-schemes. For each scheme $T$, denote by $T(k)$ its set of $k$-points. 

A reduced projective scheme of pure dimension $1$ is called a \emph{curve}. A coherent sheaf $L$ on a curve $X$ is called \emph{torsion-free, rank-one} if its associated points are the generic points of $X$, over which the sheaf has rank 1. Its \emph{degree} is $\chi(L)-\chi(\mathcal O_X)$.

Let $\pi:\mathcal{X}\rightarrow B$ be a flat projective map whose fibers over $k$-points are curves. We call $\pi$ or $\mathcal{X}/B$ a \emph{family of curves}. Given $b\in B(k)$, a subscheme $S\subseteq\mathcal X$ and a coherent sheaf $\mathcal H$ on $\mathcal{X}$, we let $S_b$ denote the intersection of $S$ with the fiber $\mathcal X_b$ of $\pi$ over $b$, and $\mathcal H_b$ the restriction of $\mathcal H$ to the fiber.

A $B$-flat coherent sheaf $\mathcal L$ on $\mathcal{X}$, for short a \emph{sheaf} on $\mathcal X/B$, is called \emph{relatively torsion-free, rank-one on $\mathcal{X}/B$} if its restriction $\mathcal L_b$ to the fiber $\mathcal{X}_b$ is a torsion-free rank-one sheaf on $\mathcal{X}_b$ for each $b\in B(k)$. It is called of (\emph{relative}) \emph{degree} $d$ if $\mathcal L_b$ has degree $d$ for each $b$. An invertible sheaf on $\mathcal X$ is relatively torsion-free, rank-one on $\mathcal{X}/B$.

A \emph{family of generalized linear series on $\mathcal X/B$} is the pair $(\L,\V)$ of a relatively torsion-free, rank-one sheaf $\mathcal L$ on $\mathcal{X}/B$ and a locally free subsheaf 
$\mathcal V\subseteq\pi_*\mathcal L$ such that the composition of the restriction with the base-change map,
$$
\mathcal V_b\to\pi_*(\mathcal L)_b\to H^0(\mathcal X_b,\mathcal L_b),
$$
is injective for each $b\in B(k)$. It is said to have \emph{rank} $r$ if $\mathcal V$ has constant rank $r+1$, and (\emph{relative}) \emph{degree} $d$ if $\mathcal L$ has relative degree $d$. 

If $B=\text{Spec}(k[[t]])$, the data of $\mathcal{X}/B$ and an isomorphism from the closed fiber of $\mathcal{X}/B$ to a curve $X$ is called a {\it regular smoothing of} $X$ if $\mathcal{X}$ is regular and the generic fiber of $\mathcal{X}/B$ is smooth over the field of Laurent series $k\{\{t\}\}$. We will often assume the isomorphism to $X$ is given without mentioning it.

Let ${\mathbf G}_{\mathbf m}:=\text{Spec}(k[t,t^{-1}])$, the one-dimensional torus over $k$. If ${\mathbf G}_{\mathbf m}$ acts on a $k$-scheme $T$, we will let $x\ast t$ denote the action of each $x\in{\mathbf G}_{\mathbf m}(k)$ on 
each $t\in T(k)$.

\textbf{Fix} a curve $X$. \textbf{Assume} $X$ has two irreducible components, denoted $Y$ and $Z$, which intersect at a single point, denoted $P$, and are smooth and intersect transversally, so that $P$ is a node of $X$.

For each $k$-scheme $B$ and integers $i,j$, an invertible sheaf $\mathcal L$ on $X\times B$ will be said to have \emph{relative bidegree $(i,j)$} on $X\times B/B$ if the restrictions of $\mathcal L$ to $Y\times b$ and $Z\times b$ have degrees $i$ and $j$, respectively, for each $b\in B(k)$. 

\subsection{The family of twisters}\label{sec:twisters}

A \emph{chain} (\emph{of rational curves}) over $k$ is a connected curve $T$ whose components are smooth, rational and can be ordered, $T_0,\dots,T_d$, in such a way that $T_i$ intersects $T_j$ only if $|i-j|\leq 1$, and the intersection is transverse, and at a single point. The ordering is unique up to reversing it. The components $T_0$ and $T_d$ are called \emph{extreme}. Given two $k$-points $N_0\in T_0(k)$ and $N_\infty\in T_d(k)$ on the smooth locus of $T$, we call $(T,N_0,N_\infty)$ a \emph{two-pointed chain}. From now on, we will only consider two-pointed chains, and we will refer to them as just chains. The nodes of $T$ and $N_0,N_{\infty}$ are called the \emph{special} points of the chain. For each $i=0,\dots,d$, denote by $T^*_i$ the set of nonspecial points of $T$ on $T_i$ and $T^*=\bigcup T_i^*$. There are actions of $\cc$ on $T$ leaving $T_i$ and the special points of $T$ on each $T_i$ fixed, and making $T_i^*$ a torsor for $i=0,\dots,d$. We call such an action a \emph{chain action}.

\textbf{Fix} a nonnegative integer $d$. \textbf{Fix} a chain $(T,N_0,N_\infty)$ having $d+1$ components. \textbf{Fix} an ordering $T_0,\dots,T_d$ of the components such that $T_i$ intersects $T_j$ only if $|i-j|\leq 1$. For each $i=1,\dots,d$, let
$N_i$ be the intersection of $T_{i-1}$ and $T_{i}$. For convenience, put $N_{d+1}:=N_\infty$. 

For each $i=0,\dots,d$, \textbf{fix} an isomorphism $T_i\xrightarrow{\cong}\IP^1$ such that, denoting by $\xi^i\colon\O_{T_i}^2\to\O_{T_i}(1)$ the tautological quotient of $T_i$, we have that $\xi^i|_{N_i}(\O_{N_i}\oplus 0)=0$ and $\xi^i|_{N_{i+1}}(0\oplus\O_{N_{i+1}})=0$. \textbf{Fix} $E_i\in T_i(k)$ such that $\xi^i|_{E_i}$ is the cokernel of the diagonal $\mathcal O_{E_i}\hookrightarrow\O_{E_i}^2$. There is a unique chain action of $\cc$ on $(T,N_0,N_\infty)$ such that $\xi^i|_{x\ast Q}=\xi^i|_Q(x,1)$, where $(x,1)$ is the natural endomorphism on $\O_{Q}^2$, for each $x\in\cc(k)$ and $Q\in T_i(k)$. \textbf{Fix} this action. Notice that for each $i=0,\dots,d$ and each $Q\in T_i^*(k)$, we have that %<$1\ast E_i=E_i$ and 
$\lim_{x\to 0}x\ast Q=N_i$, whereas $\lim_{x\to\infty}x\ast Q=N_{i+1}$.

\textbf{Fix} isomorphisms $\O_Y(P)|_P\xrightarrow{\cong}\O_P$ and $\O_Z(P)|_P\xrightarrow{\cong}\O_P$, and from these construct isomorphisms $\O_Y(iP)|_P\cong\O_P$ and $\O_Z(jP)|_P\cong\O_P$ for all integers $i$ and $j$. Using them we obtain a surjection on $X$ for each $i=0,\dots,d$:
\[
\tau^i\colon \O_Y(-iP)\oplus\O_Z(iP) \longrightarrow
\O_Y(-iP)|_P\oplus\O_Z(iP)|_P \xrightarrow{\, \, \cong\, \, } \O_P^2.
\]
Let $\O_X^{(i)}$ be the kernel of $(1,-1)\tau^i$. 
Then $\O_X^{(i)}$ is an invertible sheaf on $X$ whose restriction to $Y$ 
is $\O_Y(-iP)$ and to $Z$ is $\O_Z(iP)$. Let $\tau^i_{T_i}$ be the pullback of
$\tau^i$ and $\xi^i_X$ that of $\xi^i$ to $X\times T_i$.
Identify $P\times T_i=T_i$, and consider
the composition $\xi^i_{X}\tau^i_{T_i}$. It is a surjection of sheaves on
$X\times T_i$; let $\mathcal F_i$ denote its kernel. Identifying $X\times R=X$ for each $R\in T_i(k)$, we see that $\F_i|_{X\times E_i}=\O_X^{(i)}$. Furthermore, 
\[
\mathcal F_i|_{X\times N_i}= \O_Y(-iP)\oplus\O_Z((i-1)P)\quad\text{and}\quad
\mathcal F_i|_{X\times N_{i+1}}=\O_Y(-(i+1)P)\oplus\O_Z(iP).
\]
As $\F_i|_{X\times N_{i+1}}=\F_{i+1}|_{X\times N_{i+1}}$ for $i=0,\dots,d-1$, there is a coherent subsheaf 
\[
\mathcal F\subset p^*_1\big(\O_Y\oplus\O_Z(dP)\big)
\]
with $T$-flat quotient, where $p_1\:X\times T\to X$ is the indicated projection, such that $\mathcal F|_{X\times T_i}=\mathcal F_i$ for each $i=0,\dots,d$. Then $\F$ is a relatively torsion-free sheaf of rank $1$ and degree $0$ on $X\times T/T$. Abstractly, $\F_i|_{X\times T_i^*}\cong\O_X^{(i)}\otimes\O_{T_i^*}$. 

The sheaf $\F$ is invariant under the action of $\cc$: For each $x\in\cc(k)$, each $i=0,\dots,d$ and $Q\in T_i$, we have
\begin{equation}\label{QcQ}
\F|_{X\times(x\ast Q)}=\Ker(\xi^i|_Q(x,1)\tau^i)=(x^{-1},1)\Ker(\xi^i|_Q\tau^i)=(x^{-1},1)\F|_{X\times Q}
\end{equation}
as subsheaves of $\O_Y(-iP)\oplus\O_Z(iP)$.

We say $\mathcal F$ is the \emph{family of twisters} associated to $X$ and $d$, and \textbf{fix} the notation.

\subsection{An alternative construction}

For the sake of completeness we sketch next an alternative construction of the sheaf $\F$ through a degeneration argument. This construction will not be used elsewhere in the paper.

Let $\widehat T$ be the chain obtained from $T$ by
adding an extra rational curve at each end: one,
denoted $T_{-1}$, intersecting $T_0$ transversally
at $N_0$, the other, denoted $T_{d+1}$, intersecting
$T_d$ transversally at $N_{d+1}$. View $T\subset\widehat T$.

Let $B:=\text{Spec}(k[[t]])$. Let $\pi\colon\mathcal{X}\rightarrow B$ (resp.~$\tau:\mathcal{T}\rightarrow B$) be a regular smoothing of $X$ (resp.~$\widehat T$). Such maps exist, as the universal deformation space of a nodal curve is regular. Form the threefold
$\mathcal{X}\times_B\mathcal{T}$. It is a
smoothing of the surface $X\times\widehat T$. However,
it fails to be regular exactly at $(P,N_i)$
for $i=0,\dots, d+1$. Following the ideas in ~\cite[\S~4.2]{CEP}, 
blow up $\mathcal{X}\times_B\mathcal{T}$ successively
along the (strict transforms of)
$Y\times T_{-1}$, $Y\times T_0$, \dots, $Y\times T_d$ in this order.
Denote by $\widetilde{\mathcal X}$ the resulting space,
and by $\mathcal Y_i$ and $\mathcal{Z}_i$ the strict transforms in
$\widetilde{\mathcal X}$ of $Y\times T_i$ and $Z\times T_i$ for each
$i=-1,\dots,d+1$. It follows that $\widetilde{\mathcal X}$
is regular and the $\mathcal Y_i$ and $\mathcal Z_i$ are Cartier divisors of
$\widetilde{\mathcal X}$. Let
$\psi\colon\widetilde{\mathcal X}\to\mathcal{X}\times_B\mathcal{T}$
denote the blowup map. In the sense of ~\cite[Lem.~5.3, p.~2953]{CEP},
the map $\psi$ is a semistable modification of 
each of the two projections of the fibered product
$\mathcal{X}\times_B\mathcal{T}$ onto its factors. In particular, 
the composition of the blowup with any of 
these projections is flat.

Let $C_i:=\psi^{-1}(P,N_i)$ for each $i=0,\dots,d+1$. Then the
$C_i$ are rational smooth curves. 
The strict transforms $\mathcal Y_{i-1}$ and $\mathcal Z_i$ contain $C_i$, 
whereas $\mathcal Y_{i}$ and $\mathcal Z_{i-1}$ intersect $C_i$
transversally at unique and distinct points. We have that 
$(\mathcal Y_j+\mathcal Z_j)\cdot C_i=0$ for each $j=-1,\dots,d+1$. Also, $\mathcal Y_j\cdot C_i=0$ if $j\neq i-1,i$, whereas $\mathcal Y_i\cdot C_i=1$ and $\mathcal Y_{i-1}\cdot C_i=-1$.

Set
$$
\mathcal{G}:=\O_{\widetilde{\mathcal{X}}}
\Big(-\mathcal Y_{-1}+\mathcal Y_1+2\mathcal
Y_2+\cdots+(d+1)\mathcal Y_{d+1}\Big).
$$
Then $\mathcal{G}|_{C_i}\cong\O_{C_i}(1)$ for each
$i=0,\dots,d+1$, and hence the sheaf $\mathcal G$ is $\psi$-admissible, in the sense of \cite[\S~3]{EP}. It follows from ~\cite[Thm. 3.1, p.~63]{EP} that $\psi_{\ast}\mathcal{G}$ is a relatively torsion-free, rank-1 sheaf of degree 0 on
$\mathcal{X}\times_B\mathcal{T}/\mathcal{T}$, whose
formation commutes with base change. The restriction $\psi_{\ast}\mathcal{G}|_{X\times T}$ is isomorphic to $\mathcal F$.

\subsection{Twisted generalized linear series}

\begin{defn}
Let $B$ be a scheme. Define the \emph{category $\mathcal C_B$ of families of chains over $B$} as follows:
\begin{enumerate}
    \item A \emph{family of chains over $B$} is the data of a family of curves $\pi\:\mathcal S\to B$ and sections $\mathcal N_0,\mathcal N_\infty\subset\mathcal S$ of $\pi$ such that $(\mathcal S_b,\mathcal N_{0,b},\mathcal N_{\infty,b})$ is a chain for each $b\in B(k)$. 
    \item A \emph{morphism} from a family of chains $(\pi\colon\mathcal S\to B,\mathcal N_0,\mathcal N_\infty)$ to another, $(\pi'\colon\mathcal S'\to B,\mathcal N'_0,\mathcal N'_\infty)$, is a $B$-morphism $\gamma\:\mathcal S\to\mathcal S'$ such that $\gamma^{-1}(\mathcal N'_i)=\mathcal N_i$ for $i=0,\infty$ and $\gamma_*[\mathcal S_b]=[\mathcal S'_b]$ for each $b\in B(k)$, where $[-]$ is the fundamental class in the Chow group.
\end{enumerate}
\end{defn}

\begin{defn} Let $B$ be a scheme. Let $p_1\: B\times T \to B$ be the projection. Define the \emph{category of families of chain maps to $T$ over $B$} as the category $\mathcal C_B(T)$ of morphisms in $\mathcal C_B$ to $(p_1,B\times N_0,B\times N_{\infty})$. 
\end{defn}

The category $\mathcal C_B(T)$ is the same as the category of maps 
$\gamma=(\gamma_1,\gamma_2)\:\mathcal S\to B\times T$ where
\begin{enumerate}
    \item $(\gamma_1,\gamma_2^{-1}(N_0),\gamma_2^{-1}(N_\infty))$ is a family of chains over $B$,
    \item $\gamma_{2*}[\mathcal S_b]=[T]$ for each $b\in B(k)$,
\end{enumerate}
which we call families of chain maps to $T$ over $B$ as well. In this approach, a morphism from a family of chain maps $\gamma\:\mathcal S\to B\times T$ to another, $\eta\:\mathcal S'\to B\times T$, is simply a $(B\times T)$-morphism $f\:\mathcal S\to\mathcal S'$ such that $f_*[\gamma_1^{-1}(b)]=[\eta_1^{-1}(b)]$ for each $b\in B(k)$.

If $B=\mathrm{Spec}(k)$, a family of chain maps to $T$ over $B$ is simply called a \emph{chain map} to $T$. If $f$ is a morphism from one family of chain maps to $T$ over $B$ to another, we say that the latter is \emph{intermediate} to the former; if $f$ is an isomorphism, we say they are \emph{equivalent}. Given a chain map to $T$, the chain action on the target $(T,N_0,N_\infty)$ lifts to a chain action on the source; the lifting may not be unique.

Let $B$ be a scheme. Let $\gamma=(\gamma_1,\gamma_2)\:\mathcal S\to B\times T$ be a family of chain maps. Let $\L$ be an invertible sheaf on
$X\times B$. Consider the induced maps:
\begin{equation}
\xymatrix{&X\times\mathcal S\ar[dl]_{(1,\gamma_1)}\ar[d]_{p_{\mathcal S}} \ar[dr]^{(1,\gamma_2)}&\\
X\times B & \mathcal S & X\times T,}
\end{equation}
where $p_{\mathcal S}$ is the projection. We will call the sheaf
$$
\L(\gamma):=\L\boxtimes\F:=(1,\gamma_1)^*\L\otimes(1,\gamma_2)^*\F
$$
on $X\times\mathcal S$ the \emph{family of twists of $\L$ along} $\gamma$. Notice that $\L(\gamma)=(1,\gamma)^*\L(1_{B\times T})$.

Notice that an isomorphism $g\:\L\to\L'$ induces an isomorphism $g(\gamma)\:\L(\gamma)\to\L'(\gamma)$ in the natural way. 

\begin{defn}\label{tgls} Let $\gamma\:\mathcal S\to B\times T$ be a family of chain maps. A \emph{family of twisted generalized linear series of degree $d$ and rank $r$ on $X$ along $\gamma$} is a pair $(\L,\V)$ consisting of:
\begin{enumerate}
\item an invertible sheaf $\L$ on $X\times B$ of relative bidegree $(d,0)$ on $X\times B/B$;
\item a subsheaf $\V\subseteq p_{\mathcal S\ast}\L(\gamma)$ such that $(\L(\gamma),\V)$ is a family of generalized linear series of rank $r$ on $X\times\mathcal S/\mathcal S$.
\end{enumerate}
When $B=\text{Spec}(k)$, we say that $(\L,\V)$ is a \emph{twisted generalized linear series} on $X$ along $\gamma$.
\end{defn}

Let $L$ be an invertible sheaf on $X$ of bidegree $(d,0)$. Let $\gamma\:S\to T$ and $\gamma'\:S'\to T$ be chain maps. Let $(L,\V)$ (resp.~$(L,\V')$) be a twisted generalized linear series on $X$ along $\gamma$ (resp.~$\gamma'$). Since $\F$ is a subsheaf of $(\O_Y\oplus\O_Z(dP))\otimes\O_T$ with $T$-flat quotient, it follows that 
\[
L(\gamma),L(\gamma')\subseteq \big(L|_Y\oplus L|_Z(dP)\big)\otimes\O_{S}
\]
with $S$-flat quotients, and hence 
\[
\V,\V'\subseteq \big(\Gamma(Y,L|_Y)\oplus\Gamma(Z,L|_Z(dP))\big)\otimes\O_{S}
\]
with $S$-flat quotients. We can thus compare the vector spaces $\V_{s}$ and $\V'_{s'}$ for $s\in S(k)$ and $s'\in S(k')$, and in particular, ask whether they are equal or not; see below. 

\begin{defn}\label{inv-const} Let $(L,\V)$ be a twisted generalized linear series on $X$ along a chain map $\gamma\: S\to T$. Let $C$ be an irreducible component of $S$. We say $(L,\V)$ is \emph{constant} along $C$ if $\gamma|_C$ is constant, and $\V_{s_1}=\V_{s_2}$ for all $s_1,s_2\in C(k)$. It is said to be \emph{invariant} along $C$ if the chain action on $(T,N_0,N_\infty)$ lifts to a chain action on $(S,\gamma^{-1}(N_0),\gamma^{-1}(N_\infty))$ such that $\V_{x\ast Q}=(x^{-1},1)\V_{Q}$ for each $x\in\cc(k)$ and $Q\in C^*(k)$. We say $(L,\V)$ is \emph{everywhere nonconstant} (resp.~\emph{everywhere invariant}) if it is not constant (resp.~it is invariant) along each irreducible component of $S$. It is said to be \emph{continuous} if it is everywhere nonconstant and invariant; in this case, we call $(L,\V)$ a \emph{continuous linear series}.
\end{defn}

Let $(\L,\mathcal{V})$ be a family of twisted generalized linear series on $X$ along a family of chain maps $\gamma\:\mathcal S\to B\times T$. Let $\gamma'\:\mathcal S'\to B\times T$ be another family of chain maps and $f\colon\mathcal S'\to\mathcal S$ a morphism from $\gamma'$ to $\gamma$. Then $(\L,f^*\V)$ is a family of twisted generalized linear series on $X$ along $\gamma'$, called the \emph{expansion} of $(\L,\V)$ along $f$. 

Also, for each map $b\:B'\to B$, considering the Cartesian diagram
\[
\begin{CD}
    \mathcal S'@>b'>> \mathcal S\\
    @V\gamma'VV @V\gamma VV\\
    B'\times T @>b\times 1>> B\times T,
\end{CD}
\]
we have that $\gamma'$ is a family of chain maps, and 
\[
b^*(\L,\V):=((1\times b)^*\L,(b')^*\V)
\]
is a family of twisted generalized linear series on $X$ along $\gamma'$ called the \emph{pullback} of $(\L,\V)$ under $b$. We call $(\L,\mathcal V)$ a \emph{family of everywhere nonconstant} (resp.~\emph{everywhere invariant}, resp.~\emph{continuous}) \emph{twisted generalized linear series on $X$ along $\gamma$} if $b^*(\L,\V)$ is an everywhere nonconstant (resp.~everywhere invariant, resp.~continuous) twisted generalized linear series for each $k$-point $b\:\text{Spec}(k)\to B$.

Finally, given an invertible sheaf $M$ on $B$, we have that 
\[
(\L,\V)\otimes M:=(\L\otimes p_2^*M,\V\otimes\gamma_1^*M)
\]
is another family of twisted generalized linear series on $X$ along $\gamma$, where $p_2\:X\times B\to B$ is the indicated projection.

\begin{defn}\label{tglsequiv} Two families of twisted generalized linear series $(\L,\V)$ and $(\L',\V')$ along families of chain maps $\gamma\colon\mathcal S\to B\times T$ and $\gamma'\colon\mathcal S'\to B\times T$ are said to be \emph{equivalent} if there are an isomorphism $f\colon\mathcal S'\to\mathcal S$ from $\gamma'$ to $\gamma$, an invertible sheaf $M$ on $B$ and an isomorphism $g\:\L'\xrightarrow{\cong}\L\otimes p_2^*M$ such that the induced isomorphism $g(\gamma')$ takes $\V'$ to $f^*\V\otimes(\gamma'_1)^*M$.
\end{defn}

Pullbacks and expansions of equivalent families are equivalent families.

\begin{defn}\label{thefunctor} The (contravariant) functor of
continuous linear series of degree $d$ and
rank $r$ on $X$,
$$
\mathfrak{G}_d^r(X)\colon (\text{Schemes})
\longrightarrow(\text{Sets}),
$$
is the functor that associates to each scheme $B$ the
set $\mathfrak{G}_d^r(X)(B)$ of equivalence classes of families of 
continuous (twisted generalized) linear series of degree $d$ and rank $r$ on $X$ along families of chain maps to $T$ over $B$.
\end{defn}

\section{Parametrization of generalized linear series}\label{gls}

\subsection{The scheme of generalized linear series}\label{gls-sch}

\textbf{Fix} $J:=\Pic^{(d,0)}_X$, the component of the Picard scheme of $X$ parameterizing invertible sheaves of bidegree $(d,0)$, that is, degree $d$ on $Y$ and $0$ on $Z$. Let $\P$ be a Poincar\'e sheaf on $X\times J$. Let $D$ be an effective divisor of $X$ supported away from $P$ and denote by $D'$ its pullback to $X\times J\times T$ under the
projection. Put 
\[
\P\boxtimes\F:=p_{1,2}^*(\P)\otimes p_{1,3}^*(\F)\quad\text{and}\quad
\mathcal E:=p_{2,3\ast}\left(\P\boxtimes\F\otimes\mathcal O(D')\right),
\]
where $p_{1,2}$, $p_{1,3}$ and $p_{2,3}$ are the projections onto the indicated factors:
\begin{equation}\label{Diag:H}
\xymatrix{
&X\times J\times T\ar[dl]_{p_{1,3}}\ar[d]_{p_{2,3}}\ar[dr]^{p_{1,2}}&\\
X\times T & J\times T & X\times J.}
\end{equation}
Assume $D$ has high enough degree, so that $R^1p_{2,3\ast}\left(\P\boxtimes\F\otimes\mathcal O(D')\right)=0$, and hence $\mathcal E$ is locally free. Put 
$\mathcal G:=p_{2,3\ast}(\P\boxtimes\F\otimes\mathcal O(D')|_{D'})$. Then $\mathcal G$ is locally free. Also, restriction induces a map $\alpha\colon\mathcal E\to\mathcal G$. 

Let $r$ be a nonnegative integer. Let $\alpha'\colon\mathcal E'\to\mathcal G'$ be the pullback of $\alpha$ to 
$\text{Grass}_{J\times T}(r+1,\E)$ and $\mathcal V'\subseteq\mathcal E'$ the universal locally free subsheaf of rank $r+1$. Denote by $H^{r}_d(X)$ the degeneracy scheme of $\alpha'|_{\mathcal V'}$, the locus where the map of bundles is zero. \textbf{Fix} the notation. Then $H^r_d(X)$ is a projective scheme, called the \emph{scheme of generalized linear series}.  It is a scheme over $J\times T$.

Let $\nu=(\nu_1,\nu_2)\: H^r_d(X)\to J\times T$ be the induced map. Let $\mathcal V$ be the restriction of $\mathcal V'$ to $H^r_d(X)$. It is a locally free subsheaf of rank $r+1$ of $p_*((1\times\nu)^*(\P\boxtimes\F))$, where $p:X\times H^{r}_d(X)\rightarrow H^{r}_d(X)$ is the projection. Furthermore, the pair 
$((1\times\nu)^*(\P\boxtimes\F),\mathcal V)$ is a family of generalized linear series of degree $d$ and rank $r$ on $X\times H^r_d(X)/H^r_d(X)$.

For further use, we denote by $H^r_d(X,L)$ the fiber of $\nu_1$ over each $L\in J(k)$, and call it the \emph{scheme of generalized linear series with sections in $L$}.  

%Furthermore, $H^{r}_d(X)$ represents the functor defined
%on the category of $T$-schemes, which associates to each
%$T$-scheme $S$ the set of pairs $(\L,\V)$, where $\L$ is
%an invertible sheaf on $X\times S$ of relative multidegree
%$(d,0)$ over $S$, and $(\L\ox\F,\V)$ is a family of generalized
%linear series of rank $r$ on $X\times S/S$, modulo the usual
%equivalence relation, where $(\L,\V)\sim(\L',\V')$ if and only
%if there is an invertible sheaf $\N$ on $S$ such that
%$\L'=\L\ox\N$ and $\V'=\V\ox\N$. Here, for simplicity, we
%identify $\F$ with $(Id, f)^*\F$, where
%$(Id, f):X\times S\to X\times T$ and $f:S\to T$ the natural map.

\subsection{The action.}\label{tG}

Since $J$ is isomorphic to an Abelian variety, every action of $\cc$ on $J$ is trivial. The action of $\cc$ on $T$ lifts thus to a unique action on $J\times T$. 

We claim the action lifts naturally
to one on $H^{r}_d(X)$. Indeed, $H^r_d(X)\subseteq \text{Grass}_{J\times T}(r+1,\E)$. Now,
$$
\P\boxtimes\F\subseteq p_{1,2}^*\Big(\P\otimes p_1^*\big(\mathcal O_Y\oplus\mathcal O_Z(dP)\big)\Big),
$$
where $p_1\colon X\times J\to X$ is the projection map. Put
$$
\W_1:=p_{2\ast}\big(\P\otimes p_1^*\mathcal O_Y(D\cap Y)\big)\text{ and }
\W_2:=p_{2\ast}\big(\P\otimes p_1^*\mathcal O_Z(dP+D\cap Z)\big),
$$
where $p_2\colon X\times J\to J$ is the projection. Then $\mathcal E\subseteq q_1^*(\W_1\oplus\W_2)$, where $q_1\colon J\times T\to J$ is the projection. It follows that 
$$
H^r_d(X)\subseteq \text{Grass}_{J}\big(r+1,\W_1\oplus\W_2\big)\times T.
$$
There is now a natural action of $\cc$ on $\text{Grass}_{J}\big(r+1,\W_1\oplus\W_2\big)$, induced by the automorphism $(x^{-1},1)$ on $\W_1\oplus\W_2$ for each $x\in\cc(k)$. Pair it with the action on $T$. Then $H^r_d(X)$ is invariant and the induced map $H^r_d(X)\to T$ is equivariant. 

We linearize the action on $H^r_d(X)$ as follows. 
First, $T\subseteq(\IP^1)^{d+1}$ with equations $x_iy_j=0$ for $i<j$, where $(x_i,y_i)$ are the coordinates of the $i$-th $\IP^1$ for $i=0,\dots,d$. The embedding is equivariant, for the product action of 
$\cc$ on $(\IP^1)^{d+1}$ given by $x\ast (x_i:y_i)=(xx_i:y_i)$ for each $x\in\cc(k)$ and $i=0,\dots,d$.

Second, consider the Pl\"ucker embedding 
$\text{Grass}_{J}\big(r+1,\W_1\oplus\W_2\big)\subseteq\IP_J(\W)$, where
$$
\W:=\bigwedge^{r+1}(\W_1\oplus\W_2)=\bigoplus_{i=0}^{r+1}\Big(\bigwedge^{i}\W_1\otimes\bigwedge^{r+1-i}\W_2\Big);
$$
it is naturally equivariant. Let $\mathcal A$ be a very ample invertible sheaf on $J$, enough ample that $\W\otimes\mathcal A$ be globally generated. Then 
$\IP_J(\W)\subseteq\IP(\Gamma(J,\W\otimes\mathcal A)^*)\times J$, again naturally equivariant. Furthermore, since $\mathcal A$ is very ample, it induces an embedding $J\subseteq\IP(\Gamma(J,\mathcal A)^*)$. 

Putting everything together, we have an equivariant embedding
$$
\Lambda\colon H^r_d(X)\longrightarrow\IP(\Gamma(J,\W\otimes\mathcal A)^*)\times\IP(\Gamma(J,\mathcal A)^*)\times(\IP^1)^{d+1},
$$
which composes with the Segre embedding to give an equivariant embedding of $H^r_d(X)$ in a projective space $\mathbf P^N$.

The action on $\mathbf P^N$ induces actions on the Hilbert schemes $\text{Hilb}_{H^r_d(X)}$ and $\text{Hilb}_{\mathbf P^N}$. We will denote by $\text{Hilb}^{\phi}_{H^r_d(X)}$ the open (and closed) subscheme of $\text{Hilb}_{H^r_d(X)}$ parameterizing subschemes of $H^r_d(X)$ with multivariate Hilbert polynomial 
$$
\phi(u,v,s_0,\dots,s_d):=u(r+1)+\sum_{i=0}^d s_i+1
$$
with respect to the embedding $\Lambda$. We will also denote by $\text{Hilb}^{\phi,\cc}_{H^r_d(X)}$ and $\text{Hilb}^{\cc}_{\mathbf P^N}$ the closed subschemes of 
$\text{Hilb}^{\phi}_{H^r_d(X)}$ and $\text{Hilb}_{\mathbf P^N}$ parameterizing invariant subschemes of 
$H^r_d(X)$ and of $\mathbf P^N$, respectively. They fit in the natural Cartesian diagram:
$$
\begin{CD}
\text{Hilb}^{\phi,\cc}_{H^r_d(X)} @>>> \text{Hilb}^{\cc}_{\mathbf P^N}\\
@VVV @VVV\\
\text{Hilb}^{\phi}_{H^r_d(X)} @>>> \text{Hilb}_{\mathbf P^N}.
\end{CD}
$$
Both $\text{Hilb}^{\phi,\cc}_{H^r_d(X)}$ and $\text{Hilb}^{\cc}_{\mathbf P^N}$ are examples of \textit{invariant Hilbert schemes}. A detailed study of these schemes and their properties can be found in \cite{Br}.

\subsection{Invariant curves in the Grassmannian}\label{sec:invt-Grass}

Let $W_1$ and $W_2$ be finite-dimensional vector spaces over $k$, and put $W:=W_1\oplus W_2$. Let $\iota_i\colon W_i\to W$ be the natural injection and $\rho_i\colon W\to W_i$ be the natural surjection for $i=1,2$.

Let $n$ be a positive integer and $G:=\text{Grass}\left(n,W\right)$. 
Recall that $G$ is a nonsingular projective variety
with $\text{Pic}(G)=\z$. The degree of a curve in $G$ is its intersection number with the ample generator of $\text{Pic}(G)$.

The action of $\cc$ on $W=W_1\oplus W_2$, given by $x\ast (u_1,u_2):=(x^{-1}u_1,u_2)$ for each $x\in\cc(k)$ and $u_i\in W_i$ for $i=1,2$, induces an action on $G$. For each $x\in\cc(k)$ and $V\in G(k)$, we have $x\ast V=(x^{-1},1)V$. 

For each $V\in G(k)$, let $h_V\colon\mathbb\IP^1\to G$ be the morphism sending $(x:1)$ to $x\ast V$ for each $x\in\cc(k)$. 

\begin{prop}\label{hV} Let $V\in G(k)$ be a nonfixed point. 
Let $V_{\infty,1}:=\iota_1^{-1}(V)$ and 
$V_{0,2}:=\iota_2^{-1}(V)$. Let 
$V_{\infty,2}:=\rho_2(V)$ and $V_{0,1}:=\rho_1(V)$. Then $V_{\infty,1}\subsetneqq V_{0,1}$ and $V_{0,2}\subsetneqq V_{\infty,2}$. Furthermore,
\begin{equation}\label{limlim}
\lim_{x\to 0}x*V=V_{0,1}\oplus V_{0,2}\quad \text{and} \quad
\lim_{x\to \infty}x*V=V_{\infty,1}\oplus V_{\infty,2}.
\end{equation}
Finally, $h_V$ is an isomorphism onto its image $\Gamma$, and
$$
\deg \Gamma=
\dim_k V_{0,1}-\dim_k V_{\infty,1}=
\dim_k V_{\infty,2}-\dim_k V_{0,2}.
$$
\end{prop}

\begin{proof} The first statement is clear, where the inequalities follow from the fact that $V$ is nonfixed. As for the remaining statements, let $V'\subseteq V$ such that 
$V=V_{\infty,1}\oplus V'\oplus V_{0,2}$. Then the projection map $V'\to W_2$ is an isomorphism to a subspace $V'_2\subseteq V_{\infty,2}$ such that 
$V_{\infty,2}=V'_2\oplus V_{0,2}$. Similarly, the projection map $V'\to W_1$ is an isomorphism to a subspace $V'_1\subseteq V_{0,1}$ such that 
$V_{0,1}=V_{\infty,1}\oplus V'_1$. Choose bases $B_{\infty,1}$ of $V_{\infty,1}$, $B'$ of $V'$ and $B_{0,2}$ of $V_{0,2}$. Consider the corresponding bases 
$B'_1$ of $V'_1$ and $B'_2$ of $V'_2$, images of $B'$ under the respective isomorphisms $V'\to V'_1$ and $V'\to V'_2$. Complete $B_{\infty,1}\cup B'_1$ to a basis $B_1$ of 
$W_1$ and $B_{0,2}\cup B'_2$ to a basis $B_2$ of $W_2$. The union $B:=B_1\cup B_2$ is a basis of $W$.

For each subset $I$ of $B$, let $n_{I,1}$ be the number of vectors in $B_1$
and $n_{I,2}$ the number of those in $B_2$. For $I$ with $|I|=n$, let $p_I$ be the corresponding Pl\"ucker coordinate of
$V$ and $p_I(y,z)$ that of $V_{(y:z)}$, the image of $(y:z)$ under $h_V$, for each $(y:z)\in\IP^1(k)$ with $yz\neq 0$.
Then $p_I(y,z)=z^{n_{I,1}}y^{n_{I,2}}p_I$.

It follows from our choice of bases that $p_I\neq 0$ if and only if 
$$
B_{\infty,1}\cup B_{0,2}\subseteq I\subseteq B_{\infty,1}\cup B'_1\cup B'_2\cup B_{0,2}
$$
and the intersections $I\cap B'_1$ and $I\cap B'_2$ are identified in $B'$ with disjoint subsets covering $B'$. It thus follows that $\{(n_{I,1},n_{I,2})\,|\,p_I\neq 0\}$ is equal to
$$
\{(\dim_k V_{\infty,1},\dim_k V_{\infty,2}),(\dim_k V_{\infty,1}+1,\dim_k V_{\infty,2}-1),\dots,(\dim_k V_{0,1},\dim_k V_{0,2})\},
$$
and hence that $h_V$ is an isomorphism to its image $\Gamma$, which has degree
$$
n-\dim_k V_{\infty,1}-\dim_k V_{0,2}=\dim_k V_{\infty,2}-\dim_k V_{0,2},
$$ 
proving the third statement of the proposition.

The second statement follows as well, as it is clear that 
$p_I(y,z)/z^{\dim_k V_{\infty,1}}y^{\dim_k V_{0,2}}$ is zero for $(y:z)=(0:1)$ for all $I$ but $I=B_{\infty,1}\cup B_{0,2}\cup B'_2$, which corresponds to the subspace $V_{\infty,1}\oplus V_{\infty,2}$ of $W$, yielding the second formula in \eqref{limlim}. The first follows analogously.
\end{proof}

\begin{prop} \label{hVV} Let $V,V'\in G(k)$. Assume that $\rho_1(V')=\iota^{-1}_1(V)$ (resp.~$\rho_2(V')=\iota_2^{-1}(V)$). Then the intersection $h_{V}(\IP^1)\,\cap\,h_{V'}(\IP^1)$ is empty or scheme-theoretically a point, the latter if and only if $\rho_2(V)=\iota_2^{-1}(V')$ (resp.~$\rho_1(V)=\iota_1^{-1}(V')$).
\end{prop}

\begin{proof} By the symmetry, we may assume $\rho_1(V')=\iota^{-1}_1(V)$. Now, 
$h_V(\IP^1)$ (resp.~$h_{V'}(\IP^1)$) is contained in the subGrassmannian of $G$ parameterizing $n$-dimensional subspaces of $\rho_1(V)\oplus\rho_2(V)$ (resp.~$\rho_1(V')\oplus\rho_2(V')$).
Since $\rho_1(V')=\iota^{-1}_1(V)$, the intersection of the two subGrassmannians is contained in the Grassmannian of $n$-dimensional subspaces of 
\begin{equation}\label{eq:VV'-intersection}
\iota^{-1}_1(V)\oplus(\rho_2(V)\cap\rho_2(V')).
\end{equation}
But this space has dimension at most $n$, and thus the Grassmannian is either empty or a point. It follows that $h_{V}(\IP^1)\,\cap\,h_{V'}(\IP^1)$ is either empty or scheme-theoretically a point. 

Assume the latter. Then, the point of intersection is the space \eqref{eq:VV'-intersection}, which has dimension $n$ only if $\rho_2(V)\subseteq\rho_2(V')$, in which case \eqref{eq:VV'-intersection} is equal to $\iota^{-1}_1(V)\oplus\rho_2(V)$. This space is clearly a point on $h_V(\IP^1)$. It will be one on $h_{V'}(\IP^1)$ if and only if $\rho_2(V)=\iota_2^{-1}(V')$. \end{proof}

\section{Representability}

\subsection{Embeddedness}

Let $\mathfrak g:=(\L,\V)$ be a family of twisted generalized linear series of degree $d$ and rank $r$ 
along $\gamma$, for a family of chain maps $\gamma=(\gamma_1,\gamma_2)\colon\mathcal S\to B\times T$. Recall $\mathcal P$ from Section~\ref{gls-sch}. Then 
there are an invertible sheaf $N$ on $B$ and a map $h\colon B\to J$ 
such that $(1_X\times h)^*\mathcal P\cong \L\otimes N$. Up to replacing $\mathfrak g$ by an equivalent family, we may assume $N=\mathcal O_B$. 
The induced map $f:=(h\gamma_1,\gamma_2)\colon\mathcal S\to J\times T$ satisfies 
$(1_X\times f)^*(\mathcal P\boxtimes\mathcal F)\cong\L(\gamma)$. The locally free subsheaf $\V\subseteq p_{\mathcal S*}\mathcal L(\gamma)$ 
of rank $r+1$ induces thus a map 
$f_{\mathfrak g}\colon\mathcal S\to H^r_d(X)\times B$.

\begin{prop}\label{betaclass} If $\mathfrak g$ is a family of continuous linear series, then $f_{\mathfrak g}$ is a closed embedding. Furthermore, its image $\mathcal S_{\mathfrak g}$ is $\cc$-invariant and has relative multivariate Hilbert polynomial $\phi$ over $B$. 
\end{prop}

\begin{proof} We may assume $B=\text{Spec}(k)$. Then $f_{\mathfrak g}$ factors through the fiber $H^r_d(X,\L)$ of $H^r_d(X)$ over $\L$, which is naturally an invariant subscheme of $G\times T$, where \[
G:=\text{Grass}(r+1,W)\quad\text{for }
W:=W_1\oplus W_2:=\Gamma(Y,\L|_Y)\oplus\Gamma(Z,\L|_Z(dP)).
\]
%More precisely,
%\begin{equation}\label{unioni}
%H^r_d(X,L)\subseteq\bigcup_{i=0}^d G_i\times T_i,
%\end{equation}
%where $G_i:=\text{Grass}(r+1,W^{(i)})$, for
%$W^{(i)}:=\Gamma(Y,L|_Y(-iP))\oplus\Gamma(Z,L|_Z(iP))$. 
Let $\iota_i\colon W_i\to W$ be the natural injection and $\rho_i\colon W\to W_i$ be the natural surjection for $i=1,2$.

As a map to $G\times T$, we have $f_{\mathfrak g}=(f,\gamma_2)$. 
Let $S_0,\dots,S_m$ be the sequence of irreducible components of $\mathcal S$, ordered in such a way that 
\begin{enumerate}
    \item $\gamma_2(S_0)=T_0$ and $\gamma_2(S_m)=T_d$,
    \item $S_i\cap S_j\neq\emptyset$ if and only if $|i-j|\leq 1$.
    \end{enumerate}
Let $Q_i:=S_{i-1}\cap S_i$ for $i=1,\dots,m$, and $Q_0,Q_{m+1}\in\mathcal S$ with $\gamma_2(Q_0)=N_0$ and $\gamma_2(Q_{m+1})=N_\infty$. 

Since $\mathfrak g$ is everywhere invariant, the image under $f_{\mathfrak g}$ of each $S_i$ is invariant, thus $\mathcal S_{\mathfrak g}$ is invariant. We will argue now that $f_{\mathfrak g}$ is an embedding. First of all, since $\mathfrak g$ is also everywhere nonconstant, for each $i=0,\dots,m$, either $\gamma_2|_{S_i}$ is an embedding, or there is $s\in S_i$ such that $f(s)$ is not fixed in $G$, and hence $f|_{S_i}$ is an embedding by Proposition~\ref{hV}. In any case, 
$f_{\mathfrak g}|_{S_i}$ is an embedding.

To show that $f_{\mathfrak g}$ is injective, by contradiction, let $R_1\in S_i$ and $R_2\in S_j$ for $i\leq j$ be distinct points on $\mathcal S$ such that $f_{\mathfrak g}(R_1)=f_{\mathfrak g}(R_2)$. Then $\gamma_2(R_1)=\gamma_2(R_2)$. Since $\gamma_2$ is a chain map, $\gamma_2(S_i)=\gamma_2(S_{i+1})=\cdots=\gamma_2(S_j)=N_\ell$ for some $\ell$. As proved above, $f|_{S_i},f|_{S_{i+1}},\dots,f|_{S_j}$ are embeddings. Furthermore, Proposition~\ref{hV} yields as well that 
\[
\rho_2(f(R_1))=\rho_2(f(Q_{i+1}))\subsetneqq\rho_2(f(Q_{i+2}))\subsetneqq\cdots\subsetneqq\rho_2(f(Q_{j}))\subseteq\rho_2(f(R_2)),
\]
where the last inclusion is strict if $R_2\neq Q_j$, and
\[
\rho_1(f(R_1))\supseteq\rho_1(f(Q_{i+1}))\supsetneqq\rho_1(f(Q_{i+2}))\supsetneqq\cdots\supsetneqq\rho_1(f(Q_{j}))=\rho_1(f(R_2)),
\]
where the first inclusion is strict if $R_1\neq Q_{i+1}$. Since $f(R_1)=f(R_2)$, we must have that either $j=i+1$ and $R_1=Q_j=R_2$, or $j=i$, and thus $R_1=R_2$ because $f|_{S_i}$ is injective. In any case, $R_1=R_2$, a contradiction.

We need now only prove that the differentials $d_{Q_j}f_{\mathfrak g}|_{S_{j-1}}$ and $d_{Q_j}f_{\mathfrak g}|_{S_j}$ at $Q_j$ have different images for $j=1,\dots,m$. 

First, their images are nonzero because $f_{\mathfrak g}|_{S_{j-1}}$ and $f_{\mathfrak g}|_{S_{j}}$ are embeddings. Second, their images are distinct, if $\gamma_2|_{S_{j-1}}$ or $\gamma_2|_{S_j}$ is not constant, because then the images of $d_{Q_j}\gamma_2|_{S_{j-1}}$ and $d_{Q_j}\gamma_2|_{S_j}$ are distinct. 

Assume now that $\gamma_2(S_{j-1})$ and $\gamma_2(S_j)$ are each a point, necessarily the same point, a node $N_\ell$ for some $\ell$. Then $f$ restricts to an injection on $S_{j-1}\cup S_j$, and to embeddings $f|_{S_{j-1}}$ and $f|_{S_j}$. We need only show now that the intersection at the unique point on $f(S_{j-1})\,\cap\,f(S_j)$ is transverse. Now, since $\g$ is everywhere invariant, there are isomorphisms $\theta_1\:\cc\to S_{j-1}^*$ and $\theta_2\:\cc\to S_j^*$ such that $\V_{\theta_i(x)}=(x^{-1},1)V_i$ for each $x\in\cc(k)$ and $i=1,2$, where $V_i:=\V_{\theta_i(1)}$. Also, $f\theta_i$ extends to the map $h_i:=h_{V_i}\colon\IP^1\to G$ defined in Subsection~\ref{sec:invt-Grass}. Then $h_1(\infty)$ and $h_2(0)$ are equal to the subspace $V:=\V_{Q_j}\subseteq W$. By Proposition~\ref{hV}, 
\[
h_1(\infty)=\lim_{x\to\infty}x\ast V_1=\iota_1^{-1}(V_1)\oplus\rho_2(V_1)\quad\text{and}\quad
h_2(0)=\lim_{x\to 0}x\ast V_2=\rho_1(V_2)\oplus\iota_2^{-1}(V_2),
\]
so $h_1(\infty)=h_2(0)$ means
\[
\iota_1^{-1}(V_1)=\rho_1(V_2)\quad\text{and}\quad \rho_2(V_1)=\iota_2^{-1}(V_2).
\]
But then Proposition~\ref{hVV} yields that $h_1(\IP^1)$ and $h_2(\IP^1)$, which are $f(S_{j-1})$ and $f(S_{j})$, intersect transversally. 

It remains to show that $\mathcal S_{\mathfrak g}$ has multivariate Hilbert polynomial $\phi$. First of all, $\mathcal S_{\mathfrak g}$ is a curve, so its multivariate Hilbert polynomial $p_{\mathcal S_{\mathfrak g}}(u,v,s_0,\dots,s_d)$ is linear. Also, $\mathcal S_{\mathfrak g}$ is connected and of arithmetic genus zero, thus the constant term of $p_{\mathcal S_{\mathfrak g}}$ is $1$. Since 
$\mathcal S_{\mathfrak g}\subseteq H^r_d(X,\mathcal L)$, the coefficient of $v$ in $p_{\mathcal S_{\mathfrak g}}$ is zero. Since $\gamma_{2,*}[\mathcal S]=[T]$, the coefficient of $s_i$ is $1$ for each $i=0,\dots,d$. It remains to show that the coefficient of $u$ is $r+1$, which is equivalent to showing that $f(\mathcal S)$ has degree $r+1$ in $G$.

Let $U_j:=\V_{Q_j}\subseteq W$ for $j=0,1,\dots,m+1$. It follows from Proposition~\ref{hV} that the degree of $f(\mathcal S)$ is 
$$
\sum_{i=1}^{m+1}\Big(\dim_k\rho_1(U_{i-1})-\dim_k\rho_1(U_{i})\Big),
$$
which is equal to $\dim_k\rho_1(U_0)-\dim_k\rho_1(U_{m+1})$.
But 
$$
U_0\subseteq\Gamma(Y,L|_Y)\oplus\Gamma(Z,L|_Z(-P))\quad\text{and}\quad
U_{m+1}\subseteq\Gamma(Y,L|_Y(-(d+1)P))\oplus\Gamma(Z,L|_Z(dP)).
$$
Since $L|_Z(-P)$ and $L|_Y(-(d+1)P)$ have negative degrees, $\dim_k\rho_1(U_0)=r+1$ and $\rho_1(U_{m+1})=0$. 
Thus, the degree of $f(\mathcal S)$ is $r+1$. 
\end{proof}

\subsection{First main theorem}

\begin{thm}\label{chainmaps} For each connected $S\in \text{\rm Hilb}^{\phi,\cc}_{H^r_d(X)}(k)$, the induced map $S\to T$ is a chain map.
\end{thm}

The proof of this theorem is postponed to the next section. We use it now to prove the first main theorem of the article.

\begin{thm}\label{funct-equiv}
The functor $\mathfrak{G}_d^r(X)$ is isomorphic to the functor of points of an open and closed subscheme of $\text{\rm Hilb}_{H_d^r(X)}^{\phi,\cc}$. In particular, there exists a projective scheme $G_d^r(X)$ representing $\mathfrak{G}_d^r(X)$.
\end{thm}

\begin{proof} Proposition~\ref{betaclass} associates to a family of continuous linear series $\mathfrak g=(\L,\V)$ over a scheme $B$ a certain family of $\cc$-invariant subschemes $\mathcal S_{\mathfrak g}\subseteq H^r_d(X)\times B$ over $B$ with relative multivariate Hilbert polynomial $\phi$ 
whose fibers satisfy $h^0(\mathcal S_{\mathfrak g,b},\mathcal O_{\mathcal S_{\mathfrak g,b}})=1$ for every $b\in B$, thus a 
$B$-point of the open subscheme $U\subseteq\text{\rm Hilb}_{H_d^r(X)}^{\phi,\cc}$ parameterizing subschemes $S$ with $h^0(S,\mathcal O_S)=1$. If $\mathfrak g=(\L',\V')$ is another family equivalent to $\mathfrak G$, the corresponding subscheme $\mathcal S_{\mathfrak g'}$ is equal to $\mathcal S_{\mathfrak g}$, as it follows from Definition~\ref{tglsequiv} and the moduli properties of Jacobians and Grassmannians. 

Conversely, let $\mathcal S\subseteq H_d^r(X)\times B$ be a family of 
$\cc$-invariant closed subschemes over $B$ with relative multivariate Hilbert polynomial $\phi$. The natural map $\pi\colon\mathcal S\to B$ is flat. The fiber $\mathcal S_b$ over each $b\in B$ satisfies $h^0(\mathcal S_b,\mathcal O_{\mathcal S_b})=1$ only if $\mathcal S_b$ is connected. Theorem~\ref{chainmaps} yields the converse: if $\mathcal S_b$ is connected then $h^0(\mathcal S_b,\mathcal O_{\mathcal S_b})=1$. It follows that 
$U$ is also the subset parameterizing connected subschemes of $H_d^r(X)$, which is closed by the principle of connectedness. Furthermore, Theorem~\ref{chainmaps} yields, if the fibers of 
$\mathcal S$ are connected, that $\pi$ is a family of chains of rational curves and the natural map 
$\gamma\:\mathcal S\to T\times B$ is a family of chain maps. Since 
$J$ is isomorphic to an Abelian variety, it follows that the map $\mathcal S\to J$ factors through a map $B\to J$ yielding an invertible sheaf $\mathcal L$ on $X\times B$ of bidegree $(d,0)$ over $B$. Finally, the map $\mathcal S\to H^r_d(X)$ yields a locally free subsheaf 
$\V\subseteq p_{\mathcal S*}(\mathcal L(\gamma))$ of rank $r+1$ such that $\mathfrak g_{\mathcal S}:=(\L,\V)$ is a family of twisted generalized linear series. That these series are everywhere nonconstant follows from the fact that $\mathcal S\subseteq H_d^r(X)\times B$. That in addition they are everywhere invariant follows from the fact that $\mathcal S$ is $\cc$-invariant. Thus $\mathfrak g_{\mathcal S}$ is a family of continuous linear series over $B$. 

It is clear from the constructions that $\mathcal S_{\mathfrak g_{\mathcal S}}=\mathcal S$ and that $\mathfrak g_{\mathcal S_{\mathfrak g}}$ is equivalent to $\mathfrak g$. The last statement follows from the fact that $\text{\rm Hilb}_{H_d^r(X)}^{\phi,\cc}$ is projective because $\text{\rm Hilb}_{H_d^r(X)}^{\phi}$ is.  
\end{proof}

\section{Rigidity}

In this section we prove Theorem~\ref{chainmaps}. 

\begin{proof}[Proof of Theorem~\ref{chainmaps}] Since the coefficient of $v$ in $\phi$ is zero, we must have that 
$S_\red\subseteq H^r_d(X,L)$ for some invertible sheaf $L$ on $X$ with bidegree $(d,0)$, where ``red" denotes the reduced induced subscheme. Again, $H^r_d(X,L)$ is an invariant subscheme of $G\times T$, where $G:=\text{Grass}(r+1,W)$ for
$$
W:=W_1\oplus W_2:=\Gamma(Y,L_Y)\oplus\Gamma(Z,L_Z(dP)).
$$
Let $p_1$ and $p_2$ denote the projections of $G\times T$ onto the indicated factors.

Since $S$ is connected, so is $S_{\text{red}}$. Thus, since $S$ has linear multivariate polynomial, $S$ is of pure dimension~1, and hence so is $S_\red$. We will prove a series of claims:

\vskip0.2cm

{\bf First Claim:} {\it $S_{\text{\rm red}}$ has exactly one irreducible component $S_i$ satisfying $p_2(S_i)=T_i$ for each $i=0,\dots,d$, and $p_2$ restricts to an isomorphism $S_i\xrightarrow{\cong}T_i$. The remaining components of $S_{\text{\rm red}}$ are collapsed to the $N_\ell$ under $p_2$.}

\begin{proof}[Proof first claim] 
%For the sake of preciseness,we put $S_{-1}:=\emptyset$ and $S_{d+1}:=\emptyset$. 
Since the coefficient of each $s_i$ in $\phi$ is $1$, 
we have $p_{2*}[S_\red]=[T]$. Thus, 
for each $i=0,\dots,d$ there exists a unique irreducible 
component $S_i$ of $S_{\text{red}}$
mapping birationally onto $T_i$ via $p_2$, the
remaining components being collapsed. Since the $T_i$ are rational, so must be the $S_i$, and the induced maps $S_i\to T_i$ are isomorphisms. 
Since $S$ is invariant, and the action of $\cc$ moves the points on
$T_i^*$ transitively for each $i$, it follows that each remaining
irreducible component of $S_\red$ collapses to a point among
$N_0,\dots,N_{d+1}$. 
\end{proof}

For each $i=0,\dots, d$, let $N_{i,0}$ be the unique point on $S_i$ mapped to $N_i$ and $N_{i,\infty}$ that mapped to $N_{i+1}$. For each $i=0,\dots,d+1$, let $S_{i-\frac{1}{2}}:=(S\cap p_2^{-1}(N_i))_\red$. 

\vskip0.2cm

{\bf Second Claim:} {\it The following statements hold:
\begin{enumerate}
\item The $S_{i-\frac{1}{2}}$ are pairwise disjoint.
\item For each $i=0,\dots,d+1$, we have that $S_{i-\frac{1}{2}}$ intersects $S_j$ if and only if $j\in\{i-1,i\}$, in which case $S_{i-\frac{1}{2}}\cap S_j$ is scheme-theoretically a point.
\item For each $i=1,\dots,d$ we have that $S_{i-\frac{1}{2}}$ is a point only if $N_{i-1,\infty}= N_{i,0}$, in which case that is the point on $S_{i-\frac{1}{2}}$.
\end{enumerate}
}

\begin{proof}[Proof second claim] The first statement is clear, as the images $p_2(S_{i-\frac{1}{2}})$ are pairwise disjoint. The third is clear as well, since, by definition, $N_{i-1,\infty}$ and $N_{i,0}$ are on $S_{i-\frac{1}{2}}$, if defined. As for the second, let  $i\in\{0,\dots,d+1\}$. Since $N_i\in T_j$ if and only if $j\in\{i-1,i\}$, we have that $S_{i-\frac{1}{2}}$ intersects $S_j$ if and only if $j\in\{i-1,i\}$. Furthermore, since each $S_j$ is mapped by $p_2$ isomorphically to $T_j$, and the scheme theoretic image of $S_{i-\frac{1}{2}}$ is $N_i$, we must have that $S_{i-\frac{1}{2}}\cap S_j$ is scheme-theoretically a point for each $j\in\{i-1,i\}$.
\end{proof}

{\bf Third Claim:} {\it $S_{i-\frac{1}{2}}$ is connected for each $i=0,\dots,d+1$.}

\begin{proof}[Proof third claim] Let $i\in\{0,\dots,d+1\}$. We will prove the claim through a series of statements.

\vskip0.1cm

{\bf Statement 1:} {\it For each $i=1,\dots,d$, every connected subcurve $C$ of $\overline{S_\red-\{S_{i-1}\cup S_i\}}$ intersecting both $S_{i-1}$ and $S_i$ is contained in $S_{i-\frac{1}{2}}$.}

Indeed, the image $p_2(C)$ is connected and contains points on $T_{i-1}$ and on $T_i$. It does not contain either $T_{i-1}$ or $T_i$ because only $S_{i-1}$ and $S_i$ cover these components, by First Claim. By the same reason, it does not contain points on $T_0\cup\cdots\cup T_{i-2}$ or $T_{i+1}\cup\cdots\cup T_d$. In particular, $p_2(C)$ contains neither $N_{i-1}$ nor $N_{i+1}$, and hence must be a point on both $T_{i-1}$ and $T_i$, thus equal to $N_i$. So $C\subseteq S_{i-\frac{1}{2}}$. 

\vskip0.1cm

{\bf Statement 2:} {\it $S_{i-\frac{1}{2}}$ contains at most one isolated point, which, if it exists, is equal to $N_{0,0}$ if $i=0$, to $N_{d,\infty}$ if $i=d+1$ and to $N_{i-1,\infty}= N_{i,0}$ if $1\leq i\leq d$.}

Indeed, let $Q$ be an isolated point on $S_{i-\frac{1}{2}}$. Then $Q$ belongs to an irreducible component $C$ of $S_\red$ whose image under $p_2$ contains $N_i$. Also, since $Q$ is isolated in $S_{i-\frac{1}{2}}$, we have that $C\not\subseteq S_{i-\frac{1}{2}}$, that is, $p_2(C)\neq N_i$. But then $C=S_{i-1}$ or $C=S_i$, whence $Q=N_{i-1,\infty}$ or $Q=N_{i,0}$, respectively. The former is possible only if $i>0$ and the latter is possible only if $i\leq d$. Hence, if $i=0$ then $C=S_0$ and $Q=N_{0,0}$, whereas if $i=d+1$ then $C=S_d$ and $Q=N_{d,\infty}$.

Assume now that $1\leq i\leq d$. Either $Q=N_{i-1,\infty}$ or $Q=N_{i,0}$. By contradiction, assume $N_{i-1,\infty}\neq N_{i,0}$. We will show directly that neither $N_{i-1,\infty}$ nor $N_{i,0}$ is isolated on $S_{i-\frac{1}{2}}$. Indeed, since $N_{i-1,\infty}\neq N_{i,0}$, we have $S_{i-1}\cap S_i=\emptyset$. Since $S$ is connected, there is a connected subcurve $C$ of $\overline{S_\red-\{S_{i-1}\cup S_i\}}$ intersecting both $S_{i-1}$ and $S_i$. Statement~1 yields $C\subseteq S_{i-\frac{1}{2}}$. Since $C$ intersects $S_i$, there is an irreducible component $C'$ of $C$ intersecting $S_i$. Since $C'\subseteq S_{i-\frac{1}{2}}$, the intersection is $N_{i,0}$, and hence $N_{i,0}$ is not isolated on $S_{i-\frac{1}{2}}$. Analogously, neither is $N_{i-1,\infty}$. 

\vskip0.1cm

{\bf Statement 3:} {\it Let $C$ be a maximal connected subcurve of $S_\red$ such that $C\subseteq S_{i-\frac{1}{2}}$. Then $C$ contains 
$N_{i-1,\infty}$ if $i\geq 1$ and $N_{i,0}$ if $i\leq d$.}

Indeed, reason by contradiction. Without loss of generality, assume $i\leq d$ and $C$ does not contain $N_{i,0}$. Since $p_2(C)=N_i$, this means that $C$ does not intersect $S_i$. Since $S_\red$ is connected, there is a connected subcurve $C'$ of $\overline{S_\red-\{S_i\cup C\}}$ intersecting both $S_i$ and $C$. Let $C''$ be an irreducible component of $C'$ intersecting $C$. Then $N_i\in p_2(C'')$. Since $C'\neq S_i$, either $p_2(C' )=N_i$ or $C''=S_{i-1}$, if $i\geq 1$. The former contradicts the maximality of $C$. Thus $i\geq 1$ and $C''=S_{i-1}$. Since $p_2(C)=N_i$, we must have that $S_{i-1}$ meets $C$ at $N_{i-1,\infty}$. Now, if $S_{i-1}$  meets $S_i$, it must be at a point over $N_i$, but the unique point on $S_{i-1}$ above $N_i$ is $N_{i-1,\infty}$. Since $N_{i-1,\infty}$ lies on $C$ and $C$ does not meet $S_i$, we must have that $S_{i-1}$ does not meet $S_i$. But since $C'$ meets $S_i$ and contains $S_{i-1}$, there is a connected subcurve $C'''\subseteq\overline{C'-S_{i-1}}$ meeting $S_{i-1}$ and $S_i$. By Statement~1, we must have $p_2(C''')=N_i$. Then $C'''$ does not intersect $C$ by the maximality of $C$. In particular, $C'''$ does not contain $N_{i-1,\infty}$. But $C'''$ intersects $S_{i-1}$, and can only intersect $S_{i-1}$ at $N_{i-1,\infty}$, a contradiction.

\vskip0.1cm

{\bf Statement 4:} {\it If $S_{i-\frac{1}{2}}$ is not a point, then 
$S_{i-\frac{1}{2}}$ has no isolated point.} 

Indeed, since $S_{i-\frac{1}{2}}$ has at most one isolated point by Statement~2, if $S_{i-\frac{1}{2}}$ is not a point then it contains a nonisolated point $Q$. Let $C$ be a maximal connected subcurve of $S_\red$ containing $Q$ such that $C\subseteq S_{i-\frac{1}{2}}$. Then Statement~3 says that $C$ contains $N_{i-1,\infty}$ if $i\geq 1$ and $N_{i,0}$ if $i\leq d$. In particular, neither $N_{i-1,\infty}$ for $i\geq 1$ nor $N_{i,0}$ for $i\leq d$ is isolated in $S_{i-\frac{1}{2}}$. But then $S_{i-\frac{1}{2}}$ contains no isolated point by Statement~2.

\vskip0.1cm

{\bf Statement 5:} {\it $S_{i-\frac{1}{2}}$ is connected.}

If $S_{i-\frac{1}{2}}$ is a point, then it is clearly connected. If not, then $S_{i-\frac{1}{2}}$ has no isolated point by Statement~4. It is thus a disjoint union of connected subcurves of $S_\red$. By Statement~3, each such subcurve contains $N_{i-1,\infty}$ if $i\geq 1$ and $N_{i,0}$ if $i\leq d$. But then all these subcurves contain the same point, and hence there is just one of them. In other words, $S_{i-\frac{1}{2}}$ is connected.
\end{proof}

{\bf Fourth Claim:} {\it There are irreducible components $C_0,\dots,C_m$ of $S_\red$ such that
\begin{enumerate} 
\item $C_0=S_0$ and $C_m=S_d$;
\item $C_{i-1}\cap C_i\neq\emptyset$ for $i=1,\dots,m$. 
\item There is an increasing sequence $i_0,i_1,\dots,i_d$ of integers such that $C_{i_\ell}=S_\ell$ for $\ell=0,\dots,d$;
\item For each $\ell=1,\dots,d$, we have $i_{\ell}=i_{\ell-1}+1$ if $S_{\ell-1}$ meets $S_\ell$; otherwise, $i_{\ell}>i_{\ell-1}+1$ and $C_{i_{\ell-1}+1}\cup\cdots\cup C_{i_{\ell}-1}$ is contained in $S_{\ell-\frac{1}{2}}$.
\end{enumerate}}

\begin{proof}[Proof fourth claim]
Start with $C_0:=S_0$. Each time $C_i=S_{\ell-1}$ for $\ell=1,\dots,d$, we do the following: If $S_{\ell}$ intersects $S_{\ell-1}$, we put $C_{i+1}:=S_\ell$. If not, then $N_{\ell-1,\infty}\neq N_{\ell,0}$, and thus $S_{\ell-\frac{1}{2}}$ is not a point by Second Claim. By Third Claim, $S_{\ell-\frac{1}{2}}$ is connected. By Second Claim, there are irreducible components $C$ and $C'$ of $S_{\ell-\frac{1}{2}}$ such that $C$ intersects $S_{\ell-1}$ and $C'$ intersects $S_\ell$. Put $C_{i+1}:=C$. Since $S_{\ell-\frac{1}{2}}$ is connected, there is a sequence of irreducible components of $S_{\ell-\frac{1}{2}}$, each intersecting the next, starting with $C_{i+1}$ and ending with $C'$. Denote these components sequentially by $C_{i+1},\cdots,C_{i+j}$, where $j$ is thus the number of components in the sequence and $C_{i+j}=C'$. Then put $C_{i+j+1}:=S_{\ell}$.
\end{proof} 

We observe that $S_\red$ may not be the union of the $C_i$. Indeed, 
$S_{\ell-\frac{1}{2}}$ may not be the union $C_{i_{\ell-1}+1}\cup\cdots\cup C_{i_{\ell}-1}$.

\vskip0.2cm

{\bf Fifth Claim:} {\it The $p_1(C_i)$ which are orbit closures are either points or smooth rational curves. The sum of their degrees in $G$ is at least $r+1$, with equality only if:
\begin{enumerate}
    \item all the $p_1(C_i)$ are orbit closures, 
    \item for $i,j\in\{0,\dots,m\}$ with $i<j$ such that $p_1(C_i)$ and $p_1(C_j)$ are curves that intersect, $p_1(C_i)\cap p_1(C_j)$ is scheme-theoretically a point, and $p_1(C_{i+1}),\dots,p_1(C_{j-1})$ are all equal to this point.
\end{enumerate}}

\begin{proof}[Proof fifth claim] For each $\ell=1,2$, let $\iota_\ell\:W_\ell\to W$ be the natural injection and $\rho_\ell\:W\to W_\ell$ the natural surjection. Since $S$ is $\cc$-invariant, each $C_i$ is $\cc$-invariant, whence the $p_1(C_i)$ are $\cc$-invariant. 

For each $i=0,\dots,m$, let $V_i\subseteq W$ be the image under $p_1$ of a general point on $C_i$. Put 
\[
a_i:=\min\{\dim_k\rho_1(V)\,|\,V\in p_1(C_i)\}
\quad\text{and}\quad 
A_i:=\max\{\dim_k\rho_1(V)\,|\,V\in p_1(C_i)\}.
\]
Thus $a_i\leq A_i$. If $V_i$ is fixed by the action of $\cc$, then $V_i=\rho_1(V_i)\oplus\rho_2(V_i)$. The same holds for each $V\in p_1(C_i)$ by continuity, whence $a_i=A_i$. In this case, if $p_1(C_i)$ is an orbit closure then $p_1(C_i)$ is the point $V_i$. If $V_i$ is not fixed, then $p_1(C_i)$ is the closure of the orbit of $V_i$ in $G$, and is thus a smooth rational curve of degree $A_i-a_i>0$ by Proposition~\ref{hV}. Furthermore, the boundary of the orbit of $V_i$ consists of two points on $p_1(C_i)$ where $a_i$ and $A_i$ are attained. 
%Now, whether $V_j$ is fixed or not, $p_1(C_i\cap C_j)$ is mapped to the boundary of the orbit of $V_i$ for $j\neq i$. 

Now, for each $i=1,\dots,m$, since $C_{i-1}\cap C_i\neq\emptyset$, it follows that  
the intervals $[a_{i-1},A_{i-1}]$ and $[a_i,A_i]$ intersect, and thus 
\[
I:=\bigcup_{i=0}^m[a_i,A_i]
\]
is an interval contained in $[0,r+1]$.
Now, $C_0=S_0$ and $C_m=S_d$. Since $p_2(N_{0,0})=N_0$ and $p_2(N_{d,\infty})=N_\infty$, we have that
\[
p_1(N_{0,0})\subseteq\Gamma(Y,L|_Y)\oplus\Gamma(Z,L|_Z(-P))\quad
\text{and}\quad 
p_1(N_{d,\infty})\subseteq\Gamma(Y,L|_Y(-(d+1)P))\oplus\Gamma(Z,L|_Z(dP)).
\]
Since $L|_Z(-P)$ and $L|_Y(-(d+1)P)$ have negative degrees, it follows that $A_0=r+1$ and $a_m=0$. 
Thus $I=[0,r+1]$. Then the sum of the degrees of the $p_1(C_i)$ for $i$ such that $V_i$ is nonfixed, which is $\sum_i(A_i-a_i)$, is at least $r+1$.

Assume the sum is exactly $r+1$. Then, first, if $V_i$ is fixed, $p_1(C_i)$ must be a point, $V_i$. Thus, all the $p_1(C_i)$ are orbit closures. Second, the intervals $[a_i,A_i]$ can overlap only at the boundaries. For each $i=1,\dots,m$, since $[a_{i-1},A_{i-1}]$ and $[a_i,A_i]$ intersect, we must have that either 
$a_i=A_{i-1}$ or $A_i=a_{i-1}$. Since $A_0=r+1$ and $a_m=0$, we must have that $A_i=a_{i-1}$ for each $i=1,\dots,m$. 

Let now $i,j\in\{0,\dots,m\}$ with $i<j$ such that $p_1(C_i)$ and $p_1(C_j)$ are curves that intersect. Since $i<j$, we have $A_j\leq a_i$. Since $p_1(C_i)$ and $p_1(C_j)$ intersect, the intervals $[a_i,A_i]$ and $[a_j,A_j]$ intersect as well, and hence $A_j=a_{j-1}=A_{j-1}=\cdots=A_{i+1}=a_i$. It follows that 
$p_1(C_{i+1}),\dots,p_1(C_{j-1})$ are points, actually, the same point $V\in G$, since $C_{i+1}\cup\cdots\cup C_{j-1}$ is connected.  Furthermore, since $A_j=a_i$, and $p_1(C_i)$ and $p_1(C_j)$ are orbit closures that intersect, it follows from Proposition~\ref{hV} that 
there is a unique point of intersection, which is
\[
\iota_1^{-1}(V_i)\oplus\rho_2(V_i)=\rho_1(V_j)\oplus\iota_2^{-1}(V_j).
\]
In addition, Proposition~\ref{hVV} yields that the intersection is transverse. Since 
$C_i$ intersects $C_{i+1}$ and $C_j$ intersects $C_{j-1}$, we must have that $V$ is the intersection.
\end{proof}

Put $S^0:=C_0\cup\dots\cup C_m$. 

\vskip0.2cm

{\bf Sixth Claim:} {\it The sum of the degrees of the $p_1(C_i)$ in $G$ is $r+1$. Each $C_i$ is rational and the map $C_i\to p_1(C_i)$ is either an isomorphism or constant, the latter only if $i=i_\ell$ for some $\ell$. Furthermore, $S^0=S_\red$ and $S$ is generically reduced. In particular, $S_{-\frac{1}{2}}$ and $S_{d+\frac{1}{2}}$ are points.}

\begin{proof}[Proof sixth claim] Since $S$ has Hilbert polynomial $\phi$, and since the coefficient of $u$ in $\phi$ is $r+1$, the products of the degree of the covering $C_i\to p_1(C_i)$ with the degree of $p_1(C_i)$ sum up to at most $r+1$. On the other hand, the degrees of the $p_1(C_i)$ sum up to at least $r+1$ by Fifth Claim. It thus follows that both sums are $r+1$ and that $C_i\to p_1(C_i)$ is either constant or an isomorphism for each $i=0,\dots,m$. However, since $C_i\subseteq G\times T$, the curve $p_1(C_i)$ can only be a point if $p_2(C_i)$ is not one, that is, if $i=i_\ell$ for some $\ell$, in which case $C_i$ is isomorphic to $T_\ell$ under $p_2$. In any case, $C_i$ is a rational curve. 

Finally, the multivariate Hilbert polynomial of $S^0$ has the same linear part as that of $S$. Since $S^0\subseteq S_\red\subseteq S$, it follows that $S^0=S_\red$ and $S$ is generically reduced. 
\end{proof}

\vskip0.2cm

{\bf Seventh Claim:} {\it For each $i,j\in\{0,\dots,m\}$ with $i<j$, the curve $C_i$ intersects $C_j$ only if $j=i+1$, transversally at a single point if so.}

\begin{proof}[Proof seventh claim] Assume $C_i\cap C_j\neq\emptyset$. 
Since the $S_{\ell-\frac{1}{2}}$ are disjoint, and $S_h\cap S_q\neq\emptyset$ only if $|h-q|\leq 1$, we must have 
$i_{\ell-1}\leq i<j\leq i_{\ell}$ for some $\ell$. Now, since $p_2(S_{\ell-\frac{1}{2}})=N_\ell$ scheme-theoretically, and 
$S_{\ell-1}$ is mapped isomorphically to 
$T_{\ell-1}$, 
%for each $p\in\{i_{\ell-1}+1,\dots,i_{\ell}\}$, 
we have that $C_j$ intersects $C_{i_{\ell-1}}$ only if $j=i_{\ell-1}+1$, transversally at $N_{\ell-1,\infty}$ if so. Thus 
the claim holds if $i=i_{\ell-1}$. By analogy, it holds as well if $j=i_\ell$.

Assume $i_{\ell-1}<i<j<i_\ell$. Then the maps $C_q\to p_1(C_q)$ are isomorphisms for all $q$ with $i\leq q\leq j$ by Sixth Claim.
%$p_1(C_i),p_1(C_{i+1}),\dots,p_1(C_j)$ are subcurves of $G$ by Sixth Claim. 
Since $C_i\cap C_j\neq\emptyset$, it follows from Fifth Claim that $j=i+1$ and $p_1(C_i)\cap p_1(C_j)$ is scheme-theoretically a point. Furthermore, since the maps $C_i\to p_1(C_i)$ and $C_j\to p_1(C_j)$ are isomorphisms by Sixth Claim, also $C_i\cap C_j$ is scheme-theoretically a point.
\end{proof}

{\bf Eighth Claim:} {\it The curve $S$ is a chain of rational curves and $p_1\:S\to T$ is a chain map.}

\begin{proof} It follows from Sixth Claim that the $C_i$ are rational, and thus $S^0$ is a chain of rational curves by Seventh Claim. Then $S^0$ is connected of arithmetic genus $0$, and hence the constant part of its multivariate Hilbert polynomial coincides with that of $S$. So does its linear part by Sixth Claim. Hence $S$ and $S^0$ have the same Hilbert polynomial, and since $S^0\subseteq S$, we must have $S^0=S$. So $S$ is a chain of rational curves. First and Fourth Claims now yield that $p_1\colon S\to T$ is a chain map.
\end{proof}

The proof of Theorem~\ref{chainmaps} is now complete.
\end{proof}


\section{Level-$\delta$ limit linear series}\label{deltalls}

Let $\delta=(\delta_1,\dots,\delta_d)$ be a $d$-tuple of natural numbers. Let 
$$
\Delta:=\Delta(\delta):=\Big\{0,\frac{1}{\delta_1},\frac{2}{\delta_1},\dots,1,
1+\frac{1}{\delta_2},1+\frac{2}{\delta_2},\dots,2,\dots,d-1,
d-1+\frac{1}{\delta_d},d-1+\frac{2}{\delta_d},\dots,d\Big\}.
$$
We say that $i,j\in\Delta$ are \emph{consecutive} if $i<j$ and $\Delta\cap(i,j)=\emptyset$.

\subsection{Review}

The reader will recognize the notions and definitions below from \cite{ENR}, where we restricted to the case where all the $\delta_i$ are equal. However, the general theory is similar.

Let $L$ be an invertible sheaf on $X$. The \emph{twist $\delta$-sequence} associated to $L$ is the collection of sheaves
$L^{(i)}$ indexed by
$i\in\Delta$, where
$L^{(i)}:=L\otimes\mathcal O_X^{(i)}$ if $i\in\mathbf Z$ and, otherwise,
\begin{equation}\label{Linot}
L^{(i)}:=\Big(L|_Y\ox\O_Y\big(\lfloor-i\rfloor P\big)\Big)\bigoplus\Big(L|_Z\ox\O_Z\big(\lfloor i\rfloor P\big)\Big),
\end{equation}
where $\lfloor i\rfloor$ stands for the floor of $i$, the largest integer at most $i$.

The condition on the $L^{(i)}$ ensures that there are isomorphisms
$$
L^{(i)}|_Z\longrightarrow\text{Ker}\big(L^{(j)}\to L^{(j)}|_Y\big)\quad\text{and}\quad
L^{(j)}|_Y\longrightarrow\text{Ker}\big(L^{(i)}\to L^{(i)}|_Z\big)
$$
for each consecutive $i,j\in\Delta$, whence there are nonzero maps $\varphi^i\colon L^{(i)}\to L^{(j)}$ factoring through $L^{(i)}|_Z$ and $\varphi_i\colon L^{(j)}\to L^{(i)}$ factoring through $L^{(j)}|_Y$. The 
%<sheaves $L^{(i)}$ are defined up to isomorphism and the 
maps $\varphi^i$ and $\varphi_i$ are defined up to multiplication by $k^*$. 

If $L=\mathcal O_X$, then our notation is compatible with that used in Subsection~\ref{sec:twisters}. 

\begin{defn} 
A {\it level-$\delta$ limit linear series} of degree $d$ on $X$ is
the data $\mathfrak g=(L;V^{(i)},\,i\in\Delta)$ of an
invertible sheaf $L$ on $X$ of bidegree $(d,0)$ and vector subspaces
$V^{(i)}\subseteq\Gamma(X,L^{(i)})$ for $i\in\Delta$
of equal dimension such that
\begin{equation}\label{compatible}
\varphi^i(V^{(i)})\subseteq
V^{(j)}\text{\ and\ }
\varphi_i(V^{(j)})\subseteq
V^{(i)}\text{\ for each consecutive }i,j\in\Delta.
\end{equation}
We say that $\mathfrak g$ has rank $r$
if the $V^{(i)}$ have projective dimension $r$. We say that
$\mathfrak g$ is \emph{exact} if 
\begin{equation}\label{exact}
\text{Im}\big(\varphi^i|_{V^{(i)}}\big)=
\text{Ker}\big(\varphi_i|_{V^{(j)}}\big)
\text{\ and\ }
\text{Im}\big(\varphi_i|_{V^{(j)}}\big)=
\text{Ker}\big(\varphi^i|_{V^{(i)}}\big)\text{\ for consecutive }i,j\in\Delta.
\end{equation}
\end{defn}

%Though it may seem that the data giving $\mathfrak g$ depend on infinitely many parameters, this is not true. In fact, up to shifting, we may and will assume $L$ has multidegree $(d,0)$ for a nonnegative integer $d$, whence the $L^{(i)}$ for $i=0,\dots,d$ are the only sheaves with effective quotient invertible sheaves both on $Y$ and $Z$. It follows that $\varphi^i(V^{(i)})=V^{(i+1)}$ for each $i\geq d$ and $\varphi_i(V^{(i+1)})=V^{(i)}$ for each $i<0$. The data $(L,V^{(0)},\dots,V^{(d)})$ is thus equivalent to $\mathfrak g$. For $\delta=1$, it is precisely a limit linear series in the sense given by Osserman \cite[Def. 2.1, p.~1.168]{Oss1}.

%Observe that $\mathcal O_X^{(i)}\subseteq\mathcal O_Y(-iP)\oplus\mathcal O_Z(iP)$ for each $i=0,1\dots,d$, whence $L^{(i)}\subseteq L|_Y\oplus L|_Z(dP)$ for every $i\in\Delta$. But i
Let $\mathfrak g=(L;V^{(i)},\,i\in\Delta)$ be a level-$\delta$ limit linear series of rank $r$ and degree $d$ on $X$. It follows from \eqref{Linot} that 
$L^{(i)}\subseteq L|_Y\oplus L|_Z(dP)$ for each $i\in\Delta$. Thus, subspaces $V^{(i)}$ of $\Gamma(X, L^{(i)})$ may be viewed as subspaces of 
$U:=U_1\oplus U_2$, where $U_1:=\Gamma(Y,L|_Y)$ and $U_2:=\Gamma(Z,L|_Z(dP))$. As such, Conditions~\eqref{compatible} are equivalent to
\begin{equation}\label{compatible2}
\rho_2(V^{(i)})\subseteq\iota_2^{-1}(V^{(j)})
\quad\text{and}\quad
\rho_1(V^{(j)})\subseteq\iota_1^{-1}(V^{(i)})
\quad\text{for each consecutive }i,j\in\Delta,
\end{equation}
whereas Conditions~\eqref{exact} are equivalent to
\begin{equation}\label{exact2}
\rho_2(V^{(i)})=\iota_2^{-1}(V^{(j)})
\quad\text{and}\quad
\rho_1(V^{(j)})=\iota_1^{-1}(V^{(i)})
\quad\text{for each consecutive }i,j\in\Delta,
\end{equation}
where 
$\iota_j\: U_j\to U_1\oplus U_2$ is the natural injection and 
$\rho_j\: U_1\oplus U_2\to U_j$ is the natural surjection for $j=1,2$.

The data of the vector spaces $(V^{(i)},\,i\in\Delta)$ and the maps
$h_i:=\varphi_{i}|_{V^{(j)}}$ and 
$h^i:=\varphi^i|_{V^{(i)}}$ for consecutive $i,j\in\Delta$ 
will be called a \emph{linked chain of vector spaces}. For each $i\in\Delta$, 
let $m_i:=(r+1)-p_i-q_i$ where
$$
p_i:=\dim\Ker(h_{i-1}),\quad q_i:=\dim\Ker(h^i).
$$
Let $N_{\mathfrak g}\colon\Delta\to\mathbf Z^3$ be the map taking $i$ to $(p_i,q_i,m_i)$; it is called the \emph{numerical data} associated to the linked chain and thus to $\mathfrak g$. We say $N_{\mathfrak g}$ is \emph{exact} if $\sum m_i=r+1$. Then $\mathfrak g$ is exact if and only 
if $N_{\mathfrak g}$ is; see \cite[Prop. 3.6, p.~836]{ENR}. We say $N_{\mathfrak g}$ and $\mathfrak g$ are \emph{minimal} if $m_i>0$ for each $i\in\Delta-\mathbf Z$.

As in \cite{ENR}, Prop.~3.2, p.~834, there are a projective
scheme $G_{d,\delta}^r(X)$ parameterizing level-$\delta$
limit linear series $\mathfrak g=(L;V^{(i)},\,i\in\Delta)$ of degree $d$ and rank $r$ on $X$, and an open
subscheme
$$
G^{r,*}_{d,\delta}(X)\subseteq
G^{r}_{d,\delta}(X),
$$
parameterizing those which are exact. Also, there is a natural action of 
$\cc^\Delta$ on $G_{d,\delta}^r(X)$ such that 
$(c_i,i\in\Delta)\ast(L_1;V_1^{(i)},i\in\Delta)=(L_2; V_2^{(i)},i\in\Delta)$ if and only if $L_1\cong L_2$ and $V_2^{(i)}=V_1^{(i)}$ for $i\in\Delta\cap\mathbf Z$ and $V_2^{(i)}=(c_i,1)V_1^{(i)}$ otherwise, as subspaces of $U=U_1\oplus U_2$. The action leaves the sets of minimal and exact limit linear series invariant. Two level-$\delta$ limit linear series 
$(L_1;V_1^{(i)},i\in\Delta)$ and $(L_2;V_2^{(i)},i\in\Delta)$ are called \emph{equivalent} if they lie on the same orbit by the action of $\cc^{\Delta}$. Equivalent limit linear series have the same numerical data.

As shown in \cite{ENR}, Section~3, the numerical data yield a stratification of $G^r_{d,\delta}(X)$ into locally closed subsets, as follows: For each 
function $N\:\Delta\to\z^3$, let
\[
G_{d,\delta}^r(X;N):=\left\{\mathfrak{g}\in 
G_{d,\delta}^r(X)\,|\, N_{\mathfrak g}=N\right\}.
\]
So that $G_{d,\delta}^r(X;N)$ is nonempty, a certain condition on $N$ is necessary, called admissibility in \cite{ENR}, Def.~3.7, p.~836. In addition, the open
subscheme $G^{r,*}_{d,\delta}(X)\subseteq G^{r}_{d,\delta}(X)$ decomposes as a disjoint union of subsets which are both open and closed,
\begin{equation}\label{estrata}
G_{d,\delta}^{r,*}(X)=\coprod_{N\text{ exact}}
G_{d,\delta}^r(X;N).
\end{equation}
The $G_{d,\delta}^r(X;N)$ can be given a natural structure of moduli scheme.

For each (exact) function $N\:\Delta\to\mathbf Z^3$, let $\Delta_N\subseteq\Delta$ be the subset of those $i$ for which $i\in\mathbf Z$ or $N(i)\not\in\mathbf Z^2\times\{0\}$. Let $\delta'=(\delta'_1,\dots,\delta'_d)$ be the unique $d$-tuple of natural numbers for which there is an isomorphism of posets $f\:\Delta'\to\Delta_N$ such that $f(i)=i$ for each $i\in\z$, 
where $\Delta':=\Delta(\delta')$. Let $N':=Nf$. Then $N'$ is (exact and) minimal. Analogously to what was shown in \cite{ENR}, Prop.~4.3, p.~839, the natural forgetful morphism $G_{d,\delta}^r(X;N)\to G_{d,\delta'}^r(X;N')$, taking the limit linear series $(L;V(i),i\in\Delta)$ to $(L;V(f(i)),i\in\Delta')$, is an isomorphism.

Because of this isomorphism, we will consider the open subscheme
\[
G_{d,\delta}^{r,*}(X;\text{\rm min})\subseteq G_{d,\delta}^{r,*}(X),
\]
which is the (disjoint) union of the 
$G_{d,\delta}^r(X;N)$ for $N$ exact and minimal. Clearly, we have
$G_{d,\mathbf 1}^{r,*}(X;\text{\rm min})=G_{d,\mathbf 1}^{r,*}(X)$, where $\mathbf 1:=(1,\dots,1)$.



\subsection{Second main theorem}

Let $T^{(\delta)}$ be the chain of rational curves obtained from $T$ by splitting the branches of $T$ at $N_i$ and connecting them by a chain of rational curves of length $\delta_i-1$ for each $i=1,\dots,d$. Thus we index the components $T^{(\delta)}_i$ of $T^{(\delta)}$ with $i\in\Delta$. For each consecutive  $i,j\in\Delta$, let $N^{(\delta)}_j$ be the intersection of $T^{(\delta)}_{i}$ with $T^{(\delta)}_j$. Clearly, there is a natural map $\mu^{(\delta)}\:T^{(\delta)}\to T$ collapsing all the components $T^{(\delta)}_{i}$ for $i\not\in\mathbf Z$, and identifying $T^{(\delta)}_{i}$ with $T_i$ for $i\in\mathbf Z$.

Let $N_0^{(\delta)}$ and $N_\infty^{(\delta)}$ be the points on $T^{(\delta)}$ such that $\mu^{(\delta)}(N_0^{(\delta)})=N_0$ and $\mu^{(\delta)}(N_\infty^{(\delta)})=N_\infty$. Then $(T^{(\delta)},N_0^{(\delta)},N_\infty^{(\delta)})$ is a two-pointed chain and $\mu^{(\delta)}$ is a chain map. 

For each $i\in\mathbf Z$, let $E_i^{(\delta)}$ be the point on $T_i^{(\delta)}$ above $E_i$. For each $i\in\Delta-\mathbf Z$, fix $E_i^{(\delta)}\in(T^{(\delta)}_i)^*$. The chain action of $\cc$ on $(T,N_0,N_\infty)$ such that $1\ast E_i=E_i$ for each $i=0,\dots,d$ lifts to a unique chain action on $(T^{(\delta)},N_0^{(\delta)},N_\infty^{(\delta)})$ such that $1\ast E^{(\delta)}_i=E^{(\delta)}_i$ for each $i\in\Delta$.

\begin{thm}\label{allevels} For each 
$\delta=(\delta_1,\dots,\delta_d)\in\mathbf N^d$ and exact and minimal function $N\colon\Delta(\delta)\to\n^3$
there is a natural morphism
\[
\Psi_\delta^N\:G^{r}_{d,\delta}(X;N)\longrightarrow
G_d^r(X).
\]
The map takes a (minimal exact) level-$\delta$ limit linear series $(L;V^{(i)},i\in\Delta)$ on $X$ with numerical data $N$ to the continuous linear series $(L,\V)$ along $\mu^{(\delta)}$, where $\V|_{E_i^{(\delta)}}$ is taken to $V^{(i)}$ under the choice of an isomorphism $L(\mu^{(\delta)})|_{X\times E_i^{(\delta)}}\xrightarrow{\cong} L^{(i)}$ for each $i\in\Delta$. Each nonempty fiber of $\Psi_\delta^N$ is an equivalence class of level-$\delta$ limit linear series, and vice-versa. Finally, the images of all the $\Psi_\delta^N$ form a partition of the underlying set of $G_d^r(X)$. 
\end{thm}

(If $\delta=\mathbf 1$, then $\mu^{(\delta)}=1_T$, all the $L^{(i)}$ are invertible, and thus the $\V|_{E_i}$ do not depend on the choice of an isomorphism $L(1_T)|_{X\times E_i}\cong L^{(i)}$. Also, $\Psi^N_{\mathbf 1}$ is injective.)

%<{\color{red} Seria $\Psi_{\mathbf 1}$ um isomorfismo sobre a imagem?}

\begin{proof} Put $\Delta=\Delta(\delta)$. Consider a family of level-$\delta$ limit linear series on $X$ of degree $d$, rank $r$ and numerical data $N$ over a scheme $B$. 
%To define $\Psi_\delta$, we may assume all limit linear series in the family have numerical data $N$ for $N\:\Delta\to\mathbf Z^3$ exact and minimal. 
The family is given by an invertible sheaf $\L$ of relative bidegree $(d,0)$ on $X\times B/B$ and a family of generalized linear series $(\L^{(i)},\mathcal V^{(i)})$ on $X\times B/B$, where $\L^{(i)}:=\L(\mu^{(\delta)})|_{X\times E_i^{(\delta)}\times B}$ under the usual identification of $E_i^{(\delta)}\times B$ with $B$, for each $i\in\Delta$. The data $(\L;\V^{(i)},i\in\Delta)$ satisfies certain conditions, to be explained next. 

Let $D$ be an effective divisor of the smooth locus of $X$. Put
\[
\mathcal W_1:=p_{B*}\big(\L|_{Y\times B}((D\cap Y)\times B)\big)
\quad\text{and}\quad
\mathcal W_2:=p_{B*}\big(\L|_{Z\times B}(dP\times B+(D\cap Z)\times B)\big),
\]
and let $\W:=\W_1\oplus\W_2$, where $p_B\:X\times B\to B$ is the projection. Assume $D$ is ample enough that the $\mathcal W_i$ be locally free with formation commuting with base change, and the natural maps $\V^{(i)}\to\mathcal W_1\oplus\mathcal W_2$ be injective with locally free quotients. 

Let $\rho_\ell\:\W_1\oplus\W_2\to\W_\ell$ and $\iota_\ell\:\W_\ell\to\W_1\oplus\W_2$ be the natural maps for $\ell=1,2$. 
The conditions $(\L;\V^{(i)},i\in\Delta)$ satisfies are three. First, the map $\V^{(i)}\to \mathcal W_\ell$ induced by $\rho_\ell$ has locally free cokernel for each $i\in\Delta$ and $\ell=1,2$. It follows that 
$\mathcal A^{(i)}_\ell:=\rho_\ell(\V^{(i)})$ and 
$\mathcal B^{(i)}_\ell:=\iota_{\ell}^{-1}(\V^{(i)})$ 
are locally free. Their ranks are specified by $N$, the second condition.
%<{\color{red} Não me parece que sabendo que a restrição do conúcleo pra cada fibra tem a mesma dimensão implique que o conúcleo é localmente livre. Esta última condição teria que ser colocada na família, que precisa ser definida.}
Notice that
\[
\mathcal B_1^{(i)}\oplus\mathcal B_2^{(i)}\subseteq
\V^{(i)}
\subseteq\mathcal A_1^{(i)}\oplus\mathcal A_2^{(i)}
\subseteq \mathcal W_1\oplus\mathcal W_2,
\]
and the induced map
\[
\lambda^{(i)}_\ell\:\frac{\V^{(i)}}{\mathcal B_1^{(i)}\oplus\mathcal B_2^{(i)}}\xrightarrow{\quad\cong\quad}\frac{\mathcal A_\ell^{(i)}}{\mathcal B_\ell^{(i)}}
\]
is an isomorphism for each $\ell=1,2$. Composing the inverse of $\lambda^{(i)}_1$ with $\lambda^{(i)}_2$ gives us an isomorphism
\[
\zeta^{(i)}\: \frac{\mathcal A_1^{(i)}}{\mathcal B_1^{(i)}}\xrightarrow{\quad\cong\quad}\frac{\mathcal A_2^{(i)}}{\mathcal B_2^{(i)}}.
\]
Finally, $\mathcal A^{(i)}_2=\mathcal B^{(j)}_2$ and $\mathcal A^{(j)}_1=\mathcal B^{(i)}_1$ 
for each consecutive $i,j\in\Delta$, by the exactness of $N$, the third condition.

For each $i\in\Delta$, let $p_i\: T^{(\delta)}_i\times B\to B$ and $q_i\:T^{(\delta)}_i\times B\to T^{(\delta)}_i$ be the projections, and let $\widetilde\V^{(i)}\subseteq p_i^*(\mathcal A_1^{(i)}\oplus\mathcal A_2^{(i)})$ be the subsheaf whose quotient is the cokernel of the composition 
\[
q_i^*\mathcal O_{T^{(\delta)}_i}(-1)\bigotimes\frac{p_i^*\mathcal A_1^{(i)}}{p_i^*\mathcal B_1^{(i)}} \xrightarrow{\quad q_i^*\xi_i\otimes 1\quad}
\frac{p_i^*\mathcal A_1^{(i)}}{p_i^*\mathcal B_1^{(i)}}\bigoplus\frac{p_i^*\mathcal A_1^{(i)}}{p_i^*\mathcal B_1^{(i)}} \xrightarrow{\quad 1\oplus p_i^*\zeta^{(i)}\quad}
\frac{p_i^*\mathcal A_1^{(i)}}{p_i^*\mathcal B_1^{(i)}}\bigoplus\frac{p_i^*\mathcal A_2^{(i)}}{p_i^*\mathcal B_2^{(i)}},
\]
where $\xi_i$ is the inclusion $\mathcal O_{T^{(\delta)}_i}(-1)\hookrightarrow\mathcal O_{T^{(\delta)}_i}^2$ of the tautological subsheaf associated to an isomorphism $T^{(\delta)}_i\xrightarrow{\cong}\mathbf P^1_k$ such that 
\begin{equation}\label{cvary}
\mathcal O_{T^{(\delta)}_i}(-1)\big|_{c\ast E_i^{(\delta)}}=(c^{-1},1)\mathcal O_{T^{(\delta)}_i}(-1)\big|_{E_i^{(\delta)}}\quad\text{for each }c\in\cc(k).
\end{equation}
Then $\widetilde\V^{(i)}\subseteq p_i^*\W$ with locally free quotient, and satisfies
\begin{equation}\label{eq:VtildeV}
\widetilde\V^{(i)}\big|_{c\ast E_i^{(\delta)}}=(c^{-1},1)p_1^*\V^{(i)}\big|_{E_i^{(\delta)}}\quad\text{for each }c\in\cc(k).
\end{equation}

Moreover, there is a subsheaf $\widetilde\V\subseteq p^*\W$ with locally free quotient such that $\widetilde\V|_{T^{(\delta)}_i\times B}=\widetilde\V^{(i)}$ for each $i\in\Delta$, where $p\:T^{(\delta)}\times B\to B$ is the projection, because 
\[
\widetilde\V^{(i)}|_{N^{(\delta)}_j\times B}=\mathcal B_1^{(i)}\oplus\mathcal A_2^{(i)}=\mathcal A_1^{(j)}\oplus\mathcal B_2^{(j)}=\widetilde\V^{(j)}|_{N^{(\delta)}_j\times B}
\]
for each consecutive $i,j\in\Delta$.

Finally, $\widetilde\V\subseteq\widetilde p_*(\L(\mu^{(\delta)}))$, where $\widetilde p\: X\times T^{(\delta)}\times B\to T^{(\delta)}\times B$ is the projection. Indeed, 
\[
\widetilde p_*\big(\L(\mu^{(\delta)})\big)\subseteq\widetilde p_*\big(\L(D\times B)(\mu^{(\delta)})\big)\subseteq p^*\W
\]
with locally free quotients, if $D$ is sufficiently ample. Since $\V^{(i)}\subseteq p_{B*}\L^{(i)}$, it follows from \eqref{eq:VtildeV} that 
$\widetilde\V\subseteq\widetilde p_*(\L(D\times B)(\mu^{(\delta)}))$ over a dense open subscheme of $T^{(\delta)}\times B$, and then over the whole scheme. Also, the composition 
\[
\widetilde\V\longrightarrow\widetilde p_*\big(\L(D\times B)(\mu^{(\delta)})\big)\longrightarrow
\widetilde p_*\big(\L(D\times B)(\mu^{(\delta)})\big|_{D\times T^{(\delta)}\times B}\big)
\]
is zero because it restricts to zero on a dense open subscheme of $T^{(\delta)}\times B$, thus showing that $\widetilde\V\subseteq\widetilde p_*(\L(\mu^{(\delta)}))$. 

Thus we get a family of twisted generalized linear series $(\L,\widetilde\V)$ of degree $d$ and rank $r$ on $X$ along $\mu^{(\delta)}\times 1_B$. From the construction, $(\L,\widetilde\V)$ is clearly everywhere invariant. Furthermore, it is everywhere nonconstant because $N$ is minimal. So
$(\L,\widetilde\V)$ is a family of continuous linear series along the family of chain maps $1_B\times\mu^{(\delta)}$. It is clear that the association of $(\L,\widetilde\V)$ to $(\L;\V^{(i)},i\in\Delta)$ commutes with base change, thus defining the morphism 
$\Psi_\delta^N$ in the statement of the theorem.

The action by $\cc^{\Delta}$ on $G^{r,*}_{d,\delta}(X;N)$ changes $(\L,\widetilde\V)$ to an equivalent continuous linear series, whence 
$\Psi_\delta^N$ is $\cc^{\Delta}$-invariant. Furthermore, the equivalence class of $(\L;\V^{(i)},i\in\Delta)$ can be recovered from $(\L,\widetilde\V)$, as 
$\L(\mu^{(\delta)})|_{X\times E_i^{(\delta)}\times B}\cong\L^{(i)}$ for each $i\in\Delta$ and, under a certain choice of isomorphism, $\V^{(i)}$ is the image of $\widetilde\V|_{E^{(\delta)}_i\times B}$. It follows that each nonempty fiber of $\Psi_\delta^N$ is an equivalence class of 
level-$\delta$ limit linear series, as stated in the theorem.

Finally, let $(L,\V)$ be a continuous linear series of degree $d$ and rank $r$ on $X$ along a chain map $\mu\:S\to T$. For each $i=1,\dots,d$, let 
$\delta_i$ be the number of irreducible components of $\mu^{-1}(N_i)$. 
Put $\delta:=(\delta_1,\dots,\delta_d)$ and $\Delta:=\Delta(\delta)$. 
There is a unique isomorphism between the poset of irreducible components of $S$ and $\Delta$. For each $i\in\Delta$, let $S_i$ be the corresponding component. Also, let $N_j:=S_i\cap S_j$ for each consecutive $i,j\in\Delta$. 

For each $i\in\Delta$, fix a point $E'_i\in S^*_i$. Put 
$V^{(i)}:=\V|_{E'_i}$ and 
$L^{(i)}:=L(\mu)|_{X\times E'_i}$. Then 
$(L^{(i)},V^{(i)})$ is a linear series on $X$.

Observe that the $L^{(i)}$ form the twist $\delta$-sequence associated 
to $L$. Also, let $U_1:=\Gamma(Y,L|_Y)$ and 
$U_2:=\Gamma(Z,L|_Z(dP))$, and put $U:=U_1\oplus U_2$. Since $(L,\V)$ is everywhere invariant,
$$
\V_{N_j}=\lim_{c\to\infty}(c^{-1},1)\V|_{E'_i}=
\lim_{c\to 0}(c^{-1},1)\V|_{E'_j}
$$
as subspaces of $U$ for each 
consecutive $i,j\in\Delta$. But Proposition~\ref{hV} yields
$$
\lim_{c\to\infty}(c^{-1},1)\V|_{E'_i}=
\iota_1^{-1}(\V|_{E'_i})\oplus\rho_2(\V_{E'_i})
\text{\  \ and \  \ }
\lim_{c\to 0}(c^{-1},1)\V|_{E'_j}=
\rho_1(\V|_{E'_j})\oplus\iota_2^{-1}(\V|_{E'_j}),
$$
where $\iota_\ell\: U_\ell\to U$ is the natural injection and 
$\rho_\ell\: U\to U_\ell$ is the natural surjection for $\ell=1,2$. 
Thus the $V^{(i)}$ satisfy Conditions~\eqref{exact2}, 
and hence $(L,\V^{(i)},i\in\Delta)$ is an exact 
level-$\delta$ limit linear series on $X$. Finally, since $(L,\V)$ is everywhere nonconstant, $(L,\V^{(i)},i\in\Delta)$ is minimal, and is thus parameterized by a point on $G^{r}_{d,\delta}(X;N)$ for a unique exact, minimal $N\:\Delta\to\mathbf Z^3$. 

Since $(L,\V)$ is everywhere invariant, and since \eqref{cvary} holds for each $i\in\Delta$, it follows that there is a $T$-isomorphism $T^\Delta\to S$ taking $E_i^\Delta$ to $E'_i$ for each $i\in\Delta$ and inducing an equivalence between $(L,\widetilde\V)$, the continuous linear series corresponding to $(L,\V^{(i)},i\in\Delta)$ under $\Psi_\delta^N$, and $(L,\V)$. 
\end{proof}

\begin{prop}\label{toOss} There is a natural $J$-map $\Phi\: G^r_d(X)\to G^r_{d,\mathbf 1}(X)$ such that $\Phi\Psi_{\mathbf 1}^N$ is the open embedding of $G^{r}_{d,\mathbf 1}(X;N)$ in $G^r_{d,\mathbf 1}(X)$ for each exact, minimal $N$.
\end{prop}

\begin{proof} Let $(\L,\V)$ be a family of continuous linear series along a family of chain maps $\gamma\:\mathcal S\to T\times B$. Then $\gamma$ restricts to an isomorphism $\gamma^{-1}(E_i\times B)\to E_i\times B$ for each $i=0,\dots,d$. Under the natural identification of $E_i\times B$ with $B$, the restriction of 
$\L(\gamma)$ over $\gamma^{-1}(E_i\times B)$ is thus isomorphic to
$\L^{(i)}:=\L\otimes p_1^*\mathcal O_X^{(i)}$, where $p_1\:X\times B\to X$ is the projection, and the restriction of $\V$ to $E_i\times B$ is 
thus isomorphic to a subsheaf $\V^{(i)}$ of $p_{2*}\L^{(i)}$ such that 
$(\L^{(i)},\V^{(i)})$ is a family of linear series on $X\times B/B$, where $p_2\:X\times B\to B$ is the projection. The compatibility of the $\V^{(i)}$, that is, the fact that the image of $\V^{(i)}$ under the map induced by the natural map $\mathcal O_X^{(i)}\to\mathcal O_X^{(i\pm 1)}$ is contained in
$\V^{(i\pm 1)}$ for each $i$, follows from the argument given in the proof of the last statement of Theorem~\ref{allevels}. Thus we have a $B$-point of $G^{r}_{d,\mathbf 1}(X)$. The association commutes with base change, whence we obtain the $J$-map $\Phi$. Its property is a consequence of the construction of $\Psi_{\mathbf 1}^N$ in the proof of Theorem~\ref{allevels}.
\end{proof}

%\section*{Statements and Declarations}

%\subsection*{Ethical Approval.} Not applicable.

%\subsection*{Data Availability.} This manuscript has no associated data.

%\subsection*{Competing Interest.} The authors declare that they have no competing interests.



%{\color{red} Não modifiquei o resto. Vou tentar algo mais, que vou colocar na próxima seção, que pode via a substituir a discussão abaixo.}

%In conclusion, by Proposition 5.4 in \cite{ENR}, associated to each
%stable limit linear series $[f]\in G_d^r(X)^{\text{st}}$ we have a
%subscheme $\p(\mathfrak{g}_f)\in A_d^{-1}(L)$ in the fiber of
%the degree $d$ Abel map $A_d$. In particular,
%\begin{cor}\label{resol-map}
%For any $\mathfrak{g}=(L,V_0,\ldots,V_d)\in G_d^{r,{\rm Oss}}(X)$
%there exist a stable limit linear series
%$[f_{\widetilde{\mathfrak{g}}}]\in G_d^r(X)^{\text{st}}$
%such that $\widetilde{\mathfrak{g}}\in G_{d,\delta}^{r,\ast}(X)$
%and $\rho_{1,\delta}(\widetilde{\mathfrak{g}})=\mathfrak{g}$ for
%some $\delta$ and whose subscheme associated in the fiber of the
%Abel map $\p(\widetilde{\mathfrak{g}}_f)\in A_d^{-1}(L)$
%has ``correct'' Hilbert polynomial, i.e., the Hilbert polynomial
%$P(s,t)=\binom{s+t+r}{r}$ of the diagonal of $\p^r\times\p^r$.
%\end{cor}
%\begin{proof}
%In fact, by Theorem 5.5 in (~\cite{ENR}) there exists an
%exact level-$\delta$ limit linear series
%$\widetilde{\mathfrak{g}}\in G_{d,\delta}^{r,\ast}(X;N)$ for some %$N$.
%By Theorem \ref{allevels}, $\widetilde{\mathfrak{g}}$ corresponds to
%$[f_{\widetilde{\mathfrak{g}}}]$ the unique stable limit
%linear series on $X$, which by exactness and Proposition 5.4 in
%(~\cite{ENR}), the subscheme $\p(\widetilde{\mathfrak{g}}_f)$ has
%Hilbert polynomial $P(s,t)=\binom{s+t+r}{r}$.
%\end{proof}

\bibliographystyle{plain}
%\bibliography{../juca}
\begin{thebibliography}{11}

\bibitem{AB}
Amini, O., Baker, M.,
\newblock Linear series on metrized complexes of algebraic curves.
\newblock {\em Math. Ann.} \textbf{362} (2015), 55--106. .

%\bibitem{AEa}
%Amini, O., Esteves, E.,
%\newblock Voronoi tilings, toric arrangements and degenerations of lines bundles, I. 
%\newblock {\em ArXiv Preprint} 2012.15620 (2020).

\bibitem{AEb}
Amini, O., Esteves, E.,
\newblock Voronoi tilings, toric arrangements and degenerations of lines bundles, II. 
\newblock {\em ArXiv Preprint} 2012.15634 (2020).

\bibitem{AEc}
Amini, O., Esteves, E.,
\newblock Voronoi tilings, toric arrangements and degenerations of lines bundles, III. 
\newblock {\em ArXiv Preprint} 2102.02149 (2021).

\bibitem{AE1}
Amini, O., Esteves, E.,
\newblock Tropicalization of linear series and tilings by polymatroids. 
\newblock {\em ArXiv Preprint} 2405.04306 (2024).

\bibitem{AE2}
Amini, O., Esteves, E.,
\newblock Limit canonical series.
\newblock {\em ArXiv Preprint} 2410.13543v2 (2025).

\bibitem{BN}
Baker, M., Norine, S.,
\newblock Riemann--Roch and Abel--Jacobi theory on a finite graph.
\newblock {\em Adv. Math.} \textbf{215} (2007), 766--788.

\bibitem{Br}
Brion, M.,
\newblock Invariant Hilbert schemes. 
\newblock In: Handbook of Moduli, Farkas, G., Morrison, I. (eds),  {\em Advance Lectures in Mathematics} \textbf{24}. International Press, Somerville (2013), 64--117.

%\bibitem{ACGH}
%Arbarello, E., Cornalba, M., Griffiths, P. and Harris, J.,
%\newblock Geometry of algebraic curves, Vol.I.
%\newblock {\em Grundlehren der Mathemaitschen Wissenchaften}, Volume 267, Springer-Verlag, New York, (1985).

\bibitem{Cast}
Castorena, A., Martín-Lopez, A., Teixidor i Bigas, M.,
\newblock Invariants of the Brill--Noether curve.
\newblock  {\em Adv. Geom.} \textbf{17} (2017), 39--52.

\bibitem{CEP}
Coelho, J., Esteves, E., Pacini, M.,
\newblock Degree-2 Abel maps for nodal curves.
\newblock {\em IMRN} \textbf{2016} (2016), 2912--2973.

\bibitem{CDP}
Cools, F. et al.,
\newblock A tropical proof of the Brill--Noether theorem.
\newblock {\em Adv. Math.} \textbf{230} (2012), 759--776.

\bibitem{Cot}
Cotterill, E.,
\newblock Geometry of curves with exceptional secant planes: linear series along the general curve.
\newblock  {\em Math. Z.} \textbf{267} (2011), 549--582.

\bibitem{EH1}
Eisenbud, D., Harris, J.,
\newblock Limit linear series: Basic theory.
\newblock {\em Invent. Math.} \textbf{85} (1986), 337--371.

\bibitem{EH2}
Eisenbud, D., Harris, J.,
\newblock Existence, decomposition and limits of certain Weierstrass points.
\newblock {\em Invent. math.} \textbf{87} (1987), 495--515.

\bibitem{EH3}
Eisenbud, D., Harris, J.,
\newblock On the Kodaira dimension of the moduli space of curves.
\newblock {\em Invent. math.} \textbf{67} (1987), 23--86.

\bibitem{EH4}
Eisenbud, D., Harris, J.,
\newblock Irreducibility of some families of linear series with Brill--Noether number $-1$.
\newblock {\em Ann. scient. \'Ec. Norm. Sup.} \textbf{22} (1989), 33--53.

\bibitem{Ee}
Esteves, E.,
\newblock Linear systems and ramification points on reducible nodal curves.
\newblock {\em Matem\'atica Contempor\^anea}, \textbf{14} (1998), 21--35.

\bibitem{EM}
Esteves, E., Medeiros, N.,
\newblock Limit canonical systems on curves with two components.
\newblock {\em Invent. Math.} \textbf{149} (2002), 267--338.

\bibitem{ENR}
Esteves, E., Nigro, A., Rizzo, P.,
\newblock Level-$\delta$ limit linear series.
\newblock {\em Math. Nachr.} \textbf{291} (2018), 827--847.

\bibitem{E-Oss}
Esteves, E., Osserman, B.,
\newblock Abel maps and limit linear series.
\newblock {\em Rend. Circ. Mat. Palermo} \textbf{62} (2013), 79--95.

\bibitem{EP}
Esteves, E., Pacini, M.,
\newblock Semistable modifications of families of curves and compactified Jacobians.
\newblock {\em Ark. Mat.} \textbf{54} (2016), 55--83.

\bibitem{ES}
Esteves, E., Salehyan, P.,
\newblock Limit Weierstrass points on nodal reducible curves.
\newblock {\em Trans. Amer. Math. Soc.} \textbf{359} (2007), 5035--5056.

\bibitem{ESVI}
Esteves, E., Santos, R., Vital, E., 
\newblock Quiver representations arising from degenerations of linear series, I. 
\newblock {\em  Comm. Algebra} (2025), 1–30. https://doi.org/10.1080/00927872.2025.2530496

\bibitem{ESVII}
Esteves, E., Santos, R., Vital, E., 
\newblock Quiver representations arising from degenerations of linear series, II. 
\newblock {\em ArXiv Preprint} 2112.00514v2 (2022).

%\bibitem{GIT}
%Fogarty, J., Kirwan, F., Mumford, D.,
%\newblock Geometric Invariant Theory. Volume 34.
%\newblock {\em Results in mathematics and related areas}, Springer Verlag, Third edition. (1994).

\bibitem{Fark1}
Farkas, G., Lian, C.,
\newblock Linear series on general curves with prescribed incidence conditions.
\newblock {\em J. Inst. Math. Jussieu} \textbf{22} (2023), 2857--2877.

\bibitem{Fark2}
Farkas, G., Verra, A.,
\newblock The geometry of the moduli space of odd spin curves.
\newblock {\em Ann. Math.} \textbf{180} (2014), 927--970.

\bibitem{GKZ}
Gelfand, I., Kapranov, M., Zelevinsky, A.,
\newblock Discriminants, resultants, and multidimensional determinants. Springer (1994).

\bibitem{GW22}
Giansiracusa, N., Wu, X.,
\newblock Chow quotients of {G}rassmannians by diagonal subtori.
\newblock \newblock In: Facets of Algebraic Geometry, A Collection in Honor of William Fulton's 80th Birthday, Volume I, {\em London Mathematical Society Lecture Note Series} \textbf{ 472}. Cambridge University Press (2022), 241--264. 

%\bibitem{F}
%Fulton, W.,
%\newblock Intersection theory.
%\newblock {\em Modern surveys in mathematics}, Springer-Verlag, First edition, (1984).

%\bibitem{FPnotes}
%Fulton, W., Pandharipande, R.,
%\newblock Notes on stable maps and quantum cohomology.
%\newblock {\em Algebraic geometry, Santa Cruz, 1995}, Amer. Math. Soc., (1997).

%\bibitem{EGA}
%, James E. 
%\newblock Linear algebraic groups. 
%\newblock {\em Springer Science \& Business Media}, 
%\textbf{ Vol. 21}, (2012).

\bibitem{JP1}
Jensen, D., Payne, S.,
\newblock Tropical independence, I: shapes of divisors and a proof of the Gieseker--Petri theorem.
\newblock {\em Algebra Number Theory} \textbf{8} (2014), 2043--2066.

\bibitem{JP2}
Jensen, D., Payne, S.,
\newblock Tropical independence, II: The maximal rank conjecture for quadrics.
\newblock {\em Algebra Number Theory} \textbf{10} (2016), 1601--1640.

\bibitem{Kap-chow}
Kapranov, M.,
\newblock Chow quotients of {G}rassmannian, {I}.
\newblock In: I.~M.~Gelfand Seminar, {\em Adv. Soviet Math.} \textbf{16} (1993), 29--110.


%\bibitem{Kn}
%Knudsen, F.,
%\newblock The projectivity of the moduli space of stable curves, II: The stacks $M_{g,n}$
%\newblock {\em Math. Scand.}, 
%\textbf{52.2}, pp: 161-199. (1983).

\bibitem{Liu}
Liu, F.,
\newblock Moduli of crude limit linear series.
\newblock {\em IMRN} \textbf{2009} (2009), 4032--4050.

\bibitem{M}
Mu\~noz, G.,
\newblock Limit linear series for curves of compact type with three irreducible components.
\newblock {\em Comm. Algebra.} \textbf{48} (2020), 4457--4482.

\bibitem{Oss1}
Osserman, B.,
\newblock A limit linear series moduli scheme.
\newblock {\em Ann. I. Fourier} \textbf{56} (2006), 1165--1205.

\bibitem{Oss2}
Osserman, B.,
\newblock Limit linear series for curves not of compact type.
\newblock {\em J. Reine Angew. Math.} \textbf{753} (2019), 57--88.

\bibitem{OssT} 
Osserman, B., 
\newblock Limit linear series and the Amini--Baker construction.
\newblock {\em Math. Z.} \textbf{293} (2019) 339--369.

\bibitem{Ran}
Ran, Z.,
\newblock Degeneration of Linear Systems.
\newblock {\em Unpublished}, (1985).

\bibitem{Rizzo}
Rizzo, P.,
\newblock Level-delta and stable limit linear series on singular curves.
\newblock {\em Tese de Doutorado}, IMPA, Rio de Janeiro (2013).

%\bibitem{Ro}
%Rocha, D.,
%\newblock Abel maps in the Hilbert scheme of curves and limit linear series
%\newblock {\em Tese de Doutorado, IMPA, Rio de Janeiro (in progress)}

%\bibitem{Stacks}
%The {Stacks Project Authors},
%\newblock Stacks Project.
%\newblock {\em https://stacks.math.columbia.edu}, (2018).

\bibitem{Tar}
Tarasca, N.,
\newblock Brill--Noether loci in codimension two.
\newblock {\em Compos. Math.} \textbf{149} (2013), 1535--1668.

\bibitem{Tha}
Thaddeus, M., 
\newblock Complete colineations revisited.
\newblock {\em Math. Ann.} \textbf{315} (1999), 469--495.

\bibitem{VE}
Vital, E.,
\newblock Degenerations of linear series to curves with three components, using quiver representations.
\newblock {\em Tese de Doutorado}, IMPA, Rio de Janeiro (2022).

\end{thebibliography}

\end{document} 


%%%%%%%%%%%
\begin{prop} Let $\mu\: S\to T$ be a chain map. Then each node of $S$ is sent to 
$T_0\cup T_d\cup\{N_1,\dots,N_{d-1}\}$.
\end{prop}

\begin{proof} Let $s$ be a node of $S$, and $R_1$ and $R_2$ the irreducible components of $S$ intersecting at $P$. Suppose by contradiction that $\mu(s)\in T_i^*$ for some $i\in\{1,\dots,d-1\}$. Let $S'$ be the largest subchain of $S$ containing $R_1$ but not $R_2$ and $S''$ the largest subchain of $S$ containing $R_2$ but nor $R_1$. Then $T':=\mu(S')$ and $T'':=\mu(S'')$ are either points or subchains of $T$. Since $\mu_{\ast}[S]=[T]$, they must cover $T$ and must not have any component in common. But $\mu(s)\in T'\cap T''$, whence 
$T'\cap T''\cap T_ i^*\neq\emptyset$. Thus either $T'=T$ or $T''=T$. Suppose $T'=T$. Let $U$ and $V$ be irreducible components of $S'$ such that 
$\mu(U)=T_0$ and $\mu(V)=T_d$. Then the smallest subchain of $S'$ containing $U$ and $V$ covers $T$. It follows that the smallest subchain of $S'$ containing 
$R_1$ and $U$ or $V$ collapses under $\mu$. Since $\mu(s)\in T_i^*$, it collapses to a point on $T_i^*$, but also on $T_0\cup T_d$, a contradiction with 
$i\in\{1,\dots,d-1\}$. 
\end{proof}

%%%%%%%%%%%%%%%%%%%%%%%
\begin{prop} Let $(L,\V)$ be a twisted generalized linear series on $X$ along a chain map $\mu\:S\to T$. Let $R$ be an irreducible component of $S$. If 
$(L,\V)$ is invariant along $R$ then, given $s_1,s_2$ of $S$ on an irreducible component of $S$ with $t_1,t_2\in T^*$, there is 
$c\in\cc(k)$ such that $t_2=\sigma(c,t_1)$ and \eqref{diags1s2} commutes, where $t_i:=\mu(s_i)$ for $i=1,2$. Furthermore, if $(L,\V)$ is also nonconstant along $R$ and $\mu(R)$ is a point, then $\mu(R)\in T-T^*$. 

then for each node $R$ of $S$ or irreducible component $R$ of $S$ such that $\mu(R)$ is a point, $\mu(R)\in T-T^*$. 
\end{prop}

\begin{proof} Since $s_1,s_2$ lie on an irreducible component $R$ of $S$, we must have $t_1,t_2\in T_i^*$ for some $i$. Then $t_2=\sigma(c,t_1)$ for a unique 
$c\in\cc(k)$. Since $(L,\V)$ is everywhere invariant, it follows from continuity that $\V_{s_i}=(c_i^{-1},1)\V_s$ for a general $s\in R$ for $i=1,2$, where $c_i\in\cc(k)$ is defined by $t_i=\sigma(c_i,t)$, for $t:=\mu(s)$. But then 
$\V_{s_2}=(c^{-1},1)\V_{s_1}$ where $c:=c_2/c_1$, and $\sigma(c,t_1)=\sigma(cc_1,t)=\sigma(c_2,t)=t_2$.

As for the second statement, let $R$ be an irreducible component of $S$ 
collapsed by $\mu$, and put $t:=\mu(R)$. Then $t\in T-T^*$. 
Indeed, since $(L,\V)$ is everywhere invariant, for general $s_1$ and $s_2$ on $R$ there is $c\in\cc(k)$ such that $t=\sigma(c,t)$ and \eqref{diags1s2} commutes. 
If $t\in T^*$, it would follow that $c=1$ and thus $\V_{s_1}=\V_{s_2}$ for general $s_1$ and $s_2$ on $R$, whence for all $s_1,s_2\in R$ by continuity, contradicting the fact that $(l,\V)$ is everywhere nonconstant.

It follows that $\mu$ sends the nodes of $S$ to
$T-T^*$. Indeed, if a node $R$ of $S$ is sent to $T^*$, then we would have two intersecting irreducible components $C_1$ and $C_2$ of $S$ whose image in $T$ contains a point on $T_i^*$ for some $i$. But then $\mu$ cannot collapse either $C_1$ or $C_2$, and must map each onto $T_i$, contradicting $\mu_*[S]=[T]$. 
\end{proof}

%%%%%%%%%%%%%%%%%%%%

\begin{prop}\label{stabconst} An equivalence class of families of twisted generalized linear series along a family of chain maps 
is the pullback of a unique equivalence class of families of everywhere nonconstant twisted generalized linear series along an intermediate family of chain maps. If the former is a class of families of everywhere invariant series, then the latter is a class of families of continuous linear series. 
\end{prop}

\begin{proof} Let $(\L,\V)$ be a family of twisted generalized linear series along a family of chain maps $\gamma\:\mathcal S\to B\times T$. Clearly, 
$$
\L(1_{B\times T})\subseteq (1_X\times q)^*\big(\L_{Y\times B}(D\times B)\oplus\L|_{Z\times B}((dP+D)\times B)\big)
$$
with $B\times T$-flat cokernel $\mathcal C$ for any effective divisor $D$ of $X-P$, where $q\:B\times T\to B$ is the projection. 
Then $\mathcal E:=p'_*\mathcal C$ and $\W:=p_*(\L_{Y\times B}(D\times B)\oplus\L|_{Z\times B}((dP+D)\times B))$, for sufficiently ample $D$, are locally free sheaf whose formations commute with base change, where $p\:X\times B\to B$ and $p'\:X\times T\times B\to T\times B$ are the projections. We obtain an exact sequence 
\begin{equation}\label{sesL}
0\to p_{\mathcal R*}\L(\tau)\to\tau_1^*\W\to\tau^*\mathcal E
\end{equation}
for each family of chain maps $\tau=(\tau_1,\tau_2)\:\mathcal R\to B\times T$.
Furthermore,
since $(\L(\gamma),\V)$ is a family of generalized linear series, $\V\subseteq\pi^*\W$ with locally free quotient. Denoting by $r$ the rank of $(\L,\V)$, we have a well-defined map 
$B$-map $f\:\mathcal S\to\text{Grass}_B(r+1,\W)$. 

Let $\pi'\:\mathcal S'\to B$ be the stabilization of $\pi$ for $f\times\mu$. More precisely,
$$
\mathcal S':=\text{Proj}_B\Big(\bigoplus_{i\geq 0}\pi_*(\mathcal A^{\otimes n})\Big),
$$
where $\mathcal A$ is the pullback under $f\times\mu$ of a relatively ample invertible sheaf on $\text{Grass}(r+1,\W)\times T/B$. 
There are a $B$-map $\psi\:\mathcal S\to\mathcal S'$ collapsing the irreducible components in the fibers of $\pi$ collapsing under $f\times\mu$ and a map $f'\times\mu'\:\mathcal S'\to\text{Grass}_B(r+1,\W)\times T$ such that $f=f'\psi$ and $\mu=\mu'\psi$. Thus $\gamma':=(\pi',\mu')$ satisfies $\gamma=\gamma'\psi$ and 
is an intermediate family of chain maps. Furthermore, $\V=\psi^*\V'$ for a subsheaf $\V'\subseteq(\pi')^*\W$ with locally free quotient. 

By the exactness of \eqref{sesL}, the induced composition 
$\V'\to (\pi')^*\W\to(\gamma')^*\mathcal E$ pulls back under $\psi$ to the zero map, since $\V\subseteq p_{\mathcal S*}\L(\gamma)$. Since $\V'$ and $\mathcal E$ are locally free and 
$\psi_*\mathcal O_{\mathcal S}=\mathcal O_{\mathcal S'}$, it follows that the composition itself is zero. Thus $\V'\subseteq p_{\mathcal S'*}(\L(\gamma'))$, and hence $(\L,\V')$ is a family of twisted generalized linear series along $\gamma'$ pulling back to $(\L,\V)$ under $\psi$.

Since $f'\times\mu'$ does not collapse any irreducible component of any fiber of $\pi'$, the family $(\L,\V')$ is one of everywhere nonconstant series. If $(\L,\V)$ is a family of everywhere invariant series, since $(\L,\V')$ pulls back to $(\L,\V)$, also $(\L,\V')$ is a family of everywhere invariant series, hence a family of continuous linear series.

{\color{red} Precisa terminar a prova, explicar melhor.} {\color{red} Falta também discutir unicidade e equivalência.}
\end{proof}

%%%%%%%%%%%%%%%

\begin{prop}\label{stabconstflag} An equivalence class of families of everywhere invariant twisted flags of generalized linear series along a family of chain maps 
$\gamma=(\pi,\mu)\:\mathcal S\to B\times T$ is the pullback of a unique equivalence class of families of continuous flags of linear series along an intermediate family of chain maps. 
\end{prop}

\begin{proof} The proof is mutatis mutandi 
that of Proposition~\ref{stabconst}, the only difference being that $H^r_d(X)$ is replaced by the $(J\times T)$-product of the $H^{r_\ell}_d(X)$.
\end{proof}

%%%%%%%%%%%%%%%%%

The maps $\rho_{\Delta,\Delta'}$ can be used to glue the $G^{r,*}_{d,\Delta}(X)$ to a scheme. 

Indeed, let $\mathfrak D$ be the collection of finite subsets of $[0,d]$ 
containing $\{0,1,\dots,d\}$. For each $\Delta_1,\Delta_2\in\mathfrak D$, let $\Delta\in\mathfrak D$ with $\Delta\supseteq\Delta_1\cup\Delta_2$, and put
$$
U^{\Delta}_{\Delta_1,\Delta_2}:=\rho_{\Delta,\Delta_1}^{-1}(G^{r,\ast}_{d,\Delta_1}(X))\cap \rho_{\Delta,\Delta_2}^{-1}(G^{r,\ast}_{d,\Delta_2}(X))\subseteq 
G^{r,\ast}_{d,\Delta}(X),
$$
and
$$
U_{\Delta_1,\Delta_2}:=\rho_{\Delta,\Delta_1}(U^{\Delta}_{\Delta_1,\Delta_2})
\quad\text{and}\quad
U_{\Delta_2,\Delta_1}:=\rho_{\Delta,\Delta_2}(U^{\Delta}_{\Delta_1,\Delta_2}).
$$
Then $U_{\Delta_1,\Delta_2}$ is open in $G^{r,\ast}_{d,\Delta_1}$ and 
$\rho_{\Delta,\Delta_1}$ restricts to an isomorphism 
$\varphi^{\Delta}_{\Delta_1,\Delta_2}\:U^{\Delta}_{\Delta_1,\Delta_2}\to U_{\Delta_1,\Delta_2}$. Similarly, we have an isomorphism 
$\varphi^{\Delta}_{\Delta_2,\Delta_1}\:U^{\Delta}_{\Delta_1,\Delta_2}\to U_{\Delta_2,\Delta_1}$. Put 
$$
\varphi_{\Delta_1,\Delta_2}:=\varphi^{\Delta}_{\Delta_2,\Delta_1}(\varphi^{\Delta}_{\Delta_1,\Delta_2})^{-1}.
$$
The reader may observe that $\varphi_{\Delta_1,\Delta_2}$ does not depend on the choice of $\Delta$, and that the $\varphi_{\Delta_1,\Delta_2}$ form a glueing data for the $G^{r,*}_{d,\Delta}(X)$ for $\Delta\in\mathfrak D$. 

Let $G^{r,*}_{d,\mathfrak D}(X)$ be the glueing of the $G^{r,*}_{d,\Delta}(X)$, and let $\iota_{\Delta}\:G^{r,*}_{d,\Delta}(X)\to G^{r,*}_{d,\mathfrak D}(X)$ be the associated open embedding for each $\Delta\in\mathfrak D$. 

%%%%%%%%%%%%%%%%%%

Furthermore, the diagram
\begin{equation*}\label{diagdelta}
\xymatrix{
\rho_{\Delta,\Delta'}^{-1}(G_{d,\Delta'}^{r,\ast}(X))
\ar@{^{(}->}[r]\ar[d]^-{\rho_{\Delta,\Delta'}} &
G_{d,\Delta}^{r,\ast}(X)\ar[r]^-{\Psi_{\Delta}} &
G_d^r(X)\\
G_{d,\Delta'}^{r,\ast}(X)\ar[urr]_-{\Psi_{\Delta'}} & &
}
\end{equation*}
commutes for each $\Delta'\in\mathfrak D$ with $\Delta'\supseteq\Delta$. 

%%%%%%%%%%%%%%%%%%%

Furthermore, given $\Delta'\in\mathfrak D$ with $\Delta'\supseteq\Delta$, Diagram~\eqref{diagdelta} 
commutes, because if the morphism $B\to G^{r,*}_{d,\Delta}(X)$ given by the family of level-$\Delta$ limit linear series factors through $\rho_{\Delta,\Delta'}^{-1}(G_{d,\Delta'}^{r,\ast}(X))$, then $\overline{\Delta}=\overline{\Delta'}$.

The commutativeness of \eqref{diagdelta} for each $\Delta,\Delta'\in\mathfrak D$ with $\Delta'\supseteq\Delta$ yields, by gluing, the morphism $\Psi$ such that 
$\Psi_\Delta=\Psi\iota_\Delta$ for each $\Delta\in\mathfrak D$.

We will prove the universality of $\Psi$ now. Let $\xi\: B\to G^r_d(X)$ be the induced map. We need only prove that $\xi$ lifts \'etale locally to a map to $G^{r,*}_{d,\mathfrak D}(X)$, uniquely defined up to the action of $\mathbb G_m^\Delta$. 

Let $(\L,\V)$ be a family of continuous linear series of degree $d$ and rank $r$ on $X$ along a family of chain maps $\gamma:=(\pi,\mu)\:\mathcal S\to B\times T$ defining $\xi$. Let $b\in B(k)$. Let $\Delta\in\mathfrak D$ such that the number of irreducible components of $\mu_b^{-1}(N_i)$ is $\#\Delta\cap(i-1,i)$ for each $i=1,\dots,d$. Thus there is a unique isomorphism between the poset of irreducible components of $\mathcal S_b$ and $\Delta$. For each $i\in\Delta$, let $\mathcal S_{b,i}$ be the corresponding component. Also, let $N_{b,j}=\mathcal S_{b,i}\cap\mathcal S_{b,j}$ for each consecutive $i,j\in\Delta$. 

There are an \'etale map $\epsilon\: B'\to B$ and $b'\in B'$ with $\epsilon(b')=b$, and quasi-sections $s_i\:B'\to\mathcal S$ of $\pi$ through its smooth locus such that $s_i(b')$ lies on $\mathcal S_{b,i}$ for each $i\in\Delta$. Let $E_{i}$ be the image of $s_i$ and put 
$\V^{(i)}:=\pi_{*}(\V|_{E_{i}})$ and 
$\L^{(i)}:=(1_X\times\pi_b)_*(\L(\gamma)|_{X\times E_{b,i}})$ for each $i$, where $\pi'\:\mathcal S\times_BB'\to B'$ is the induced map. Then 
$(\L^{(i)},\V^{(i)})$ is a family of linear series on $X\times B'/B'$.

Observe that $\V^{(i)}_x=\V_{s_j(x)}$ for each $x\in B'(k)$. Also, the $\L^{(i)}_x$ form the twist $\Delta$-sequence associated 
to $\L_x$. Let $U_{1,x}:=\Gamma(Y,\L_x|_Y)$ and 
$U_{2,x}:=\Gamma(Z,\L_x|_Z(dP))$, and put $U_x:=U_{1,x}\oplus U_{2,x}$. Since $(\L,\V)$ is everywhere invariant,
$$
\V_{N_{x,j}}=\lim_{c\to\infty}(c^{-1},1)\V_{s_i(x)}=
\lim_{c\to 0}(c^{-1},1)\V_{s_j(x)}
$$
as subspaces of $U_x$ for each 
consecutive $i,j\in\Delta$. But Proposition~\ref{hV} yields
$$
\lim_{c\to\infty}(c^{-1},1)\V_{s_i(x)}=
\iota_{1,x}^{-1}(\V_{s_i(x)})\oplus\rho_{2,x}(\V_{s_i(x)})
\text{\  \ and \  \ }
\lim_{c\to 0}(c^{-1},1)\V_{s_j(x)}=
\rho_{1,x}(\V_{s_j(x)})\oplus\iota_{2,x}^{-1}(\V_{s_j(x)}),
$$
where $\iota_{\ell,x}\: U_{\ell,x}\to U_x$ is the natural injection and 
$\rho_{\ell,x}\: U_x\to U_{\ell,x}$ is the natural surjection for $\ell=1,2$. Thus the $V^{(i)}_x$ satisfy Conditions~\eqref{exact2}, 
and hence $(\L_x,\V^{(i)}_x,i\in\Delta)$ is an exact 
level-$\Delta$ limit linear series on $X$. We have thus obtained a map 
$\phi_b\:B_b\to G^{r,\ast}_{d,\Delta}(X)$. 

{\color{red} Prova antiga, abaixo.}

Let now $N\:\Delta\to\z^3$ be exact and minimal. Set-theoretically, the continuous linear series along chain maps $\mu$ parameterized by 
$\Psi_\Delta(G^r_{d,\Delta}(X;N))$ are such that the number of irreducible components of $\mu^{-1}(N_i)$ is $\#\Delta\cap(i-1,i)$ for $i=1,\dots,d$. 

Consider thus a family $(\L,\V)$ of continuous linear series of degree $d$ and rank $r$ on $X$ along a family of chain maps $\gamma:=(\pi,\mu)\:\mathcal S\to B\times T$. Assume the number of irreducible components of $\mu_b^{-1}(N_i)$ is $\#\Delta\cap(i-1,i)$ for each $b\in B(k)$ and $i=1,\dots,d$. Let $\xi\: B\to G^r_d(X)$ be the induced map. 

{\color{red} O $\Delta$ e o $N$ tem que ser definidos a partir de $(\L,\V)$. O $N$ não vai ser minimal para todas as fibras em uma vizinhança de um ponto escolhido.}

For each $b\in B(k)$, there is a unique isomorphism between the poset of irreducible components of $\mathcal S_b$ and $\Delta$. For each $i\in\Delta$, let $\mathcal S_{b,i}$ be the corresponding component. Also, let $N_{y,j}=\mathcal S_{b,i}\cap\mathcal S_{b,j}$ for each consecutive 
$i,j\in\Delta$. 

For each $b\in B(k)$ there are an \'etale map $\epsilon_b\: B_b\to B$ with image containing $b$ and quasi-sections $s_i\:B_b\to\mathcal S$ of $\pi$ through its smooth locus such that $s_i(b)$ lies on $\mathcal S_{b,i}$ for each $i\in\Delta$. Since the number of irreducible components of the geometric fibers of $\pi$ does not vary, 
$s_i(x)\in\mathcal S_{x,i}^*$ for each $i\in\Delta$ and $x\in B_b(k)$. Let $E_{b,i}$ be the image of $s_i$ and put 
$\V^{(i)}:=\pi_{b,*}(\V|_{E_{b,i}})$ and 
$\L^{(i)}:=(1_X\times\pi_b)_*(\L(\gamma)|_{X\times E_{b,i}})$ for each $i$, where $\pi_b\:\pi^{-1}(B_b)\to B_b$ is the induced map. Then 
$(\L^{(i)},\V^{(i)})$ is a family of linear series on $X\times B/B$.

Observe that $\V^{(i)}_x=\V_{s_j(x)}$ for each $x\in B_b(k)$. Also, the $\L^{(i)}_x$ form the twist $\Delta$-sequence associated 
to $\L_x$. Let $U_{1,x}:=\Gamma(Y,\L_x|_Y)$ and 
$U_{2,x}:=\Gamma(Z,\L_x|_Z(dP))$, and put $U_x:=U_{1,x}\oplus U_{2,x}$. Since $(\L,\V)$ is everywhere invariant,
$$
\V_{N_{x,j}}=\lim_{c\to\infty}(c^{-1},1)\V_{s_i(x)}=
\lim_{c\to 0}(c^{-1},1)\V_{s_j(x)}
$$
as subspaces of $U_x$ for each 
consecutive $i,j\in\Delta$. But Proposition~\ref{hV} yields
$$
\lim_{c\to\infty}(c^{-1},1)\V_{s_i(x)}=
\iota_{1,x}^{-1}(\V_{s_i(x)})\oplus\rho_{2,x}(\V_{s_i(x)})
\text{\  \ and \  \ }
\lim_{c\to 0}(c^{-1},1)\V_{s_j(x)}=
\rho_{1,x}(\V_{s_j(x)})\oplus\iota_{2,x}^{-1}(\V_{s_j(x)}),
$$
where $\iota_{\ell,x}\: U_{\ell,x}\to U_x$ is the natural injection and 
$\rho_{\ell,x}\: U_x\to U_{\ell,x}$ is the natural surjection for $\ell=1,2$. Thus the $V^{(i)}_x$ satisfy Conditions~\eqref{exact2}, 
and hence $(\L_x,\V^{(i)}_x,i\in\Delta)$ is an exact 
level-$\Delta$ limit linear series on $X$. We have thus obtained a map 
$\phi_b\:B_b\to G^{r,\ast}_{d,\Delta}(X)$. 

{\color{red} A prova mostra que um certo mapa de funtores \é um epimorfismo. Argumentar que um contínuo equivalente a $(\L,\V)$ resulta em um discreto equivalente módulo ação do grupo. Falta argumentar que o mapa de funtores é um monomorfismo.}

Notice that $\Psi_\Delta\phi_b=\xi\epsilon_b$. Also, changing $(\L,\V)$ to an equivalent continuous linear series or changing the quasi-sections $s_i$ changes $\phi_b$ to an equivalent map under the action by $\cc^{\Delta}$. 

Suppose now that $\xi$ factors through $\Psi_\Delta(G^r_{d,\Delta}(X;N))$. Then 
$\phi_b$ factors through $G^r_{d,\Delta}(X;N)$. Let $\eta\:G^r_{d,\Delta}(X;N)\to M$ be a $\cc^\Delta$-invariant map to a scheme $M$. Then the composition $\eta\phi_b$ does not depend on the choices of the $s_i$. By descent, the maps $\eta\phi_b$ are induced by a map $\lambda\:B\to M$. If $\xi=\Psi_\Delta\xi'$ then 
$\lambda=\eta\xi'$. 

We have thus a map $\lambda\: G^r_d(X)\to M$
 such that $\eta=\lambda\Psi_\Delta$. 

 {\color{red} Algo estranho na argumentação, que não terminou, da última afirmaçpão.}

%%%%%%%%%%%%%%%%%

\subsection{The correspondence, I}

Let $T^\Delta$ be the chain of rational curves of length $\#\Delta$ obtained from $T$ by splitting the branches of $T$ at $N_i$ and connecting them by a chain of rational curves of length 
$\#\Delta\cap (i-1,i)$ for each $i=1,\dots,d$. Thus we index the components $T^\Delta_i$ of $T^\Delta$ with $i\in\Delta$. We obtain a chain map $\mu^\Delta\:T^\Delta\to T$. For each consecutive 
$i,j\in\Delta$, let $N^{(\delta)}_j$ be the intersection of $T^{(\delta)}_{i}$ with $T^{(\delta)}_j$. Let $N_0^\Delta$ and $N_\infty^\Delta$ on $T^\Delta$ above $N_0$ and $N_{d+1}$. 
%For convenience, put $N_{d+\delta^{-1}}^{(\delta)}:=N_\infty^{(\delta)}$. 

For each $i\in\Delta$, let $E_i^\Delta\in(T^\Delta_i)^*$ such that $\mu^\Delta(E_i^\Delta)=E_i$ if $i\in\mathbb Z$. The chain action on $(T,N_0,N_\infty)$ lifts to a unique chain action 
on $(T^\Delta,N_0^\Delta,N_\infty^\Delta)$ for which $(T^\Delta_i)^*$ is a torsor and $1\ast E^\Delta_i=E^\Delta_i$ for each $i\in\Delta$.

\begin{prop}\label{dlevel} For each invertible sheaf $L$ of bidegree $(d,0)$ on $X$ there is an isomorphism $L(\mu^\Delta)_{E_i^\Delta}\cong L^{(i)}$ for each $i\in\Delta$. Furthermore, for each everywhere invariant twisted
generalized linear series $(L,\V)$ along $\mu^\Delta$, letting $V^{(i)}$ be the image of 
$\V_{E_i^\Delta}$ under any isomorphism $L(\mu^\Delta)_{E_i^\Delta}\cong L^{(i)}$ for each $i\in\Delta$, the data
$\mathfrak g(L,\V):=(L;V^{(i)},i\in\Delta)$ is an exact 
level-$\Delta$ limit linear series on $X$ whose equivalence class depends only on that of $(L,\V)$. Conversely, for each exact level-$\Delta$ limit linear series $\mathfrak{g}=(L;V^{(i)},i\in\Delta)$ on $X$, there is a locally free subsheaf $\V\subseteq p_*(L(\mu^\Delta))$,
where $p\:X\times T^\Delta\to T^\Delta$ is the projection, such that:
\begin{enumerate}
\item $(L,\V)$ is an everywhere invariant twisted generalized
linear series along $\mu^\Delta$;
\item $\mathfrak g(L,\V)$ is equivalent to $\mathfrak g$.
\end{enumerate}
Furthermore, for any other locally free subsheaf $\V'\subseteq p_*(L(\mu^\Delta))$ satisfying the same two conditions as $\V$, the equivalence classes of $(L,\V')$ and $(L,\V)$ are equal.
\end{prop}

\begin{proof} {\color{red} Prova a ser removida após a conclusão da prova do Thm. 5.5} For each $i\in\Delta$, let $t_i:=\mu^{(\delta)}(E_i^\delta)$. If $i\in\mathbb Z$, 
then $t_i=E_i$ and $\F_{t_i}\cong\mathcal O^{(i)}_X$. If $i\not\in\mathbb Z$ then $t_i$ is not on the smoth locus of $T$, but $\F_{t_i}\cong\mathcal O^{(i)}_X$ all the same. Clearly, $L(\mu^\Delta)_{E_i^\Delta}=L\otimes\F_{t_i}$, and thus the first statement follows.

As for the second statement, let $(L,\V)$ be an invariant twisted
generalized linear series $(L,\V)$ along $\mu^\Delta$. Then 
$\V\subseteq p_*(L(\mu^\Delta))$. Now, 
$\F_t\subseteq \mathcal O_Y\oplus\mathcal O_Z(dP)$ for each $t\in T$, whence 
$p_*(L(\mu^\Delta))\subseteq U\otimes\mathcal O_{T^\Delta}$, where 
$U:=U_1\oplus U_2$, for 
$$
U_1:=\Gamma(Y,L|_Y)\quad\text{and}\quad U_2:=\Gamma(Z,L|_Z(dP)).
$$
Since $(L,\V)$ is everywhere invariant, 
$$
\V_{N^\Delta_i}=\lim_{c\to\infty}(c^{-1},1)\V_{E_{i-1}^\Delta}=
\lim_{c\to 0}(c^{-1},1)\V_{E_{i}^\Delta},
$$
as subspaces of $U_1\oplus U_2$ for each $i\in\Delta-\{0\}$. But Proposition~\ref{hV} yields
$$
\lim_{c\to 0}(c^{-1},1)\V_{E_{i-1}^\Delta}=
\rho_1(\V_{E_{i-1}^\Delta})\oplus\iota_2^{-1}(\V_{E_{i-1}^\Delta})
\text{\  \ and \  \ }
\lim_{c\to\infty}(c^{-1},1)\V_{E_i^\Delta}=\iota_1^{-1}(\V_{E_i^\Delta})\oplus\rho_2(\V_{E_i^\Delta}),
$$
where $\iota_j\: U_j\to U_1\oplus U_2$ is the natural injection and 
$\rho_j\: U_1\oplus U_2\to U_j$ is the natural surjection for $j=1,2$. Thus the $V^{(i)}$ satisfy Conditions~\eqref{exact2}, 
and hence $\mathfrak g(L,\V)$ is an exact 
level-$\Delta$ limit linear series on $X$. 

We have that $\mathfrak g(L,\V)$ depends only on the choice of isomorphisms $L(\mu^{(\delta)})_{E_i^{\delta}}\cong L^{(i)}$, whence its equivalence class is well-defined. Let $f$ be an automorphism of $T^\Delta$ inducing an equivalence of $\mu^\Delta$ with itself, and consider the 
induced twisted generalized linear series $(L,f^*\V)$. For each $i\in\Delta$, the point $f(E_i^\Delta)$ is on $T^\Delta_i$ and on the smooth locus of $T^\Delta$, with 
$f(E_i^\Delta)=E_i^\Delta$ if $i\in\mathbb Z$. For each $i\in\Delta$, let $\widetilde V^{(i)}$ be the image of 
$(f^*\V)_{E_i^\Delta}$ under the same isomorphism $L(\mu^\Delta)_{E_i^\Delta}\cong L^{(i)}$ we used to define $V^{(i)}$. If $i\in\mathbb Z$ then $\widetilde V^{(i)}=V^{(i)}$. If $i\not\in\mathbb Z$, at any rate, since $(L,\V)$ is everywhere invariant, there is $c\in\cc(k)$ such that 
$\V_{f(E_i^\Delta)}=(c^{-1},1)V_{E_i^\Delta}$. But $L^{(i)}$ is not invertible in this case, whence there is an automorphism $\psi_i$ of $L^{(i)}$ taking $V^{(i)}$ to $\widetilde V^{(i)}$. Thus $(L;\widetilde V^{(i)},i\in\Delta)$ is equivalent to $\mathfrak g(L,\V)$.

As for the third statement, given $\mathfrak g$, choose an isomorphism 
$\mathcal O^{(i)}_X\cong\F_{t_i}$ for each $i\in\Delta$, which induces an isomorphism $L^{(i)}\cong L\otimes\F_{t_i}$. As 
$\F_t\subseteq\mathcal O_Y\oplus\mathcal O_Z(dP)$ for each $t\in T$, we have that $L\otimes\F_t\subseteq L|_Y\oplus L|_Z(dP)$ for each $t\in T$. The image of $V^{(i)}$ under the induced isomorphism 
$L^{(i)}\cong L\otimes\F_{t_i}$ is thus a subspace of 
$U$ contained in 
$\Gamma(X,L\otimes\F_{t_i})$, which we will still denote by $V^{(i)}$. It follows that $(c^{-1},1)V^{(i)}$ is a subspace of 
$U$ contained in 
$\Gamma(X,L\otimes(c^{-1},1)\F_{t_i})=\Gamma(X,L\otimes\F_{c\ast t_i})$
for each $c\in\cc(k)$ and $i\in\Delta$. 

For each $i\in\Delta$, as $c$ varies in $\cc(k)$ the 
$(c^{-1},1)V^{(i)}$ form a locally free subsheaf of $U\otimes\mathcal O_{\cc}$. 
Identifying $(T^\Delta_i)^\ast$ with $\cc$, and extending, we obtain a locally free subsheaf $\V^{(i)}\subseteq U\otimes\mathcal O_{T^\Delta_i}$ such that 
$\V^{(i)}_{E^\Delta_i}=V^{(i)}$ for each $i\in\Delta$ and
$$
\V^{(i)}_{N^\Delta_i}=\lim_{c\to 0}(c^{-1},1)V^{(i)}\quad\text{and}
\V^{(i)}_{N^\Delta_j}=\lim_{c\to \infty}(c^{-1},1)V^{(i)}
$$
for each consecutive $i,j\in\Delta$. 
Since $(c^{-1},1)V^{(i)}\subseteq\Gamma(X,L\otimes(c^{-1},1)\F_{t_i})$ for each 
$c\in\cc(k)$, it follows that $\V^{(i)}$ is a subsheaf of 
$p_*L(\mu^\Delta)|_{T^\Delta_i}$ for each $i\in\Delta$.

The $\V^{(i)}$ will patch to a locally free subsheaf $\V\subseteq U\otimes\mathcal O_{T^\Delta}$ if and only if 
$\V^{(i)}|_{N_j^\Delta}=V^{(j)}_{N_j^\Delta}$ for each consecutive $i,j\in\Delta$. If so, since $\V^{(i)}\subseteq p_*L(\mu^\Delta)|_{T^\Delta_i}$ for each $i$, we have that $\V\subseteq p_*L(\mu^\Delta)$ and $(L(\mu^\Delta),\V)$ is a family of generalized linear series. Clearly, $(L,\V)$ is an everywhere invariant twisted generalized linear series along $\mu^\Delta$. Also, the image of $\V_{E_i^\Delta}$ under any isomorphism 
$L(\mu^\Delta)_{E_i^\Delta}\cong L^{(i)}$ is $V^{(i)}$ up to an automorphism of $L^{(i)}$ for each $i$, and hence 
$\mathfrak g(L,\V)$ is equivalent to $\mathfrak g$. 

It remains to prove that 
$\V^{(i)}|_{N_j^\Delta}=V^{(j)}_{N_j^\Delta}$ for consecutive $i,j\in\Delta$, or that
\begin{equation}\label{twolimit}
\lim_{c\to\infty}(c^{-1},1)V^{(i)}=\lim_{c\to 0}(c^{-1},1)V^{(j)}
\quad\text{for each consecutive }i,j\in\Delta
\end{equation}
as subspaces of $U_1\oplus U_2$. Again, Proposition~\ref{hV} yields
$$
\lim_{c\to 0}(c^{-1},1)V^{(i)}=
\rho_1(V^{(i)})\oplus\iota_2^{-1}(V^{(i)})
\text{\  \ and \  \ }
\lim_{c\to\infty}(c^{-1},1)V^{(j)}=
\iota_1^{-1}(\V^{(j)})\oplus\rho_2(V^{(j)}),
$$
hence \eqref{twolimit} follows from \eqref{exact2}, that is, from the exactness of $\mathfrak g$.

As for the last statement, let $\widetilde{\V}\subseteq p_*(L(\mu^\Delta))$ be a locally free subsheaf satisfying the same two conditions as $\V$. For each $i\in\Delta$, let $V^{(i)}$ (resp.~$\widetilde V^{(i)}$) be the image of $\V_{E_i^\Delta}$ (resp.~$\widetilde{\V}_{E_i^\Delta}$) under an isomorphism $L(\mu^\Delta)_{E_i^\Delta}\cong L^{(i)}$. Since 
$\mathfrak g(L,\V)$ and $\mathfrak g(L,\widetilde{\V})$ are equivalent to $\mathfrak g$, we have that $\widetilde V^{(i)}=V^{(i)}$ for $i\in\mathbb Z$, whence $\widetilde{\V}_{E_i^\Delta}=\V_{E_i^\Delta}$, and there is $c_i\in\cc(k)$ such that $\widetilde V^{(i)}=(c_i^{-1},1)V^{(i)}$ for each $i\in\Delta-\mathbb Z$, whence $\widetilde{\V}_{E_i^\Delta}=(c_i^{-1},1)\V_{E_i^\Delta}$. Since 
$(L,\widetilde\V)$ and $(L,\V)$ are everywhere invariant, it follows from Proposition˜\ref{invequiv} that $(L,\widetilde\V)$ and $(L,\V)$ are equivalent.
\end{proof}

\begin{rem} Each exact 
level-$\{0,\dots,d\}$ limit linear series gives rise to an exact
level-$\Delta$ limit linear series; see \cite{ENR}, Prop.~4.3, p.~??. 
The latter corresponds by Proposition~\ref{dlevel} to an everywhere twisted generalized linear series along $\mu^\Delta$, which is not everywhere constant if $\Delta\neq\{0,\dots,d\}$. Indeed, it is the pullback of a twisted generalized linear series along $1_T$, this one a continuous linear series.

A continuous linear series along a chain map $\mu\colon S\to T$ is equivalent to a continuous linear series along $\mu^{\Delta}$ for some $\Delta$, hence corresponds to a minimal exact level-$\Delta$ limit linear series.
\end{rem}

\section{Degenerations of linear series}\label{dls}

A {\it curve} is a reduced projective scheme over an algebraically closed field of pure dimension 1. A {\it family of curves}, which we denote by $\mathcal{X}/B$, is a morphism $\pi\colon\mathcal X\rightarrow B$ of schemes which is flat, projective and whose geometric fibers are curves. We say $\mathcal X/B$ is a \emph{smoothing} of its special fiber if $B=\text{Spec}(k[[t]])$ for an algebraically closed field $k$, the special fiber of $\mathcal{X}/B$ is a curve and its generic fiber is smooth. We say the smoothing is \emph{regular} if $\mathcal X$ is regular.

Let $\pi\colon\mathcal{X}\rightarrow B$ be a regular smoothing and $X$ its special fiber, a curve over a field $k$. Assume $X$ has only two irreducible components $Y$ and $Z$ and that they intersect transversally and at a unique point $P$.

Let $d$ be a nonnegative integer. Let $T$ be a {\it chain} of $d+1$ (smooth) rational curves over $k$. Thus $T$ has $d+1$ components, all isomorphic to $\mathbf P^1_k$, which can be ordered, $T_0,\dots,T_{d+1}$, such that $T_i\cap T_j\neq\emptyset$ if and only if
$|i-j|\leq 1$. Furthermore, $T_{i-1}$ and $T_{i}$ intersect transversally and at a unique point $N_i$ for each $i=1,\dots,d$. 

Let $\widehat T$ be the chain obtained from $T$ by
adding an extra rational curve at each end: one,
denoted $T_{-1}$, intersecting $T_0$ transversally, and at a unique point $N_0$, the other, denoted $T_{d+1}$, intersecting
$T_d$ transversally, and at a unique point $N_{d+1}$, which we will also denote by $N_\infty$. View $T\subset\widehat T$. For each $i=0,\dots,d$, let 
$E_i\in T_i-\{N_i,N_{i+1}\}$.

Let $\tau\colon\mathcal{T}\rightarrow B$ be a regular smoothing of $\widehat T$. It exists as the universal deformation space of a nodal curve is regular. Form the threefold
$\mathcal{X}\times_B\mathcal{T}$. It is flat and projective over $B$, with the surface $X\times\widehat T$ as special fiber. However,
it fails to be regular exactly at $(P,N_i)$
for $i=0,\dots, d+1$. Following the ideas in ~\cite[Subsection 4.2]{CEP}, 
blow up $\mathcal{X}\times_B\mathcal{T}$ successively
along the (strict transforms of)
$Y\times T_{-1}$, $Y\times T_0$, \dots, $Y\times T_d$ in this order. 
Let
$\psi\colon\widetilde{\mathcal X}\to\mathcal{X}\times_B\mathcal{T}$
denote the blowup map, and $\mathcal Y_i$ and $\mathcal{Z}_i$ the strict transforms in $\widetilde{\mathcal X}$ of $Y\times T_i$ and $Z\times T_i$ for each
$i=-1,\dots,d+1$. Then $\widetilde{\mathcal X}$
is regular and the $\mathcal Y_i$ and $\mathcal Z_i$ are Cartier divisors of
$\widetilde{\mathcal X}$. In the sense of ~\cite[Lemma 5.3, p.~2953]{CEP},
the map $\psi$ is a semistable modification of 
each of the two projections of $\mathcal{X}\times_B\mathcal{T}$ onto $\mathcal X$ and $\mathcal T$. In particular, the composition of the blowup with any of 
these projections is flat.

Let $R_i:=\psi^{-1}(P,N_i)$ for each $i=0,\dots,d+1$. Then the
$R_i$ are rational smooth curves. 
The strict transforms $\mathcal Y_{i-1}$ and $\mathcal Z_i$ contain $R_i$, 
while $\mathcal Y_{i}$ and $\mathcal Z_{i-1}$ intersect $R_i$
transversally at unique and distinct points. We have that 
$(\mathcal Y_j+\mathcal Z_j)\cdot R_i=0$ for each $j=-1,\dots,d+1$. 
Furthermore, $\mathcal Y_j\cdot R_i=0$ if $j\neq i-1,i$, whereas $\mathcal Y_i\cdot R_i=1$ and $\mathcal Y_{i-1}\cdot R_i=-1$.

Set
$$
\mathcal{G}:=\O_{\widetilde{\mathcal{X}}}
\Big(-\mathcal Y_{-1}+\mathcal Y_1+2\mathcal
Y_2+\cdots+(d+1)\mathcal Y_{d+1}\Big).
$$
Then $\mathcal{G}|_{R_i}\cong\O_{R_i}(1)$ for each
$i=0,\dots,d+1$ and hence $\mathcal G$ is $\psi$-admissible. It follows from ~\cite[Thm. 3.1, p.~63 ]{EP} that $\psi_{\ast}\mathcal{G}$ is a relatively torsion-free, rank-1 sheaf of degree 0 on
$\mathcal{X}\times_B\mathcal{T}/\mathcal{T}$, whose
formation commutes with base change. Put $\mathcal F:=\psi_{\ast}\mathcal{G}|_{X\times T}$. 

An invertible sheaf on the generic fiber of $\pi$ extends to an invertible sheaf on $\mathcal X$. Let $\mathcal L$ be an invertible sheaf on $\mathcal X$ whose restriction to the generic fiber has degree $d$. For each $j\in\mathbf Z$, the sheaves $\mathcal L$ and $\mathcal L(jY)$ have isomorphic restrictions to the generic fiber, but $\deg(\mathcal L(jY)|_Z)=\deg(\mathcal L|_Z)+j$. Assume that $\mathcal L|_Z$ has degree $0$, and thus $\mathcal L|_Y$ has degree $d$. Let $\mathcal L\boxtimes\mathcal G$ the tensor product with $\mathcal G$ of the pullback of $\mathcal M$ to $\widetilde{\mathcal{X}}$. Let $L$ be the restriction of $\mathcal L$ to the sepcial fiber $X$, and $L\boxtimes\mathcal F$ the tensor product with $\mathcal F$ of the pullback of $L$ to $X\times T$. Then 
$(L\boxtimes\mathcal F)|_{Y\times E_i}\cong L|_Y(-iP)$ and $(L\boxtimes\mathcal F)|_{Z\times E_i}\cong L|_Z(+iP)$ for $i=0,\dots,d$. 




A vector subspace of the space of sections of the restriction of $\mathcal L$ to the generic fiber is the restriction of a unique saturated free $k[[t]]$-submodule of $H^0(\mathcal X,\mathcal L)$. Let $\mathcal V\subseteq $


For each $i=0,\dots,d$, choose isomorphisms $T_i\cong\IP^1$ such that, denoting by 
$\xi_i\colon\O_{T_i}^2\to\O_{T_i}(1)$ the tautological quotient of $T_i$,
we have that $\xi_i|_{N_i}(\O_{N_i}\oplus 0)=0$ and
$\xi_i|_{N_{i+1}}(0\oplus\O_{N_{i+1}})=0$. Let $E_i\in T_i$ such that 
$\xi_i|_{E_i}$ is the cokernel of the diagonal $\mathcal O_{E_i}\hookrightarrow\O_{E_i}^2$. 
There is an action of $\cc$ on $T_i$ such that $N_i$ and $N_{i+1}$
are fixed points and $T_i^*:=T_i-\{N_i,N_{i+1}\}$ is a torsor for
$\cc$; choose the action such that $\xi_i|_{c\ast Q}=\xi_i|_Q(c,1)$,
where $(c,1)$ is the natural endomorphism on $\O_{Q}^2$, for each
$c\in\cc(k)$ and $Q\in T_i$. There is a unique action of $\cc$ on $T$ restricting to the actions on $T_i$ for each $i$. We call the action a \emph{chain action} on the two-pointed chain $(T,N_0,N_\infty)$. Notice that $T^*:=\bigcup T_i^*$ is the set of nonfixed points. Also, for each $i=0,\dots,d$ and each $Q\in T_i^*$, we have that 
$\lim_{c\to 0}c\ast Q=N_i$ whereas $\lim_{c\to\infty}c\ast Q=N_{i+1}$.


For a family of curves $\mathcal{X}/B$ a sheaf $\mathcal{G}$ on $\mathcal{X}/B$ is called {\it relatively torsion-free} (resp. {\it relatively rank} 1) on $\mathcal{X}/B$ if $\mathcal{G}_b$ is a torsion-free (resp. rank-1) sheaf on $\mathcal{X}_b$ for every geometric point $b$ of $B$. 

%%%%%%%%%%%%%%%%%%%%%%

It is now enough to show that the differentials $d_\infty h_1$ and $d_0h_2$ have different images in the tangent space $T_VG$ of $G$ at $V$, which satisfies
\[
T_VG=\mathrm{Hom}_k\Bigg(\rho_1(V_2)\oplus\rho_2(V_1),\frac{W_1}{\rho_1(V_2)}\oplus\frac{W_2}{\rho_2(V_1)}\Bigg).
\]
But for $i=1,2$, the map $h_i$ factors through the subscheme of $G$ parameterizing subspaces of $\rho_1(V_i)\oplus\rho_2(V_i)$. Thus the images of 
$d_\infty h_1$ and $d_0h_2$ are in the subspaces
\[
\text{Hom}_k\Bigg(\rho_1(V_2)\oplus\rho_2(V_1),\frac{\rho_1(V_1)}{\rho_1(V_2)}\oplus 0\Bigg)\quad\text{and}\quad \text{Hom}_k\Bigg(\rho_1(V_2)\oplus\rho_2(V_1),0\oplus\frac{\rho_2(V_2)}{\rho_2(V_1)}\Bigg)
\]
of $T_VG$, respectively. As they are nonzero, the two images are clearly different. Thus $f_{\mathfrak g}$ is an embedding.

It remains to show that $\mathcal S_{\mathfrak g}$ has multivariate Hilbert polynomial $\phi$. First of all, $\mathcal S_{\mathfrak g}$ is a curve, so its multivariate Hilbert polynomial $p_{\mathcal S_{\mathfrak g}}(u,v,s_0,\dots,s_d)$ is linear. Also, $\mathcal S_{\mathfrak g}$ is connected and of arithmetic genus zero, thus the constant term of $p_{\mathcal S_{\mathfrak g}}$ is $1$. Since 
$\mathcal S_{\mathfrak g}\subseteq H^r_d(X,L)$, the coefficient of $v$ in $p_{\mathcal S_{\mathfrak g}}$ is zero. Since $\mu_*[\mathcal S_{\mathfrak g}]=[T]$, the coefficient of $s_i$ is $1$ for each $i=0,\dots,d$. It remains to show that the coefficient of $u$ is $r+1$, which is equivalent to showing that $f(\mathcal S)$ has degree $r+1$ in $G$.

Let $U_j:=\V_{Q_j}\subseteq W$ for $j=0,1,\dots,m+1$. It follows from Proposition~\ref{hV} that the degree of $f(\mathcal S)$ is 
$$
\sum_{i=1}^{m+1}\Big(\dim_k\rho_1(U_{i-1})-\dim_k\rho_1(U_{i})\Big),
$$
which is equal to $\dim_k\rho_1(U_0)-\dim_k\rho_1(U_{m+1})$.
But 
$$
U_0\subseteq\Gamma(Y,L|_Y)\oplus\Gamma(Z,L|_Z(-P))\quad\text{and}\quad
U_{m+1}\subseteq\Gamma(Y,L|_Y(-(d+1)P))\oplus\Gamma(Z,L|_Z(dP)).
$$
Since $L|_Z(-P)$ and $L|_Y(-(d+1)P)$ have negative degrees, $\dim_k\rho_1(U_0)=r+1$ and $\rho_1(U_{m+1})=0$. 
Thus the degree of $f(\mathcal S)$ is $r+1$. 
%%%%%%%%%%%%%%%%

%%%%%%%%%%%%%%%
\section{Flags of continuous linear series}

\subsection{The moduli space}

\begin{defn} Let $\underline r:r_0,\dots,r_m$ be an increasing sequence of nonnegative integers. A \emph{family of twisted flags of generalized linear series of rank $\underline r$ and degree $d$ over $B$} consists of the data of families $(\L,\V_i)$ of twisted generalized linear series on $X$ of degree $d$ and rank $r_i$ along families of chains maps $\gamma_i\:\mathcal S_i\to B\times T$ for $i=0,\dots,m$ for which there are a family of chain maps $\gamma\:\mathcal S\to B\times T$ and maps $f_i\:\mathcal S\to \mathcal S_i$ of families of chain maps such that $f_0^*\V_0\subseteq f_1^*\V_1\subseteq\cdots\subseteq f_m^*\V_m$. We say that the data is a \emph{twisted flag of generalized linear series} if $B=\text{Spec(k)}$.  {\color{red} Mudar $m$ pra $p$.}
\end{defn} 

%With the notation as above, notice that $(1_X\times f)^*\L(\gamma)=\L(\gamma f)$. Since $(\L(\gamma),\V)$ is a family of generalized linear series, so is its pullback under $f$, which is $(\L(\gamma f),f^*\V)$, and thus so is  $(\L(\gamma f),\V')$. In particular, $(\L,\V')$ is a family of twisted generalized linear series on $X$ along $\gamma f$.

\begin{defn} A twisted flag of generalized linear series on $X$ is said to be
\emph{everywhere nonconstant} (resp.~\emph{everywhere invariant}) if all the twisted generalized linear series that compose it are everywhere nonconstant (resp.~everywhere invariant). It is said to be a \emph{continuous flag of linear series} if it is everywhere nonconstant and invariant. 
\end{defn}

\begin{prop}\label{stabconstflag2} For each $i=1,\dots,m$, let $(\L,\V_i)$ be a family of continuous linear series of degree $d$ and rank $r_i$ along a 
family of chain maps $\gamma_i\:\mathcal S_i\to B\times T$. Assume that $\mathfrak g:=(\L,\V_0,\dots,\V_m)$ is a family of continuous flags of linear series. Then there are a family of chain maps $\gamma\:\mathcal S\to B\times T$ and maps $f_i\:\mathcal S\to \mathcal S_i$ of families of chain maps such that: 
\begin{enumerate}
    \item $f_0^*\V_0\subseteq f_1^*\V_1\subseteq\cdots\subseteq f_m^*\V_m$;
    \item for each $b\in B(k)$ and each irreducible component $R$ of $\mathcal S_b$, there is $i$ such that $(\L_b,f_i^*\V_i|_{\mathcal S_b})$ is nonconstant along $R$.
\end{enumerate}
\end{prop}

\begin{proof} As in Proposition~\ref{stabconst}, let $\mathcal S$ be as in the definition of a flag, and consider the stabilization of $\mathcal S/B$ with respect to $(f_1,\dots,f_m)\:\mathcal S\to \prod\mathcal S_i$, the product being fibered over $B\times T$. We may replace $\mathcal S/B$ by the stabilization, thus assuming that $(f_1,\dots,f_m)$ is stable {\color{red} como um mapa. Ainda. é necessário argumentar (1) para a estabilização}. But thus, for each $b\in B(k)$ and each irreducible component $R$ of $\mathcal S_b$, there is $i$ such that $f_i(R)$ is not a point. Since $(\L_b,\V_i|_{\mathcal S_{i,b}})$ is everywhere nonconstant, it follows that $(\L_b,f_i^*\V_i|_{\mathcal S_b})$ is nonconstant along $R$.
\end{proof}

\begin{defn} Two families of twisted flags of generalized linear series of the same rank $\underline r$ are \emph{equivalent} if the corresponding families of twisted generalized linear series of rank $r_i$ are equivalent. 
\end{defn}

\begin{defn} The (contravariant) functor of continuous flags of linear series of 
rank $\underline r$ and degree $d$ on $X$,
$$
\mathfrak{G}_d^{\underline r}(X)\colon (\text{Schemes})
\longrightarrow(\text{Sets}),
$$
is the functor that associates to each scheme $B$ the
set of equivalence classes of families of 
continuous flags of linear series of degree $d$ and rank $\underline r$ over $B$. 
\end{defn}

It is clear that $\mathfrak{G}_d^{\underline r}(X)$ is a subfunctor of the 
product of the $\mathfrak{G}_d^{r_i}(X)$ over $J$. Now, $\mathfrak{G}_d^{r_i}(X)$ is isomorphic to the functor of points of an open and closed subscheme of $\text{\rm Hilb}_{H_d^{r_i}(X)}^{\phi_i,\cc}$ for $i=0,\dots,p$, where 
$\phi_i:=(r_i+1)u_i+\sum_{j=0}^d s_j+1$. 

Let 
$$
H^{\underline r}_d(X)\subseteq H^{r_0}_d(X)\times_{(J\times T)} H^{r_1}_d(X)\times_{(J\times T)}\cdots\times_{(J\times T)}H^{r_p}_d(X)
$$
be the closed subscheme parameterizing tuples $(L,t,V_0,V_1,\dots,V_p)$ with $V_0\subseteq V_1\subseteq\cdots \subseteq V_p\subseteq\Gamma(X,L\otimes\mathcal F_t)$ for $L\in J$ and $t\in T$, where each $V_i$ has dimension $r_i+1$. The product of the actions of $\cc$ on the $H^{r_i}_d(X)$ leaves $H^{\underline r}_d(X)$ invariant. 

The action on $H^{\underline r}_d(X)$ induces an action on the Hilbert scheme of $H^{\underline r}_d(X)$. We will denote by $\text{Hilb}^{\underline\phi}_{H^{\underline r}_d(X)}$ the open (and closed) subscheme of the Hilbert scheme parametrizing subschemes $S$ of $H^{\underline r}_d(X)$ with multivariate Hilbert polynomial with respect to the product of the $H^{r_i}_d(X)$ satisfying
$$
h^0(S,\mathcal O_S(u_0,u_1,\dots,u_p,v,s_0,\dots,s_d))=\underline\phi(u,v,s_0,\dots,s_d):=\sum_{i=0}^pu_i(r_i+1)+\sum_{i=0}^d s_i+1.
$$
We will also denote by $\text{Hilb}^{\underline\phi,\cc}_{H^{\underline r}_d(X)}$ the closed subscheme of 
$\text{Hilb}^{\underline\phi}_{H^{\underline r}_d(X)}$ parameterizing invariant subschemes of $H^{\underline r}_d(X)$.

Keep the setup as in Proposition~\ref{stabconstflag2}. There is a map $f_{\mathfrak g}\:\mathcal S\to H^{\underline r}_d(X)$ induced by the filtration 
$f_0^*\V_0\subseteq f_1^*\V_1\subseteq\cdots\subseteq f_m^*\V_m$. 

\begin{prop}\label{betaclassflag} If $\mathfrak g$ is a family of continuous flags of linear series then $f_{\mathfrak g}$ is a closed embedding. Furthermore, its image $\mathcal S_{\mathfrak g}$ is $\cc$-invariant and has relative multivariate Hilbert polynomial $\underline\phi$. 
\end{prop}

\begin{proof} As in the proof of Proposition~\ref{betaclass} we may assume that $B=\text{Spec}(k)$. Let $\mathfrak g_i:=(\L,\V_i)$ for each $i=0,\dots,p$. Then $f_{\mathfrak g}=(f_{\mathfrak g_1}f_1,\dots,f_{\mathfrak g_p}f_p)$ as maps to the product 
$\prod H^{r_i}_d(X)$. Since for each irreducible component $R$ of $\mathcal S$ there is $i$ such that $f_i^*\mathfrak g_i$ is not constant along $R$, it follows as in the proof of Proposition~\ref{betaclass} that $f_{\mathfrak g_i}f_i|_R$ and hence $f_{\mathfrak g}|_R$ are embeddings. 

{\color{red} Argumentar a injetividade}

To show $f_{\mathfrak g}$ is an embedding, we need only show that the differentials $d_Nf_{\mathfrak g}|_{C_1}$ and $d_Nf_{\mathfrak g}|_{C_2}$ have different images for each node $N$ of $S$, where $C_1$ and $C_2$ are the irreducible components of $\mathcal S$ containing $N$. If there is $i$ such that $f_i^*\mathfrak g_i$ is nonconstant along $C_1$ and $C_2$ then it follows as in the proof of Proposition~\ref{betaclass} that the images are indeed different. If not, there are $i,j$ such that $f_i^*\mathfrak g_i$ is nonconstant along $C_1$ but constant along $C_2$ and $f_j^*\mathfrak g_j$ is constant along $C_1$ but nonconstant along $C_2$. Then 
$\mu(C_1)=\mu(N)=\mu(C_2)$. And since $f_{\mathfrak g_i}f_i$ restricts to an embedding on $C_1$ but is constant on $C_2$, and $f_{\mathfrak g_j}f_j$ restricts to an embedding on $C_2$ but is constant on $C_1$, the images of $d_Nf_{\mathfrak g}|_{C_1}$ and $d_Nf_{\mathfrak g}|_{C_2}$ must be different.

Since the $\mathfrak g_i$ are everywhere invariant, the image $\mathcal S_{\mathfrak g}$ is invariant in the product $\prod H^{r_i}_d(X)$, thus in $H^{\underline r}_d(X)$. The proof that $\mathcal S_{\mathfrak g}$ has relative multivariate Hilbert polynomial $\underline\phi$ is as in Proposition~\ref{betaclass}.
\end{proof}

\begin{thm}\label{chainmapsflag} For each connected 
$S\in \text{\rm Hilb}^{\underline\phi,\cc}_{H^{\underline r}_d(X)}(k)$, the induced map $S\to T$ is a chain map.
\end{thm}

\begin{proof} The proof is essentially that of Theorem~\ref{chainmaps}, we highlight the difference. First of all, since the coefficient of $v$ in $\underline\phi$ is zero, we have that 
\begin{equation}\label{Sprod}
S_\red\subseteq H^{r_0}_{d}(X,L)\times\cdots \times H^{r_p}_{d}(X,L)\subseteq 
\text{Grass}(r_0+1,W)\times\cdots\times\text{Grass}(r_p+1,W)\times T
\end{equation}
equivariantly for some invertible sheaf $L$ on $X$ of bidegree $(d,0)$, where $W$ is as in the proof of Theorem~\ref{chainmaps}. For each $\ell=0,\dots,p$, let $q_\ell$ denote the projection of the righ-hand side product in \eqref{Sprod} to $G_\ell:=\text{Grass}(r_\ell+1,W)$. Recall that, the space $H^{r_i}_d(X,L)$ denotes the scheme of generalized linear series of degree $d$ and rank $r_i$ with sections in $L$.

The first and second claims and their proofs apply verbatim. The third claim holds modified: for each $\ell=0,\dots,p$, the $q_\ell(C_i)$ are orbit closures in $G_\ell$ and the sum of their degrees is at most $r_\ell+1$, with equality only if $q_\ell(C_i)\cap q_\ell(C_j)$ is set-theoretically at most a point for distinct $i,j$. The proof is the same. The fourth claim holds modified as well: each $C_i$ is rational; for each $\ell=0,\dots,p$ the map $C_i\to q_\ell(C_i)$ is either an isomorphism or constant and the sum of the degrees of the $q_\ell(C_i)$ in $G_\ell$ is $r_\ell+1$; if the maps $C_i\to q_\ell(C_i)$ are constant for all $\ell$ then $i=i_v$ for some $v$. The proof is the same. 

The fifth claim holds verbatim. Its proof must be modified only in the last paragraph: assuming that $i>i_{\ell-1}$, we have that $q_\ell(C_i)$ is a subcurve of $G_\ell$ for some $\ell$. Since 
$q_\ell(C_i)\cap q_\ell(C_j)$ is a point and $C_i\to q_\ell(C_i)$ is an isomorphism by the third and fourth claims, we must have that $C_i\cap C_j$ is set-theoretically a point.

The remaining of the proof is the same.  

{\color{red} Verbatim: Checar que nada precisa ser modificado.} 
\end{proof}

\begin{thm}\label{moduliflag} The functor $\mathfrak{G}_d^{\underline r}(X)$ is a closed subfunctor of the product of the $\mathfrak{G}_d^{r_i}(X)$ over $J$. In particular, there is a projective scheme representing $\mathfrak{G}_d^{\underline r}(X)$.
\end{thm}

\begin{proof} A $B$-point of the fibered product 
$\prod\mathfrak{G}_d^{r_i}(X)$ over $J\times T$ corresponds to a collection of $B$-flat subschemes $\mathcal S_i\subseteq H^{r_i}_d(X)\times B$ whose fibers over $B$ are connected. If it is a $B$-point of $\mathfrak{G}_d^{\underline r}(X)$ then there is a subscheme $\mathcal S\subseteq H^{\underline r}_d(X)\times B$ such that $\mathcal S\subseteq \prod\mathcal S_i$, the product being fibered over $B\times T$. Conversely, the condition that the fibered product $\prod\mathcal S_i$ over $T\times B$ contain a 
$B$-flat subscheme $\mathcal S\subseteq H^{\underline r}_d(X)\times B$ whose fibers over $B$ are connected is closed in $B$, and implies that the $B$-point is a $B$-point of $\mathfrak{G}_d^{\underline r}(X)$. 
{\color{red} A ser precisada melhor.} 
\end{proof}

\subsection{Forgetful and Abel maps}

Let $\underline r$ be a nondecreasing sequence of nonnegative integers. Let $G^{\underline r}_d(X)$ be a projective scheme representing the functor $\mathfrak G^{\underline r}_d(X)$. If $\underline s$ is a subsequence of $\underline r$, there is a natural forgetful morphism $\Phi^{\underline r}_{\underline s}\:G^{\underline r}_d(X)\to G^{\underline s}_d(X)$.    

For each nonnegative integer $d$, let $\text{Hilb}_X^d$ be the corresponding component of the Hilbert scheme of points of $X$. In  
\cite{E-Oss} is described an Abel map $S^d(X)\to J$, where $S^d(X)$ is the $d$-th symmetric product of $X$. Compose it with the Hilbert-Chow map to get an Abel map $A_d\:\text{Hilb}_X^d\to J$. 
%The map $A_d$ takes a subscheme $W$ of $X$ of length $d$ to the invertible sheaf $L$ of degree $(d,0)$ on $X$ whose dual $L^*$ is equivalent to the sheaf of ideals of $W$ in the Grothendieck group of $X$.  

\begin{thm} Let $d$ be a nonnegative integer. Then the map $G^0_{d,1}(X)\to J$ factors through a natural isomorphism $G^0_{d,1}(X)\to\text{\rm Hilb}^d_X$.    
\end{thm}

\begin{proof} Let $D$ be an effective divisor on $X$. The natural maps $\mathcal O_X^{(i)}\to\mathcal O_X^{(i\pm 1)}$ induce a linked chain of vector spaces
$$
\mathfrak V_{L,D}:\Gamma(X,L^{(0)}(D))\longleftrightarrow \Gamma(X,L^{(1)}(D))))\longleftrightarrow \dots\longleftrightarrow \Gamma(X,L^{(d)}(D)).
$$
Assume $D$ is sufficiently ample. Then the $\Gamma(X,L^{(i)}(D))$
are vector spaces of dimension $d+\deg D+1-g$ for each $i=0,\dots,d$ and each $L\in J$, where $g$ is the arithmetic genus of $X$, and 
$\mathfrak V_{L,D}$ is exact. Thus we obtain a $J$-scheme $\mathbb L\mathbb P_{J,D}$ whose fiber over $L\in J$ is the scheme $\mathbb L\mathbb P(\mathfrak V_{L,D})$ of subchains of vector spaces of dimension $1$ of $\mathfrak V_{L,D}$. The fiber $\mathbb L\mathbb P(\mathfrak V_{L,D})$ is a reduced subscheme of the product 
$\prod\mathbb P(\Gamma(X,L^{(i)}(D))$ with multivariate Hilbert polynomial equal to that of diagonal; see \cite[Thm. 4.3, p. 84]{E-Oss}. Then $\mathbb L\mathbb P_{J,D}$ is $J$-flat and reduced. Now, there is a natural map 
$\mathbb L\mathbb P_{J,D}\to\text{Hilb}^{d+\deg D}_X$, given by taking the intersection of the zero loci of the sections of $L^{(i)}(D)$ generating each subchain of dimension 1 of $\mathfrak V_{L,D}$ for each $L\in J$. Indeed, by \cite[\S 5-6]{ESVII}, the intersection is a subscheme of $X$ of length $d$ whose sheaf of ideals is isomorphic to $L^{(i)}(D)^*$ for some $i\in 2^{-1}\mathbf Z$. Since $\mathbb L\mathbb P_{J,D}$ is reduced, the association is indeed a morphism.

Now, there is clearly an embedding 
$G^0_{d,1}\to\mathbb L\mathbb P_{J,D}$. The composition 
$G^0_{d,1}\to\mathbb L\mathbb P_{J,D}\to\text{Hilb}^{d+\deg D}_X$ maps to the subscheme parameterizing subschemes of $X$ containing $D$, and thus factors through $\text{Hilb}^d_X$.  Furthermore, the sheaf of ideals of a subscheme in the image of the fiber of $G^0_{d,1}$ over each $L\in J$ is isomorphic to $(L^{(i)})^*$ for some $i\in 2^{-1}\mathbf Z$, whence the subscheme lies on the fiber over $L$ of the Abel map $A\:\text{Hilb}^d_X\to J$. Thus 
$G^0_{d,1}\to\text{Hilb}^{d}_X$ is a $J$-map. 

We claim that $G^0_{d,1}(X)\to\text{Hilb}^d_X$ is a closed embedding. Indeed, it is enough to show that $\mathbb L\mathbb P_{J,D}\to \text{Hilb}^{d+\deg D}_X$ is an embedding, which is reduced to showing that $\mathbb L\mathbb P(\mathfrak V_{L,D})\to\text{Hilb}^{d+\deg D}_X$ is an embedding for each $L\in J$. First, the map is injective. Indeed each subchain of $\mathfrak V_{L,D}$ of one-dimensional subspaces $W^{(j)}\subseteq\Gamma(X,L^{(j)}(D))$ is generated by two consecutive, 
$W^{(i)}$ and $W^{(i+1)}$, the corresponding subscheme of $X$ being the intersection of the zero schemes of generators $s_i$ and $s_{i+1}$ of them, with sheaf of ideals thus isomorphic to $(L^{(i)}(D))^*$, 
$(L^{(i+1)}(D))^*$ or  $(L^{(i+1/2)}(D))^*$, according to whether $W^{(i)}$ only or $W^{(i+1)}$ only or both generate the subchain. As no two sheaves 
$(L^{(j)}(D))^*$ for $j\in 2^{-1}\mathbf Z$ are isomorphic, knowledge of the subscheme of $X$ recovers the subchain.

Second, the differential of the map is injective at each point. Indeed, consider a subchain of $\mathfrak V_{L,D}$ of one-dimensional subspaces $W^{(j)}\subseteq\Gamma(X,L^{(j)}(D))$ generated by $W^{(i)}$ and $W^{(i+1)}$. 
Sections $s_i$ and $s_{i+1}$ generating $W^{(i)}$ and $W^{(i+1)}$ correspond to a map $(s_i^*,s_{i+1}^*)\:(L^{(i)}(D))^*\oplus (L^{(i+1)}(D))^*\to\mathcal O_X^2$ whose  
composition with the sum has image $\mathcal I$, where $\mathcal I$ is the sheaf of ideals of the corresponding subscheme of $X$. Were the subchain  generated by $W^{(i)}$ only or $W^{(i+1)}$ only then $s_i^*$ or 
$s_{i+1}^*$ only would have image $\mathcal I$, whence their infinitesimal variation would vary $\mathcal I$ infinitesimally as well. Otherwise, an infitesimal variation of the subchain corresponds to a 
map $(L^{(i)}(D))^*\oplus (L^{(i+1)}(D))^*\to\mathcal O_X^2$. Assuming it does not vary the subscheme of $X$, its composition with the sum has image $\mathcal I$. Since $\mathcal I$ is not invertible, we must have that the induced map $(L^{(i)}(D))^*\oplus (L^{(i+1)}(D))^*\to\mathcal O_X^2$ factors through an isomorphism  $(L^{(i)}(D))^*|_Y\oplus (L^{(i+1)}(D))^*|_Z\to\mathcal I$, and hence the infinitesimal variation of the subchain is trivial.

{\color{red} A explicar melhor ainda.}

Furthermore, $G^0_{d,1}(X)\to\text{Hilb}^d_X$ is set-theoretically surjective. Indeed, given a subscheme $W\subseteq X$ of length $d$, let $L\in J$ be the invertible sheaf corresponding to it by the Abel map. Then $W$ is the zero locus of a section of $L^{(i)}$ for some $i\in 2^{-1}\mathbf Z$. The section generates a subchain of dimension 1 of $\mathfrak V_{L,D}$, which corresponds to a point on $G^0_{d,1}(X)$. 

So $G^0_{d,1}(X)\to\text{Hilb}^d_X$ is an isomorphism of $G^0_{d,1}(X)$ onto a closed subscheme of $\text{Hilb}^d_X$ supported on the same set, but $\text{Hilb}^d_X$ is reduced, so that subscheme must be $\text{Hilb}^d_X$ itself.
\end{proof}
 Recall the definiton of $J$-map $\Phi\: G^0_d(X)\to G^0_{d,1}(X)$ in Proposition \ref{toOss}.
%{\color{red} Recordar $\Phi$.}

\begin{prop}\label{toOss0} The natural $J$-map $\Phi\: G^0_d(X)\to G^0_{d,1}(X)$ is an isomorphism.
\end{prop}

\begin{proof} A continuous limit linear series of degree 0 is the data $(L,\V)$ of a twisted generalized linear series along a chain map $\mu\:S\to T$ where $S$ has at most on irreducible component more than $T$. 

If $S$ has the same number of components as $T$, then 
$(L,\V)$ is determined by its image in $G^0_{d,1}$. If not, assume $S$ has a extra component over the node $N_i$. Let $N_{i,0}$ (resp.~$N_{i,\infty}$ be the intersection of this component with that above $T_i$ (resp.~$T_{i+1}$). Write $\V_{N_{i,0}}=V_{0,Y}\oplus V_{0,Z}$ and $\V_{N_{i,\infty}}=V_{\infty,Y}\oplus V_{\infty,Z}$. Then $V_{\infty,Y}=0$ and $V_{0,Z}=0$. 
Furthermore, the extra component of $S$ corresponds to the orbit of a subspace $W\subseteq\Gamma(X,L^{(i+1/2)})$ which is contained in 
$V_{0,Y}\oplus V_{\infty,Z}$. Since the latter has dimension 2 and the orbit has dimension 1, the orbit consists of all one-dimensional subspaces of $V_{0,Y}\oplus V_{\infty,Z}$, being thus determined by the image of $(L,\V)$ in $G^0_{d,1}$. It follows that $\Phi$ is injective.

On the other hand, given a subchain of on-edimensional vector spaces $V^{(i)}$ of $\Gamma(X,L^{(i)})$, there are two cases. If the subchain is exact, it corresponds to a point on $G^0_d(X)$ under the map $G^{0,*}_d(X)\to G^0_d(X)$. If the chain is not exact, then it is the image of a point on $G^{0,*}_{d,2}(X)$, which corresponds to a point on $G^0_d(X)$ under the map $G^{0,*}_{d,2}(X)\to G^0_d(X)$. 




{\color{blue} Sketch: To show it is a closed embedding, we restrict to the fiber over a $L\in J$, and then it should be a linear matter. Finally, use that $G^0_{d,1}(X)$ is reduced, because equal to $\text{Hilb}^d_X$.}
\end{proof}

\begin{thm} The map $\Phi^{0,r}_r\:G^{0,r}_d(X)_{\mathrm{red}}\to G^r_d(X)_{\mathrm{red}}$ is flat and the embedding  
$$
\Phi^{0,r}_0\times\Phi^{0,r}_r\:G^{0,r}_d(X)\longrightarrow G^0_d(X)\times_J G^r_d(X)
$$
induces an injective map $G^r_d(X)_{\mathrm{red}}\to\text{\rm Hilb}_{\text{\rm Hilb}^d_X}$.
\end{thm}

\begin{proof} {\color{red} A colocar. Pode estar errado.}  {\color{blue} Sketch: An element $(L,\V)$ of 
$G^r_d(X)$ corresponds to a level-$\Delta$ limit linear series $(L;V^{(i)},i\in\Delta)$ for a 
finite subset $\Delta\subseteq [0,d]\cap\mathbb Q$ containing $\{0,1,2,\dots,d\}$. Using the isomorphism 
$G^0_d(X)\to G^0_{d,\{0,\dots,d\}}(X)$ we can identify the fiber of $\Phi^{0,r}_r$ over $(L,\V)$ with the image in the linked Grassmannian $\mathbb L\mathbb P(V^{(i)},i\in\{0,\dots,d\})$ of the linked Grassmnannian 
$\mathbb L\mathbb P(V^{(i)},i\in\Delta)$. As $(V^{(i)},i\in\Delta)$ is an exact chain of vector spaces, $\mathbb L\mathbb P(V^{(i)},i\in\Delta)$ is a degeneration of the diagonal by ??, and thus so is its image in $\mathbb L\mathbb P(V^{(i)},i\in\{0,\dots,d\})$. It follows that all the fibers of $\Phi^{0,r}_r$ have the same Hilbert polynomial. Flatness would ensue if we knew $G^r_d(X)$ were reduced. So far, we cannot conclude this.}
\textcolor{red}{Pra pensar com um pouco mais de cuidado: poderiamos impor a estrutura de esquema reduzido induzido sobre $G^r_d(X)$ e dai, sob essa condicao, garantir que a construçao aqui é possivel}. 
\end{proof}
%%%%%%%%%%%%%%%%

%%%%%%%%%%%%%%%
\begin{prop}\label{stabconst} An equivalence class of families of everywhere invariant twisted generalized linear series along a family of chain maps 
is the pullback of a unique equivalence class of families of continuous linear series along an intermediate family of chain maps. 
\end{prop}

{\color{red} Será que não é melhor colocar mais pra frente na secao 6?}

\begin{proof} Let $\mathfrak g=(\L,\V)$ be a family of everywhere invariant twisted generalized linear series along a family of chain maps $\gamma\:\mathcal S\to B\times T$.  Let $f_\mathfrak g\:\mathcal S\to H^r_d(X)\times B$ be the induced map. Let $\pi'\:\mathcal S'\to B$ be the stabilization of $f_{\mathfrak g}$ (cf. \cite[\S 2]{Kn} and \cite[proof of Prop. 6]{FPnotes}). More precisely, 
$$
\mathcal S':=\text{Proj}_B\Big(\bigoplus_{n\geq 0}\gamma_{1,*}(\mathcal A^{\otimes n})\Big),
$$
where $\mathcal A$ is the pullback under $f_{\mathfrak g}$ of a relatively ample invertible sheaf on $H^r_d(X)\times B/B$. There are a $B$-map $f'\:\mathcal S\to\mathcal S'$ collapsing the same irreducible components in the fibers of $\gamma_1$ collapsed under $f_{\mathfrak g}$ and a $B$-map $f'_{\mathfrak g}\:\mathcal S'\to H^r_d(X)\times B$ such that $f_{\mathfrak g}=f'_{\mathfrak g}f'$. Let $\gamma'$ be the composition of $f'_{\mathfrak g}$ with the natural map $H^r_d(X)\times B\to T\times B$. Then $\gamma'\:\mathcal S'\to B\times T$ is a family of chains maps satisfying $\gamma=\gamma'f'$. Furthermore, the map yields a family of twisted generalized linear series $\mathfrak g':=(\L,\V')$ whose pullback under $f$ is $(\L,\V)$ and such that $f_{\mathfrak g'}=f'_{\mathfrak g}$. 


Clearly, $(\L,\V')$ is a family of everywhere invariant series. 
Since $f_{\mathfrak g'}$ does not collapse any 
irreducible component of any fiber of $\pi'$, the family $(\L,\V')$ is thus one of continuous linear series. Then $f_{\mathfrak g'}$ is an embedding by Proposition~\ref{betaclass}. If there were a family $(\L,\V')$ of continuous linear series along a family of chain maps $\gamma''\:\mathcal S''\to B\times T$ and a morphism of families of chain maps $f''\:\mathcal S\to\mathcal S''$ such that $(\L,\V)$ is the pullback of $(\L,\V'')$ under $f''$, then $f_{\mathfrak g''}$ is an embedding by Proposition~\ref{betaclass}, and $f_{\mathfrak g}=f_{\mathfrak g''}f''$. There is thus an isomorphism $f\:\mathcal S'\to\mathcal S''$ such that $f''=ff'$ and $f_{\mathfrak g'}=f_{\mathfrak g''}f$. It follows that $(\L,\V'')$ is equivalent to $(\L,\V')$. 

Finally, a family equivalent to $(\L,\V)$ is clearly the pullback of a family equivalent to $(\L,\V')$.
\end{proof}
%%%%%%%%%%%%%%%

%%%%%%%%%%%%%%%

\subsection{Notation} We will fix notation along the text. Here we make a short list of the important ones:
\begin{itemize}
    \item $k$, a given algebraically closed field;
    \item $X$, a given nodal curve, $Y$ and $Z$, all of its components, $P=Y\cap Z$, its only singularity;
    \item $d$, a given nonnegative integer;
    \item $(T,N_0,N_{\infty})$, a chain, $T_0,\dots,T_d$, its ordered components, $N_i:=T_{i-1}\cap T_i$ for $i=1,\dots,d$;
    \item $T_i\xrightarrow{\cong}\mathbf P^1$, chosen isomorphism, taking $N_i$ to $(0:1)$ and $N_{i+1}$ to $(1:0)$, for $i=0,\dots,d$;
    \item $E_i\in T_i(k)$, point taken to $(-1:1)$;
    \item $\O_Y(P)|_P\xrightarrow{\cong}\O_P$ and $\O_Z(P)|_P\xrightarrow{\cong}\O_P$, chosen isomorphisms;
    \item $p_S\:X\times S\to S$, indicated projection;
    \item $J:=\mathrm{Pic}^{(d,0)}_X$, component of Picard scheme of $X$, and $\mathcal P$ Poincaré sheaf on $X\times J$;
    \item $H^r_d(X)$, scheme over $J\times T$ of generalized linear series;
    \item $H^r_d(X,L)$, scheme over $T$ of generalized linear series with sections in $L$;
    \item $\phi(u,v,s_0,\dots,s_d):=u(r+1)+\sum_{i=0}^d s_i+1$, Hilbert polynomial on $H^r_d(X)$.
\end{itemize}

%%%%%%%%%%%%%%%%%%%%%

\begin{prop}\label{invequiv} Let $L$ be an invertible sheaf on $X$ of bidegree $(d,0)$ and $\gamma\:S\to T$ a chain map. Fix a chain action of $\cc$ on $S$ lifting the chain action on $T$. Let $(L,\V)$ and $(L,\V')$ be everywhere invariant twisted generalized linear series along $\gamma$. Then 
$(L,\V)$ and $(L,\V')$ are equivalent if and only if for each 
irreducible component $C$ of $S$ there are $s_1,s_2\in C^*(k)$ and $x\in\cc(k)$ such that $\gamma(s_2)=x\ast\gamma(s_1)$ and $\V'_{s_2}=(x^{-1},1)\V_{s_1}$.
\end{prop}

\begin{proof} Assume $(L,\V)$ and $(L,\V')$ are equivalent. Then there is an automorphism $f$ of $\gamma$ such that $\V'=f^*\V$. For each irreducible component $C$ of $S$, let $s_1\in C^*(k)$ be any $k$-point, and $s_2:=f(s_1)$. Then $s_2\in C^*(k)$ and 
$\V'_{s_1}=\V_{s_2}$. Since $f$ is an automorphism  of $\gamma$, we have that 
$\gamma(s_2)=\gamma(s_1)=1\ast\gamma(s_1)$. 

Conversely, assume that for each irreducible component $C$ of $S$ there are $s^C_1,s^C_2\in C^*(k)$ and $x_C\in\cc(k)$ such that $\gamma(s^C_2)=x_C\ast\gamma(s^C_1)$ and $\V'_{s^C_2}=((x_C)^{-1},1)\V_{s^C_1}$. Since $(L,\V)$ and $(L,\V')$ are everywhere invariant, there are isomorphisms $\theta_C,\theta_C'\:\cc\to C^*$ such that $\gamma\theta_C$ and $\gamma\theta'_C$ are equivariant and $\V_{\theta_C(x)}=(x^{-1},1)\V_{\theta_C(1)}$ and $\V'_{\theta'_C(x)}=(x^{-1},1)\V'_{\theta'_C(1)}$ for each $x\in\cc(k)$. Let $x^C_1,x^C_2\in\cc(k)$ such that 
$\theta_C(x_1^C)=s^C_1$ and $\theta'_C(x_2^C)=s^C_2$. Let $f$ be the unique automorphism of $S$ fixing the nodes and the points over $N_0$ and $N_\infty$ but letting $f(\theta'_R(x))=\theta_R((x^C_2)^{-1}x_Cx^C_1x)$ for each $C$ and 
$x\in\cc(k)$. It is an automorphism of $\gamma$, because 
$$
\gamma f(\theta'_C(x))=
\gamma\theta_C((x^C_2)^{-1}x_Cx^C_1x)=(x_C(x^C_2)^{-1}x)\ast\gamma(s^C_1)=
((x^C_2)^{-1}x)\ast\gamma(s^C_2)=\gamma(\theta'_C(x))
$$
for each $C$ and $x\in\cc(k)$. Furthermore,
$$
\V'_{\theta'_C(x)}=(x^{-1}x^C_2,1)\V'_{s^C_2}=((xx_C)^{-1}x^C_2,1)\V_{s_1^C}=
\V_{\theta_C((x_2^C)^{-1}x_Cx_1^Cx)}=\V_{f(\theta'_C(x))}=(f^*\V)_{\theta'_C(x)}
$$
for each $C$ and $x\in\cc(k)$. It follows that $\V'=f^*\V$, and hence that 
$(L,\V)$ and $(L,\V')$ are equivalent.
\end{proof}

%%%%%%%%%%%%%%%
